\documentclass[11pt]{amsart}
\usepackage[T1]{fontenc}
\usepackage{lmodern,amsmath,amssymb,amsthm,mathtools,microtype,booktabs,array,longtable}
\usepackage[margin=1.06in]{geometry}
\usepackage[hidelinks]{hyperref}
\emergencystretch=2em
\numberwithin{equation}{section}
\newtheorem{theorem}{Theorem}[section]
\newtheorem{lemma}[theorem]{Lemma}
\newtheorem{proposition}[theorem]{Proposition}
\newtheorem{corollary}[theorem]{Corollary}
\theoremstyle{definition}\newtheorem{definition}[theorem]{Definition}
\theoremstyle{remark}\newtheorem{remark}[theorem]{Remark}
\newcommand{\Q}{\mathbb Q}\newcommand{\R}{\mathbb R}\newcommand{\C}{\mathbb C}\newcommand{\N}{\mathbb N}\newcommand{\Z}{\mathbb Z}\newcommand{\E}{\mathbb E}
\newcommand{\ip}[2]{\langle #1,#2\rangle}\newcommand{\norm}[1]{\lVert #1\rVert}
\newcommand{\epsc}{\varepsilon_{\mathrm c}}
\DeclareMathOperator{\osc}{osc}\DeclareMathOperator{\spanop}{span}\DeclareMathOperator{\End}{End}\DeclareMathOperator{\diag}{diag}\DeclareMathOperator{\tr}{tr}\DeclareMathOperator{\rad}{rad}\DeclareMathOperator{\TV}{TV}
\title[Stochastic realization and orbit geometry]{Stochastic Realization, Finite Quotients,\\and Uniform Orbit Geometry}
\author{Qian Qi}
\date{28 September 2026}
\subjclass[2020]{Primary 93B15, 68Q45; Secondary 22C05, 60J10, 11J13}
\keywords{stochastic realization, finite quotient, Ramsey extraction, compact group, spectral gap, nonuniform width}
\begin{document}
\begin{abstract}
This complete research edition contains the structural and quantitative
stochastic-realization theory and all retained mathematical refinements.
Same-width, same-error stationarization under eventual approximate returns
and physical-equivariant purification give equality-inclusive finite
quotient criteria. The least orbit dimension supplies a uniform cap
exponent for all unit directions. A fixed rational qubit experiment has
an exponential exact profile, robust under exponentially small positive
error. Spectral entropy removes the logarithmic loss in the noisy
converse; least-orbit and target-orbit dimensions give sharp matching
laws when they agree. An additional rational Bloch experiment has
explicit spectral constants and exact profile $5^t$. Its positive upper realization is constructed by a uniform
rational compiler. The focused structural and quantitative articles are
independently complete; this edition preserves the combined development,
including arithmetic, profinite and existential higher-dimensional
examples. Repeatable probes additionally give a finite-response-quotient
classification for uniform causal process error and a constructive
rational approximation under adaptive control. Orthogonal covariance removes the quantitative return assumption and
gives joint subexponential-accuracy laws. A classical tail-capacity
argument sharpens the process minimum to every error below one.
No independent analytic pipeline closure is asserted.
\end{abstract}
\maketitle
\noindent\textit{Complete research edition. The focused quantitative and structural articles are independently complete; this edition retains the whole active mathematical development.}

\noindent The principal additional conclusions are
Theorem~\ref{thm:driftoccupation62} (return-free occupation),
Theorem~\ref{thm:jointlaw62} (joint accuracy law),
Theorem~\ref{thm:causalstrong62} (full-error process minimum),
Theorem~\ref{thm:noidle62} (six-command numerical example),
Theorem~\ref{thm:sharpoccupation61} (synchronized occupation),
Theorem~\ref{thm:effectivelps61} (effective spectral example),
Theorem~\ref{thm:minorbit58} (all-direction orbit caps),
Theorem~\ref{thm:rationalseparation58} (explicit rational robust separation),
Theorem~\ref{thm:rationalcompiler58} (constructed finite tables),
Theorem~\ref{thm:causalclassification59} (repeatable-observation
classification), and Theorem~\ref{thm:adaptivecompiler59} (adaptive
path-law construction).
\section{Introduction}
A finite stochastic register has a continuum of probability distributions. Its number of labels is therefore not determined by the cardinality of the numerical orbit it realizes. Allowing an external clock creates a further difficulty: the command rows may change at every position, and the entire realization may be redesigned at each horizon. The central question is whether these freedoms change the least width needed to preserve a fixed numerical experiment with a prescribed uniform error.

Our first theorem identifies the limiting clocked width with the stationary minimum, without a loss of labels or accuracy. The synchronization hypothesis is eventual approximation of the physical identity at every sufficiently large length. Under it, arbitrarily many cuts of width at most $k$ force a stationary realization on $k$ labels, even when the intervening registers are unrestricted. The proof selects approximately homogeneous composite transition tables by finite Ramsey theory. It applies to irreversible topological monoids as well as groups. It does not assume that a physical return acts trivially on the hidden register.

For compact groups a second reduction makes the stationary realization physically equivariant. Transition randomness can be removed, and hidden permutations above the physical identity can be averaged out, without increasing width or uniform error. The stationary minimum is therefore intrinsic to a finite continuous group action with randomized initialization. Combining the two reductions gives an equality-inclusive finite-quotient criterion for clocked boundedness, including positive critical radii on groups with infinitely many components. A moving-frame construction gives the exact limiting minimum in each return-length phase, again without additional labels.

\subsection{The interface and its two resources}
Unless a section explicitly specifies a different physical space, $H$ is a compact Hausdorff group and $\mathcal A$ is a finite nonempty alphabet whose letters $u_a\in H$ generate a dense subgroup. The execution convention is
\[
                  u_{a_1\cdots a_n}=u_{a_n}\cdots u_{a_1}.
\]
There are finite nonempty seed and query sets $\mathcal S,\mathcal J$ and continuous maps $F_{s,j}:H\to Y_j$. In the structural results $Y_j$ is a nonempty norm-compact convex subset of a real Banach space. The localization and algebraic sections specify finite-dimensional output spaces. For a categorical law, $Y_j=\Delta_{M_j}$ and the norm on differences is $\frac12\|\cdot\|_1$, namely total variation. A prescribed joint law is one query; incompatible measurements are not assigned an artificial joint law.

\begin{definition}[Stationary realization]\label{def:stationary55}
A realization on $k$ labels consists of initialization rows $\alpha_s\in\Delta_k$, row-stochastic matrices $T_a$ of order $k$, and decoder values $D_j(i)\in Y_j$, all independent of time and word length. Its error is
\[
       \sup_{s,j,w\in\mathcal A^*}
        \|\alpha_sT_{a_1}\cdots T_{a_n}D_j-F_{s,j}(u_w)\|_j.
\]
The empty word is included. Let $M_\varepsilon$ be the least $k$ at error at most $\varepsilon$, or $\infty$ when none exists. A permutation realization has permutation command matrices; its initialization need not be deterministic.
\end{definition}
The clocked resource $W_{N,\varepsilon}$ instead minimizes the maximum size of the registers $S_0,\ldots,S_N$ at a fixed horizon, with arbitrary row-stochastic transitions depending on time and horizon, and a terminal query decoder. Definition~\ref{def:machine54} states the model formally. Every advertised persistent label counts, including unused padding. Private information retained for future commands is part of the current label. The original width model does not charge the external clock, real table descriptions, their construction, exact arithmetic, or exact atomic sampling. Section~\ref{sec:complexity56} gives separate finite-description and fair-random-bit bounds; it does not reinterpret those resources as free computational memory. A numerical decoder vector is not an additional sampled-answer register.

For an open normal subgroup $U\le H$ put
\begin{equation}\label{eq:quotientradius55}
 r(U)=\max_{s,j,h\in H}\rad_{Y_j}(F_{s,j}(hU)),
 \qquad\rad_Y(A)=\min_{y\in Y}\sup_{x\in A}\|x-y\|.
\end{equation}
The coset images are compact by continuity. The radius depends continuously on the translate, and compactness gives the displayed extrema. The restriction of the center to the legal output set is essential. The component radius $R_0$ uses $U=H^\circ$, whether or not this subgroup is open; it is the infimum of the finite-quotient radii. At the critical error, the relevant issue is attainment of that infimum, not a limit from larger errors.

\subsection{The principal conclusions}
Theorem~\ref{thm:stationarization56} proves $W_{N,\varepsilon}\to M_\varepsilon$ under eventual approximate returns. Finite limits are eventually equal, and every $k<M_\varepsilon$ has a uniform narrow-cut occupation bound. Theorems~\ref{thm:equivariant56} and \ref{thm:completeboundary56} then give the intrinsic group-action minimum and its finite-quotient existence criterion. Theorem~\ref{thm:phaseminimum56} treats each length phase by a time-dependent change of physical frame. Table~\ref{tab:map56} separates these quantifiers and hypotheses.

The quantitative consequences have additional hypotheses. For a specified Diophantine planar alphabet the full-subcritical matching polynomial law is retained. A fixed $2$-adic interface has linear exact width and bounded width at every positive error. For a spherical interface in dimension $d$ with a spectral-gap alphabet, Theorem~\ref{thm:spherical56} gives a block estimate; Theorem~\ref{thm:sharpoccupation61} strengthens it to the matching order $N^{(d-1)/2}$. A rational alphabet generating $\mathrm{SO}(3)$ supplies a genuinely noncommuting example. Its spectral gap is imported from Bourgain--Gamburd, not inferred from a finite numerical test. Theorem~\ref{thm:ETR56} gives an $\exists\R$-complete finite input problem, while Theorem~\ref{thm:bits56} prices dyadic descriptions and fair random bits.


The quantitative theory also extends beyond spheres. Theorem~\ref{thm:uniformorbit57} separates a small-ball exponent uniform over all unit feature directions from the dimension of the target orbit. For conjugation on traceless Hermitian matrices a spectral split gives the uniform exponent $2(d-1)$, including directions with repeated eigenvalues. The pure-state orbit has the same dimension. Theorem~\ref{thm:matrixwidth57} therefore gives nearly matching width exponents $d-1$ for full legal density-matrix outputs in every dimension. Its exact upper construction stays inside the positive cone, rather than replacing the legal output set by an ambient Euclidean ball.

At an extreme target, the exact answer can be very different. Theorem~\ref{thm:extreme57} identifies the entire optimal exact width profile with the reachable orbit cardinalities. Theorem~\ref{thm:freeendpoint57} applies it to one fixed five-letter algebraic dense free alphabet in $SU(d)$: exact pure-state width is $2\cdot3^N-1$, while each fixed positive subcritical error has the polynomial bounds just described. The specified seed is transcendental, not a rational-input instance. These assertions concern a full matrix-valued terminal interface, not adaptive observations or the zero-error endpoint of a single scalar readout.

The compact nonnegative-matrix group structure used in purification is classical, and is separated from the external numerical error constraints throughout Section~\ref{sec:comparison54}. Neutral-letter collapse for advice automata is also a relevant predecessor to stationarization. The additional assertions proved here concern the same stochastic label bound and the same closed numerical error balls, arbitrary intervening widths, and approximate rather than exact returns. Neither classical theorem is relabeled as a new result.

\clearpage
\begin{table}[p]
\centering\small
\caption{Resources, hypotheses, and conclusions. All widths count persistent labels; all errors are worst-word numerical errors.}\label{tab:map56}
\begin{tabular}{p{.20\textwidth}p{.35\textwidth}p{.35\textwidth}}
\toprule
Object & Hypotheses & Conclusion\\
\midrule
$M_\varepsilon$ & Compact $H$; finite dense alphabet; compact convex outputs; no returns needed & Least finite continuous $H$-action with randomized initialization (Theorem~\ref{thm:equivariant56})\\[3pt]
$r(U)$ & $U$ open normal; arbitrary number of components & $M_\varepsilon<\infty$ iff some $r(U)\le\varepsilon$, equality included (Theorem~\ref{thm:finitequotient55})\\[3pt]
$W_{N,\varepsilon}$ & Eventual approximate returns; even an irreversible topological monoid & $W_{N,\varepsilon}\to M_\varepsilon$; finite limits eventually equal; all-profile occupation (Theorem~\ref{thm:stationarization56})\\[3pt]
$R_0$ & Compact group and eventual approximate returns & $R_0=\inf_Ur(U)$; critical boundedness iff the infimum is attained (Theorem~\ref{thm:completeboundary56})\\[3pt]
$M_{r,\varepsilon}$ & Block length $p$, phase subgroup $J$; block approximate returns & $W_{N,\varepsilon}\to M_{r,\varepsilon}$ for $N\equiv r$; no extra phase labels (Theorem~\ref{thm:phaseminimum56})\\[3pt]
$\varepsilon_r$ & Finitely many components; nonempty exact-return language & Component radius in the reachable phase; earlier localized occupation proof (Theorem~\ref{thm:phases55})\\[3pt]
Exact boundary & Cofinite exact returns; finite-dimensional outputs & Finite translate rank yields a finite observable quotient (Theorem~\ref{thm:zero55})\\[3pt]
Quantitative width & Separate arithmetic, profinite, or uniform-orbit spectral-gap hypotheses & Planar and dyadic laws; sharp spherical and matrix-orbit power laws (Theorems~\ref{thm:sharpoccupation61}, \ref{thm:matrixwidth57})\\[3pt]
Extreme outputs & Affine group action; target orbit consists of legal extreme points & Exact profile equals orbit cardinalities; free example $2\cdot3^N-1$ (Theorems~\ref{thm:extreme57}, \ref{thm:freeendpoint57})\\[3pt]
Finite encodings & Explicit finite group and rational categorical tables & $\exists\R$-complete stationary decision; dyadic precision and random-bit budgets (Theorems~\ref{thm:ETR56}, \ref{thm:bits56})\\
\bottomrule
\end{tabular}
\end{table}
\clearpage

\section{Repeatable observations and causal realization}\label{sec:causal59}
Terminal numerical accuracy and accuracy of a generated observation process
are different requirements. Repeating a nondisturbing observation can
separate arbitrarily close, but distinct, numerical responses. We now
formulate this distinction within the same finite-label resource model.

Let $H$ be a compact Hausdorff group, let $\mathcal A$ be a finite nonempty
command alphabet generating a dense subgroup, and let $\mathcal S$ and
$\mathcal J$ be finite nonempty sets. For each $s,j$ fix a continuous law
\[
 p_{s,j}:H\longrightarrow\Delta_{M_j}.
\]
The physical state starts at $g=e$. A command $a$ replaces $g$ by $u_ag$
and emits a fixed null symbol. A probe $j$ leaves $g$ unchanged and emits
a fresh draw from $p_{s,j}(g)$, conditionally on the public past
and current physical state. At each step the next command or probe may be chosen from the entire
observed transcript by a randomized policy. The seed is specified at
initialization. Probes are repeatable statistical observations; when
$p_{s,j}$ is obtained from a density matrix, this is an oracle model, not
repeated nondisturbing measurement of one unknown quantum system.

\begin{definition}[Causal finite-label machine]\label{def:causal59}
At horizon $N$ a machine has registers $S_0,\ldots,S_N$, initialization
rows $E_s$, and, for each control $c$ and possible output $y$, a
nonnegative matrix $K_{t,c,y}:S_{t-1}\to S_t$ such that
$\sum_yK_{t,c,y}$ is row-stochastic. Its entries jointly specify the
emission and next label. The matrices may depend on $t,N$ but on the past
only through the retained label and current control. Its error is
\[
 \sup_{s,\pi}\TV\bigl(\mathbb P^{\pi}_{s,N},
                            \mathbb Q^{\pi}_{s,N}\bigr),
\]
where the laws include the complete control--output transcript under the
same policy $\pi$. Write $\mathsf{PW}_{N,\varepsilon}$ for the minimum
peak register size at error at most $\varepsilon$.
\end{definition}
Every advertised label and all retained private randomness are counted.
The horizon, external clock and arbitrary real tables remain free. The
policy is an external experimenter, not an uncharged private memory for
the machine. Outputs may influence future controls, but are not supplied
again as a free history input. A sampled output register, when charged,
adds the separate fixed output-alphabet cost. The error above is one
budget for the entire transcript, not a per-step error budget.

For $(s,g)$ define its future response function
\begin{equation}\label{eq:futuretype59}
 \Theta_{s,g}(x,j)=p_{s,j}(xg),\qquad x\in H,
 \qquad \mathcal X=\{\Theta_{s,g}:s\in\mathcal S,\ g\in H\}.
\end{equation}
Equip the finite product of continuous-function spaces with the uniform
norm. Put $n=|\mathcal X|$ when this set is finite, and $n=\infty$
otherwise. This is a behavioral quotient, not the dimension of a linear
predictive representation. In particular, infinitely many distinct
predictive rows do not by themselves imply infinite nonnegative
realization size for a general hidden Markov process
\cite[Theorem~2 and Section~4]{SJR2004}. Repeatable observations are the
additional input in the theorem below.

\begin{theorem}[Uniform process error and finite physical actions]
\label{thm:causalclassification59}
For every fixed $0\le\varepsilon<1$ the following are equivalent:
\begin{enumerate}
\item $\sup_N\mathsf{PW}_{N,\varepsilon}<\infty$;
\item $n<\infty$;
\item some open normal subgroup $U\lhd H$ makes every $p_{s,j}$ constant
      on each coset of $U$;
\item a stationary exact causal realization exists with permutation
      command updates and state-preserving probe emissions.
\end{enumerate}
If $n<\infty$ and $0\le\varepsilon<1/2$, then
\begin{equation}\label{eq:causalminimum59}
                   \mathsf{PW}_{N,\varepsilon}=n
                   \quad\hbox{for all sufficiently large }N.
\end{equation}
If $n=\infty$, then, for every $\varepsilon<1$ and every $k\ge1$,
there is $L_{k,\varepsilon}<\infty$ such that every admissible realization
satisfies
\begin{equation}\label{eq:causaloccupation59}
 \#\{t\in\{1,\ldots,N\}:|S_t|\le k\}\le L_{k,\varepsilon}.
\end{equation}
The same occupation conclusion holds for finite $n$, $\varepsilon<1/2$
and $k<n$. Intervening registers are unrestricted. No return condition
is imposed on the command alphabet: the probes themselves leave the
physical state unchanged and supply the padding used in the proof.
\end{theorem}
The threshold $1$ in the boundedness criterion is sharp: total variation
never exceeds one, so at tolerance one a single label always suffices.
The elementary one-continuation proof of \eqref{eq:causalminimum59}
uses $\varepsilon<1/2$. Theorem~\ref{thm:causalstrong62} below
proves the full minimum for every $\varepsilon<1$ by a different
tail-event argument; the present proof and its stronger one-cut test
remain useful in their stated range.

\begin{lemma}[A common separating itinerary]\label{lem:itinerary59}
Choose $m\ge2$ different functions in \eqref{eq:futuretype59}, represented
by executable seed--prefix pairs. For every $\delta>0$ there is one
finite deterministic control itinerary with disjoint output events
$A_1,\ldots,A_m$ such that its target law from the $i$th state assigns
probability at least $1-\delta$ to $A_i$.
\end{lemma}
\begin{proof}
Equality of two future response functions is preserved and reflected by
left multiplication of their physical states by any $u\in H$: replace
$x$ in \eqref{eq:futuretype59} by $xu$, a bijection of $H$. Positive
executable products are dense in $H$. Indeed positive powers approximate
the inverse of every generator in a compact group, so the closed positive
semigroup contains the dense generated group.

List the $\binom m2$ pairs. At each stage the current two states for that
pair still have different future response functions. By continuity and
positive-word density, a finite executable word followed by one probe
separates some coordinate of their output laws. Append that word and
record the coordinate, then proceed to the next pair. This constructs
one itinerary, not a family of incompatible pair-dependent suffixes.
No reset or inverse word is required. Let $J=\binom m2$, and let
$\mu_{i,r}$ be the selected coordinate probability at the $r$th probe
location when the initial state is the $i$th representative. There is
$\Delta>0$ such that every pair of rows of this finite matrix differs
by at least $\Delta$ in some coordinate.

At each location repeat the chosen probe $L$ times. These repetitions
leave the physical state fixed. The empirical coordinate frequencies
$\widehat\mu_r$ obey
\[
 \mathbb P_i\{\max_r|\widehat\mu_r-\mu_{i,r}|\ge\Delta/3\}
       \le 2J\exp(-2L\Delta^2/9).
\]
For completeness, the scalar bound follows from
$\E e^{z(B-\E B)}\le e^{z^2/8}$ for a Bernoulli variable, exponential
Markov inequality, and optimization in $z$; a union bound gives the
display. The moment bound follows by differentiating the logarithmic
moment generating function twice: its second derivative is the variance
of a tilted Bernoulli variable and is at most $1/4$.
Choose $L\ge1$ with the displayed upper bound at most $\delta$.
The events
$A_i=\{\max_r|\widehat\mu_r-\mu_{i,r}|<\Delta/3\}$ are disjoint,
because $2\Delta/3<\Delta$, and have the required probabilities.
\end{proof}

\begin{lemma}[Testing across a finite register]\label{lem:testingcut59}
Suppose $m$ equal-length seed--prefix pairs at cut $t$ admit an itinerary
of length $\ell$ as in Lemma~\ref{lem:itinerary59}. For every
$N\ge t+\ell$, an error-$\varepsilon$ causal realization satisfies
\begin{equation}\label{eq:testingcut59}
                  |S_t|\ge m(1-\varepsilon-\delta).
\end{equation}
If $\varepsilon+\delta<1/2$, the stronger bound $|S_t|\ge m$ holds.
\end{lemma}
\begin{proof}
Use the same deterministic suffix for all representatives and pad its
end with probes if necessary. Marginalize all prefix outputs. Write
$\alpha_i(z)$ for the resulting cut-label distribution and $Q_z$ for
the suffix-output law conditional on label $z$. Causality makes $Q_z$
independent of the chosen prefix. Total variation contracts under
marginalization, so correctness gives
\[
 1-\delta-\varepsilon
     \le\sum_z\alpha_i(z)Q_z(A_i).
\]
This uses no conditional correctness assertion after a rare observed
prefix. Summing in $i$, using $\alpha_i(z)\le1$, and then disjointness,
yields
\[
 m(1-\delta-\varepsilon)
 \le\sum_z\sum_iQ_z(A_i)\le |S_t|.
\]
If $1-\delta-\varepsilon>1/2$, each $i$ has a positive-mass label $z_i$
with $Q_{z_i}(A_i)>1/2$. Disjointness forces the labels $z_i$ to be
different. This proves the stronger assertion.
\end{proof}

\begin{proof}[Proof of Theorem~\ref{thm:causalclassification59}]
The map $g\mapsto\Theta_{s,g}$ is continuous in the uniform norm.
This follows from continuity on the compact product $H\times H$ and a
finite-cover argument, uniformly in the first coordinate. Thus each
seed's image is compact. Executable images are dense in it. If
$\mathcal X$ is infinite, arbitrarily many distinct types therefore
have executable representatives; if it is finite, every type has one.

Choose finitely many such representatives and pad their prefixes to a
common length $t_0$ by any fixed probe. Their physical states are
unchanged. Apply Lemma~\ref{lem:itinerary59}. For every
$t_0\le t\le N-\ell$, pad further with the same probe, marginalize its
outputs, and apply Lemma~\ref{lem:testingcut59}. The itinerary and
$\ell$ do not depend on $t,N$ or on the machine.

If $n=\infty$, choose $m$ with $m(1-\varepsilon)>k$ and then choose
$\delta>0$ with $m(1-\varepsilon-\delta)>k$. Every interior cut has
width greater than $k$. At most $t_0+\ell$ positive cuts lie outside
this interval, including when it is empty. This proves
\eqref{eq:causaloccupation59}, divergence, and the implication (i)$\Rightarrow$(ii).
For finite $n$ and $\varepsilon<1/2$, use $m=n$ and
$0<\delta<1/2-\varepsilon$ when $n\ge2$. All sufficiently interior
cuts require $n$ labels. The case $n=1$ follows from nonempty registers.

For finite $\mathcal X$, the maps
\[
 \Theta(x,j)\longmapsto\Theta(xu,j),\qquad u\in H,
\]
define a continuous permutation action of $H$ on $\mathcal X$.
Continuity follows from uniform continuity of each of the finitely many
functions. Its kernel $U$ is open and normal. Evaluating at $x=e$
shows that every $p_{s,j}$ is constant on $U$-cosets, proving
(ii)$\Rightarrow$(iii). Conversely (iii) gives at most
$|\mathcal S|[H:U]$ response types. More economically, use the $n$
distinct types as labels, initialize at $\Theta_{s,e}$, update by the
permutation above, and emit $\Theta(e,j)$ on probe $j$ without changing
the label. This produces the exact conditional response at each step.
Induction on transcript length proves equality for every adaptive
policy, hence (iv) and (i). It also supplies the upper bound $n$ in
\eqref{eq:causalminimum59}. The occupation assertion for $k<n$ follows
from the stronger testing bound just established.
\end{proof}

\begin{corollary}[Connected groups]\label{cor:connectedprocess59}
If $H$ is connected, uniformly bounded causal width at any fixed total
transcript error below one is possible exactly when every $p_{s,j}$ is
constant on $H$. If at least one is nonconstant, every fixed width has
the occupation bound \eqref{eq:causaloccupation59}.
\end{corollary}
\begin{proof}
A continuous action of a connected group on a finite discrete set is
trivial. The assertion follows from the theorem.
\end{proof}

\begin{remark}[Why terminal purification is insufficient]\label{rem:correlation59}
A two-label variable chosen once with equal masses and then repeatedly
reported has the correct Bernoulli$(1/2)$ marginal at every time, but its
length-$L$ law is supported on the two constant strings. Its total
variation distance from $L$ independent fair bits is $1-2^{1-L}$.
Thus even exact numerical marginals do not supply a path-law realization.
The causal theorem concerns an actual nonnegative emission instrument,
not a list of separately correct decoders. It neither classifies
arbitrary irreversible hidden-state systems nor asserts a measurement
procedure for noncommuting quantum observables.
\end{remark}

\section{A strong converse for repeatable observation processes}
\label{sec:tail62}
The one-continuation testing bound distinguishes the full number of
response types only when the error is below one half. Repeating a
common diagnostic cycle gives a stronger obstruction: a finite hidden
register cannot support more positive-probability observation-tail types
than it has labels. The underlying tail-event principle is classical
\cite{Blackwell62,Cohn62,GS62}. We give the short argument needed here,
including the hidden-emission and nonhomogeneous conventions, before
applying compactness to separately redesigned finite-horizon machines.

\begin{lemma}[Tail capacity at narrow cuts]\label{lem:tailcapacity62}
Let $(S_t,Y_t)_{t\ge0}$ be a hidden emission process with finite
registers and finite output alphabets, whose joint next-state and
emission law depends on the entire past only through $S_t$ and the
external time. Its transition kernels may be nonhomogeneous and
irreversible. Suppose $|S_t|\le k$ at infinitely many deterministic
times. Then there cannot be more than $k$ pairwise disjoint observation
tail events having positive probability. Here an observation tail event
belongs to $\bigcap_{t\ge0}\sigma(Y_{t+1},Y_{t+2},\ldots)$.
\end{lemma}
\begin{proof}
Let $B_1,\ldots,B_r$ be such events. Include the initialization variable,
all past hidden labels, and all past emissions in $\mathcal F_t$.
The Markov property for future emissions gives a version
\[
 \mathbb P(B_i\mid\mathcal F_t)=b_{i,t}(S_t),\qquad
 b_{i,t}(s)=\mathbb P(B_i\mid S_t=s).
\]
Only labels of positive probability are relevant. Since each $B_i$ is
measurable with respect to $\mathcal F_\infty=\sigma(\bigcup_t\mathcal F_t)$,
upward martingale convergence gives
\[
 b_{i,t}(S_t)\longrightarrow\boldsymbol1_{B_i}
                                  \quad\hbox{almost surely and in }L^1.
\]
Consequently, for every sufficiently large $t$ and each $i$, some
positive-probability label $s_i$ satisfies $b_{i,t}(s_i)>1/2$.
There is one common sufficiently large time for the finite collection
of events. Disjointness implies $\sum_i b_{i,t}(s)\le1$, so these
labels are distinct. At an arbitrarily late narrow cut, $r\le k$.
This is the conditional-probability zero--one mechanism of
\cite[Theorem~1 and Section~5.1]{GS62}, applied to the emitted tail;
no stationarity, mixing rate, or lower bound on state masses is used.
\end{proof}

Fix pairwise distinct laws $P_1,\ldots,P_n$ on one finite alphabet
$\Omega$. An $R$-step $k$-label instrument has one initialization row
$\alpha_i$ for each $i$ and common, possibly time-dependent,
nonnegative emission matrices $K_{t,y}$ of size $k\times k$ with
$\sum_yK_{t,y}$ row-stochastic. Let $Q_{i,R}$ be its output law, and set
\begin{equation}\label{eq:producterror62}
 e_{k,R}=\min_{\alpha_i,K_{t,y}}
             \max_{1\le i\le n}\TV(Q_{i,R},P_i^{\otimes R}).
\end{equation}
The minimum is attained in a finite product of compact simplexes.
Initialization may be randomized and may depend on $i$; subsequent
kernels may not depend on an unrecorded value of $i$.

\begin{proposition}[Uniform finite witnesses for product laws]
\label{prop:productconverse62}
If $k<n$, then $e_{k,R}$ is nondecreasing in $R$ and
\begin{equation}\label{eq:productconverse62}
                         \lim_{R\to\infty} e_{k,R}=1.
\end{equation}
In particular, for every $\varepsilon<1$ there is an integer
$R(k,\varepsilon,P_1,\ldots,P_n)$ for which no such instrument has
error at most $\varepsilon$. When the laws and the tolerance are
rational, one can find such an integer by a terminating sequence of
finite real-algebraic feasibility decisions. No effective bound on the
number of decisions, or polynomial-time claim, is made.
\end{proposition}
\begin{proof}
Truncation and contraction of total variation under marginalization
show that $e_{k,R}\le e_{k,R+1}$. Suppose its limit were some $e<1$.
Consider the compact countable product consisting of the initialization
rows and all kernels $(K_{t,y})_{t\ge1}$. For every $R$, impose all
closed constraints
$\TV(Q_{i,r},P_i^{\otimes r})\le e$ for $i\le n,r\le R$.
Each finite system is feasible by attainment in \eqref{eq:producterror62}.
Nested compactness gives one infinite hidden process satisfying every
finite constraint, with common kernels and $k$ labels at every time.

Let $Q_i$ denote its infinite output law and $P_i^\infty$ the product
law. The finite constraints imply $\TV(Q_i,P_i^\infty)\le e$.
Indeed, for any event, approximate its indicator in $L^1$ under the
finite measure $Q_i+P_i^\infty$ by indicators of cylinder events;
these generate the product sigma-field. The bound on event-probability
differences passes to the limit. Equivalently, one can apply martingale
convergence to conditional probabilities under this finite measure.

For each $i$, let $B_i$ be the event that empirical frequencies of all
symbols converge to $P_i$. These are disjoint observation tail events,
and the strong law gives $P_i^\infty(B_i)=1$. Hence
$Q_i(B_i)\ge1-e>0$. Mix the initialization index uniformly, retaining
the same common transition kernels. The resulting hidden process has
output law $Q=n^{-1}\sum_iQ_i$, so $Q(B_i)>0$ for every $i$.
Lemma~\ref{lem:tailcapacity62} contradicts $n>k$. This proves the limit.

For the last assertion, introduce the finitely many stochastic rows as
real variables. Every word probability is a polynomial in their entries;
the targets are rational. Absolute-value slack variables express each
finite total-variation inequality by polynomial equalities and
inequalities. Real quantifier elimination decides feasibility. Test
$R=1,2,\ldots$ until it is infeasible. The limit just proved guarantees
termination. This is a finite-witness search, not an implementation of a
general minimum-width algorithm.
\end{proof}

\begin{lemma}[A diagnostic cycle]\label{lem:diagnosticcycle62}
Suppose the response-type set $\mathcal X$ in
\eqref{eq:futuretype59} is finite, of cardinality $n\ge2$.
There is a finite deterministic control word $B$ of positive length
$\ell$ such that its action on $\mathcal X$ is the identity and its
output laws $P_1,\ldots,P_n$ from the $n$ types are pairwise distinct.
Successive executions of $B$ give independent, identically distributed
$P_i$-blocks from initial type $i$. Any intervening fixed probe controls
may be inserted if their outputs are discarded.
\end{lemma}
\begin{proof}
Every type has an executable seed--prefix representative, as proved in
Theorem~\ref{thm:causalclassification59}. Apply
Lemma~\ref{lem:itinerary59} to all $n$ types with, for example,
$\delta=1/4$. The resulting finite deterministic itinerary $T$ has
pairwise distinct output laws: its disjoint diagnostic events have
probability at least $3/4$ under the corresponding laws. Its command
product induces a permutation $\pi$ of the finite set $\mathcal X$.
Let $d\ge1$ be the order of this permutation and take $B=T^d$.
Its action on types is the identity, and its output laws remain
distinct because their first-$T$ marginals are distinct.

The output law of any fixed continuation depends only on the response
type, not on its chosen representative. Probes emit fresh conditionally
independent observations and have no physical effect. Since $B$ returns
the type deterministically, its conditional output law after any number
of previous blocks is again $P_i$. Induction gives the product law.
Intervening probes also preserve the type; their ignored outcomes do not
change the deterministic controls or the conditional block laws. The
cycle is a return on the \emph{finite response quotient}, not an assumed
identity-product word in the physical group.
\end{proof}

\begin{theorem}[The full-error causal minimum]\label{thm:causalstrong62}
In the repeatable-probe model of Definition~\ref{def:causal59}, let
$n=|\mathcal X|\in\N\cup\{\infty\}$. For every fixed
$0\le\varepsilon<1$,
\begin{equation}\label{eq:causalstrong62}
                 \lim_{N\to\infty}\mathsf{PW}_{N,\varepsilon}=n.
\end{equation}
When $n<\infty$, equality with $n$ holds at every sufficiently large
horizon. For every integer $k<n$ there is $L_{k,\varepsilon}<\infty$
such that every error-$\varepsilon$ realization has
\begin{equation}\label{eq:causalstrongoccupation62}
          \#\{1\le t\le N:|S_t|\le k\}\le L_{k,\varepsilon},
\end{equation}
with no restriction on intervening widths. Thus the finite-response
cardinality is the eventual minimum throughout the whole interval
below the maximal total-variation error, not only below one half.
The full-transcript and fresh nondisturbing probe assumptions are
unchanged. No command return condition is imposed.
\end{theorem}
\begin{proof}
For $n=\infty$, Theorem~\ref{thm:causalclassification59} already proves
\eqref{eq:causalstrongoccupation62} for every $k$ and
$\varepsilon<1$. For $n=1$, nonempty registers and the exact one-type
machine prove the assertion. Suppose $2\le n<\infty$ and $k<n$.
Choose executable representatives of all types, and pad their prefixes
with a fixed probe to one common length $b\ge0$. Let $B$, $\ell$ and
$P_i$ be as in Lemma~\ref{lem:diagnosticcycle62}. Fix an integer $R$
for which $e_{k,R}>\varepsilon$, supplied by
Proposition~\ref{prop:productconverse62}.

Consider an advertised profile. Discard narrow cuts with $t<b$, and
greedily select from the others
$t_1<\cdots<t_J$ with gaps at least $\ell$. If there are $R+1$
selected cuts, start from each of the equal-length representatives,
pad to $t_1$ with the same probe, and put one copy of $B$ immediately
after each of $t_1,\ldots,t_R$. Fill the remaining positions before
the next selected cut with that probe. Complete the horizon with any
fixed controls. For each representative the retained $R$ block outputs
have target law $P_i^{\otimes R}$.

Marginalize all prefix, filler, and post-block outputs. In the proposed
machine, the actual transition from one selected cut to the next,
jointly with the retained block output, is a nonnegative instrument
between at most $k$ labels. Summing over the discarded emissions gives
its stochastic normalization. All its entries are common to the
representatives: the schedule is deterministic and the retained label
is the only private past information available to a row. Intervening
registers may be arbitrarily wide. The cut-$t_1$ distributions are
allowed to be the arbitrary initialization rows $\alpha_i$ in
\eqref{eq:producterror62}. After padding unused labels, these composite
kernels therefore define an $R$-step $k$-label instrument.

Whole-transcript correctness under these deterministic policies and
contraction of TV under marginalization give its error at most
$\varepsilon$. This contradicts $e_{k,R}>\varepsilon$. Thus $J\le R$.
The discarded cuts number at most $b$, and each selected cut accounts
for at most $\ell$ remaining narrow cuts. Consequently one may take
\[
                         L_{k,\varepsilon}=b+\ell R.
\]
This count includes empty eligible intervals and excludes cut zero.
The prefix, cycle, and witness integer depend only on the interface,
$k$ and $\varepsilon$, not on the horizon or machine. No correctness
conditional on a rare observed prefix was used.

An exact stationary $n$-type machine works at every horizon by
Theorem~\ref{thm:causalclassification59}. Applying the occupation bound
to $k=n-1$ shows that a smaller peak is impossible once
$N>L_{n-1,\varepsilon}$. This proves eventual equality and the theorem.
\end{proof}

\begin{corollary}[Strong converse below the response cardinality]
\label{cor:processerrorone62}
If $k<n$, the least worst-policy whole-transcript error achievable by
peak width at most $k$ tends to one as the horizon tends to infinity.
At error one, one label always suffices. When the finite group action
and all probe tables are explicitly rational, the proof supplies a
terminating search for a valid occupation constant via the finite
product-law feasibility problems.
\end{corollary}
\begin{proof}
For each $\varepsilon<1$, the preceding occupation bound excludes a
$k$-peak realization for every $N>L_{k,\varepsilon}$. TV is at most one,
which proves the limit and the endpoint. For explicit finite data,
response types and their command permutations can be enumerated. The
common finite itinerary and its permutation order can be found by finite
search, its cycle laws are rational, and
Proposition~\ref{prop:productconverse62} finds $R$. The claimed bound
$b+\ell R$ then follows. No efficiency conclusion follows from this
termination proof.
\end{proof}

\begin{proposition}[Why one half is not the process boundary]
\label{prop:binarytail62}
For $P_1=\delta_0$, $P_2=\delta_1$ and one hidden label,
\begin{equation}\label{eq:binarytail62}
                         e_{1,R}=1-2^{-R}.
\end{equation}
Nevertheless, a single width-one cut, followed by two-label storage,
can have worst error one half for these two product laws at arbitrarily
large remaining horizons.
\end{proposition}
\begin{proof}
A one-label instrument has independent time-dependent emissions with
probabilities $p_t$ of zero. Its two target-string probabilities are
$u=\prod_t p_t$ and $v=\prod_t(1-p_t)$. The two TV errors are $1-u$
and $1-v$. Since $uv\le4^{-R}$, their maximum is at least $1-2^{-R}$,
attained by $p_t=1/2$ for every $t$. For the last assertion, after the
one-label cut choose a fair bit, store it in two labels, and emit it
repeatedly. Its law is the equal mixture of the two deterministic
strings and is at distance one half from each. Thus for errors above
one half the theorem needs control of \emph{repeated} narrow cuts;
it does not assert the false pointwise lower bound at every interior cut.
\end{proof}

The tail-capacity lemma is a necessary condition for arbitrary hidden
emission processes, including irreversible ones. The diagnostic-cycle
application, however, still uses the stated repeatable-probe model.
It is neither a general HMM/POMDP realization classification nor a
claim about nondisturbing repeated quantum measurements. The new
conclusion is the all-error finite-horizon strong converse and occupation
bound obtained from the classical tail principle, finite-product
compactness, and one common physical diagnostic cycle.

\section{Stationarization from recurrent narrow cuts}\label{sec:ramsey56}
The first reduction does not require reversible physical dynamics. In this section only, replace $H$ by a topological monoid $G$ with identity $1$, retain a finite command alphabet, and let $F_{s,j}:G\to Y_j$ be continuous. The output sets are norm-compact convex subsets of real Banach spaces. Define $M_\varepsilon$ and $W_{N,\varepsilon}$ by the same stochastic models. Compactness of the physical monoid is not required.

\begin{definition}[Eventual approximate returns]\label{def:returns56}
The alphabet has eventual approximate returns if, for every neighborhood $V$ of $1$, there is an integer $g(V)\ge0$ such that every integer $n\ge g(V)$ is the length of a word with product in $V$. Cofinite exact returns mean that a single $g$ works with product exactly $1$.
\end{definition}
The statement concerns physical products only. The hidden rows on a return word need not be close to an identity matrix. Approximation here is eventual at every length, not merely along a return subsequence.

\begin{theorem}[Width-preserving stationarization]\label{thm:stationarization56}
Assume eventual approximate returns. For every $\varepsilon\ge0$ and integer $k\ge1$, the following are equivalent:
\begin{enumerate}
\item a stationary realization with at most $k$ labels has error at most $\varepsilon$;
\item among error-$\varepsilon$ clocked realizations, the number of positive cuts having width at most $k$ is unbounded.
\end{enumerate}
Consequently, if $M_\varepsilon>k$, there is $L_{k,\varepsilon}<\infty$, independent of the horizon, such that every error-$\varepsilon$ realization satisfies
\begin{equation}\label{eq:occupation56}
 \#\{t\in\{1,\ldots,N\}:|S_t|\le k\}\le L_{k,\varepsilon}.
\end{equation}
All other registers may be arbitrarily large. Moreover,
\begin{equation}\label{eq:limitminimum56}
                  \lim_{N\to\infty}W_{N,\varepsilon}=M_\varepsilon
                  \quad\hbox{in }\N\cup\{\infty\}.
\end{equation}
If the right side is finite, equality holds for every sufficiently large $N$.
\end{theorem}
The theorem preserves the proposed width, including randomized initialization. It does not assert that the microscopic rows of the clocked machine become stationary. They can be entirely unrelated at different times.

\subsection{Finite witnesses and homogeneous transition tables}
For a matrix $A$ write $\|A\|_{\mathrm r}=\max_i\sum_j|A_{ij}|$. Put $a=|\mathcal A|$ and $D_* =\max(1,\max_j\operatorname{diam}Y_j)$. Let $\mathcal X_k$ be the compact parameter space of stationary $k$-label machines. For $m\ge0$ define
\[
 e_{k,m}=\min_{X\in\mathcal X_k}\max_{s,j,\,|w|\le m}
               \|X_{s,j}(w)-F_{s,j}(u_w)\|_j.
\]
The finite maximum is continuous, so the minimum is attained. If no stationary error-$\varepsilon$ machine exists, compactness of $\mathcal X_k$ and the finite intersection property give an $m$ with
\begin{equation}\label{eq:finitewitness56}
                       e_{k,m}=\varepsilon+\gamma,\qquad\gamma>0.
\end{equation}
Indeed the nested closed sets defined by the constraints $|w|\le m$ otherwise have a common point. This argument also applies when the full infinite-word error is not a continuous function of the parameters.

We shall use the finite edge-coloring theorem of Ramsey \cite{Ramsey}. Denote by $R_q(h)$ an integer such that every coloring of the edges of a complete graph on $R_q(h)$ vertices with at most $q$ colors has a monochromatic $h$-vertex subgraph. Only finiteness is essential.

\begin{remark}[An explicit but coarse Ramsey bound]
For definiteness one may take
\begin{equation}\label{eq:ramseybound56}
                   R_q(h)\le q^{qh+1}\qquad(q\ge2,\ h\ge2).
\end{equation}
To see this elementary bound, successively select a vertex and retain a largest color class among its remaining neighbors. With $q^{qh+1}$ initial vertices this produces at least $q(h-1)+1$ selected vertices, each of whose edges to later selected vertices has one color. At least $h$ of their associated colors agree; these vertices form the required subgraph. The deliberately excessive exponent avoids any endpoint rounding issue.
\end{remark}

Partition $[0,1]$ into at most $\lceil k/\delta\rceil+1$ intervals of length at most $\delta/k$. Two tuples in the same entrywise cell of
\[
                  \mathsf{Stoch}(k)^{\{0\}\cup\mathcal A}
\]
have corresponding matrices within $\delta$ in $\|\cdot\|_{\mathrm r}$. The number of occupied cells is at most
\begin{equation}\label{eq:colors56}
             q=\max\bigl(2,(\lceil k/\delta\rceil+1)^{(a+1)k^2}\bigr).
\end{equation}
Each representative used below is chosen from an occupied cell and is therefore stochastic; an entrywise rounded matrix need not be stochastic and is not used as a representative.

\begin{proof}[Proof of Theorem~\ref{thm:stationarization56}]
The forward implication follows by using one stationary machine at every horizon. For the converse, suppose \eqref{eq:finitewitness56}. Choose $0<\eta<\gamma/4$. Continuity of multiplication and of the finitely many $F_{s,j}$ gives a neighborhood $V$ of $1$ with the following property. In any word of length at most $m$, inserting at most $m+3$ factors from $V$, including before and after the word, changes every output by less than $\eta$. There are only finitely many relevant command words and insertion patterns; continuity at the tuple of identity factors, followed by a finite intersection of neighborhoods, proves the assertion. This uses neither an invariant metric nor commutativity.

Let $g=g(V)$, put $b=g+1$, and choose
\[
       0<\delta\le \min\bigl(1,\gamma/(4D_*(m+1))\bigr).
\]
Take a clocked realization with so many narrow cuts that, after discarding cuts $t<g$ or $t>N-g$ and greedily retaining cuts separated by at least $b$, at least $R_q(m+2)$ remain. Such a choice is possible by hypothesis; at most $2g$ cuts are discarded and each retained cut accounts for at most $b$ of the others.

Identify each retained register with a subset of $\{1,\ldots,k\}$. Extend rows from dummy labels arbitrarily to stochastic rows without altering any row on the actual register. For retained cuts $t_i<t_j$ write $\ell=t_j-t_i\ge g+1$. Choose a length-$\ell$ return word with product in $V$, and let $A_{ij}$ be its actual composite kernel between these cuts. For each letter $c\in\mathcal A$, choose the word consisting of $c$ followed by a length-$(\ell-1)$ return word with product in $V$, and let $B_{ij,c}$ be its composite kernel. Thus the physical products of these words are respectively in $V$ and of the form $v u_c$ with $v\in V$. All the matrices are of order $k$, however large the intermediate registers may be.

Color the edge $\{i,j\}$, oriented from the earlier cut to the later cut, by the cell of $(A_{ij},(B_{ij,c})_c)$. Ramsey's theorem gives $m+2$ ordered cuts
\[
                         t_0<t_1<\cdots<t_{m+1}
\]
whose edges have one color. Fix a stochastic representative $(E,(B_c)_c)$ from that color. Choose fixed return words for the prefix of length $t_0$ and the suffix of length $N-t_{m+1}$, both with products in $V$. Let $\alpha_s$ be the actual state distribution after the prefix. Let $D_j(i)$ be the conditional expected original decoder after the suffix, given label $i$ at its start, with arbitrary legal values on dummy labels. Convexity gives $D_j(i)\in Y_j$.

Consider the stationary machine
\begin{equation}\label{eq:extractedmachine56}
                  \alpha_s,\qquad T_c=B_c,\qquad d_j=ED_j.
\end{equation}
Its decoder is legal. To test $w=c_1\cdots c_\ell$, $0\le\ell\le m$, use the first $\ell$ consecutive clique edges with kernels $B_{i-1,i,c_i}$ and then the single edge from $t_\ell$ to $t_{m+1}$ with kernel $A_{\ell,m+1}$. For $\ell=0$ there is only this last edge. Add the fixed prefix and suffix. This is a genuine word of length $N$; its physical product is obtained from $u_w$ by at most $\ell+3$ insertions from $V$. Its target output is consequently within $\eta$ of $F_{s,j}(u_w)$.

Each of its $\ell+1$ edge matrices differs from the corresponding matrix in \eqref{eq:extractedmachine56} by at most $\delta$. Stochastic multiplication contracts the $\ell^1$ distance between probability rows. A telescoping expansion, and subtraction of any fixed point of $Y_j$ from the decoder, therefore bound the difference between the actual output and $\alpha_s B_{c_1}\cdots B_{c_\ell}E D_j$ by
\[
                           D_*(\ell+1)\delta.
\]
Wordwise correctness of the original machine implies $e_{k,m}\le\varepsilon+\eta+D_*(m+1)\delta<\varepsilon+\gamma$, a contradiction. The representative $E$ is used only as part of the terminal decoder. No idempotence of $E$, consistency between different filler words, or hidden-state identity on a physical return was assumed.

The same argument shows that one may take
\begin{equation}\label{eq:ramseyloss56}
              L_{k,\varepsilon}=2g+(g+1)R_q(m+2).
\end{equation}
For exact cofinite returns, use exact fillers and omit $\eta$ and the neighborhood choice; $g$ is the given exact-return bound. In general $g$ depends on the finite witness and the continuity neighborhood. The remaining parameters depend only on the interface, the alphabet, $k$, and $\varepsilon$, never on a machine or horizon.

If $M_\varepsilon=\infty$, \eqref{eq:occupation56} for each $k$ gives $W_{N,\varepsilon}\to\infty$. If $M_\varepsilon=K<\infty$, the stationary machine gives $W_{N,\varepsilon}\le K$ for every $N$, and the bound for $k=K-1$ gives eventual equality when $K>1$. When $K=1$ equality holds for every $N$ because registers are nonempty.
\end{proof}

\begin{proposition}[Approximate returns without exact returns]\label{prop:approxreturns56}
Let the two circle commands have angles $\theta,\phi$, where $1,\theta/(2\pi),\phi/(2\pi)$ are rationally independent. They have no nonempty exact-return word, but satisfy Definition~\ref{def:returns56}.
\end{proposition}
\begin{proof}
A length-$n$ product has angle $n\phi+j(\theta-\phi)$, $0\le j\le n$. The positive orbit of the irrational rotation $\theta-\phi$ is dense. Compactness supplies, for every accuracy, a finite initial orbit segment that is a net of the circle. All longer initial segments remain nets; rotation by $n\phi$ preserves their accuracy. Thus a product is arbitrarily close to the identity at every sufficiently large length. An exact return would be a nontrivial integer relation among $1,\theta/(2\pi),\phi/(2\pi)$.
\end{proof}
A single irrational letter fails the eventual condition: there is only one word at each horizon. Its target can be placed in a one-label horizon-dependent decoder even when no finite stationary approximation at the same tolerance exists. Thus mere recurrence along a subsequence is not a substitute for Definition~\ref{def:returns56}.

\begin{remark}[Quantitative interpretation]
The Ramsey estimate gives a finite occupation constant, not a useful
asymptotic width rate. The homogeneous neutral representative remains
inside the terminal decoder in the extraction proof. It is neither an
identity transition nor an assumed idempotent.
\end{remark}

\section{Stationary stochastic purification}\label{sec:purification55}
Write $\mathsf{Stoch}(k)$ for the compact semigroup of row-stochastic matrices of order $k$. The identity of a group inside this semigroup can be a singular idempotent; it must not be confused with the identity matrix.

\begin{lemma}[An idempotent over the physical identity]\label{lem:idempotent55}
Let $\mathcal T$ be a compact subsemigroup of $H\times\mathsf{Stoch}(k)$ whose projection onto $H$ is surjective. There is an element $e=(1,E)\in\mathcal T$ with $E^2=E$ and rank equal to the minimum matrix rank in $\mathcal T$. The monoid $e\mathcal T e$ is a compact group with identity $e$ and still projects onto $H$.
\end{lemma}
\begin{proof}
Only the integer ranks $1,\ldots,k$ occur, so their nonempty set has a minimum which is attained. Choose $(h,T)$ of that minimum rank $r$. The powers of $T$ are bounded, since they are stochastic. Thus its complex eigenvalues have modulus at most one and all Jordan blocks on the unit circle have size one: a larger peripheral Jordan block would produce polynomially growing powers. The closure of positive powers of an element of a compact group is a group: arbitrarily large positive powers approach its identity, and the preceding powers approach its inverse. Apply this observation to the compact monothetic group generated by $(h,\lambda_1,\ldots,\lambda_b)$ in $H\times(S^1)^b$, where the $\lambda_i$ are the peripheral eigenvalues of $T$. There is a net of integers tending to infinity for which $h^n\to1$ and all peripheral eigenvalues raised to the $n$th power tend to one. Along it $T^n$ tends to the spectral projection $E$ onto the peripheral subspace. This is a real stochastic idempotent. Its rank is at most $r$ and at least $r$ by minimality, so it is $r$. Nets avoid any metrizability assumption on $H$.

Every matrix $B$ in $e\mathcal T e$ satisfies $B=EBE$ and has rank $r$. Its action on the row space of $E$ is consequently invertible. Restriction embeds $e\mathcal T e$ as a compact submonoid of $H\times\mathrm{GL}(r,\R)$, with identity $(1,I_r)$; injectivity follows because the restriction determines $(xE)B=xB$ for every row $x$, by $B=EBE$, and hence determines every row of $B$. A compact submonoid of a topological group is a subgroup, by the same positive-power return argument. The inverse images of its inverses are still stochastic. Finally, sandwiching an element with physical coordinate $h$ between $e$'s leaves that coordinate unchanged. Surjectivity follows.
\end{proof}

\begin{lemma}[The recurrent simplex]\label{lem:simplex55}
If $E$ is a stochastic idempotent of rank $r$, then
\[
 \Delta_k E=\{p\in\Delta_k:pE=p\}
\]
is a simplex with exactly $r$ vertices. Every group of stochastic matrices with identity $E$ permutes these vertices.
\end{lemma}
\begin{proof}
Regard $E$ as a transition matrix. Its closed communicating classes each have a unique stationary probability row supported on that class. Every stationary row is a convex combination of these rows: mass on a transient class is zero, and stationarity within a closed irreducible class determines a one-dimensional space. For completeness, transient mass vanishes because each transient state has a positive probability of reaching a closed class in a bounded number of steps; a common bound over the finitely many states makes transient mass decrease by a fixed factor unless it is zero. The closed-class rows have disjoint supports, hence are linearly independent.

Each row of $E$ is stationary because $E^2=E$. Conversely, every stationary row equals its product with $E$. The span of the closed-class stationary rows is therefore precisely the row space of $E$, so there are $r$ classes and $r$ simplex vertices. If $B$ has a stochastic inverse relative to $E$, multiplication by $B$ and by this inverse are inverse affine maps of this simplex. Affine bijections permute its vertices.
\end{proof}

\begin{theorem}[Purification without a width or error loss]\label{thm:purification55}
For every compact-group interface in Definition~\ref{def:stationary55}, a stationary $k$-label realization of error at most $\varepsilon$ admits a permutation realization of error at most $\varepsilon$ on at most $k$ labels. There is no return-word hypothesis and no restriction on the number of connected components of $H$.
\end{theorem}
\begin{proof}
For the given machine form
\[
 \mathcal T=\overline{\{(u_w^{-1},T_w):w\in\mathcal A^*\}}
       \subset H\times\mathsf{Stoch}(k).
\]
The inverse on the physical coordinate is important: concatenation satisfies $u_{vw}^{-1}=u_v^{-1}u_w^{-1}$ and $T_{vw}=T_vT_w$. Thus $\mathcal T$ is a compact semigroup. Positive-word density makes its projection onto $H$ surjective. By continuity, its every element $(h,B)$ satisfies
\begin{equation}\label{eq:closederror55}
             \norm{\alpha_sBD_j-F_{s,j}(h^{-1})}_j\le\varepsilon
             \quad(s\in\mathcal S,\ j\in\mathcal J).
\end{equation}

Apply Lemma~\ref{lem:idempotent55} and denote the vertices of $\Delta_k E$ by $\nu_1,\ldots,\nu_r$. For each letter the element
\[
        e(u_a^{-1},T_a)e=(u_a^{-1},ET_aE)
\]
belongs to the group $e\mathcal T e$. By Lemma~\ref{lem:simplex55}, $B_a=ET_aE$ permutes the $\nu_i$. Use that permutation as the new command. If $V$ has rows $\nu_i$ and the induced row permutation is $\Pi_a$, the precise intertwining identity is $VB_a=\Pi_aV$. Write uniquely
\[
          \alpha_sE=\sum_{i=1}^r\beta_s(i)\nu_i,
          \qquad \widetilde D_j(i)=\nu_iD_j\in Y_j.
\]
The last inclusion uses convexity of the legal output set. For a nonempty word the output of the resulting $r$-label machine is $\alpha_s B_{a_1}\cdots B_{a_n}D_j$. For the empty word it is $\alpha_sED_j$. In both cases the corresponding joint element lies in $\mathcal T$ and has physical coordinate $u_w^{-1}$. Equation~\eqref{eq:closederror55} proves the error bound. Since $r\le k$, the width has not increased.
\end{proof}
The new machine need not reproduce the old machine's erroneous output word by word; it remains in the same closed error balls about the target. Nor do we assert $T_{w}=I$ on a physical return, invertibility of the original $T_a$, or deterministic initialization. These three incorrect shortcuts would invalidate the reduction.

\subsection{Finite quotients and the unrestricted stationary minimum}
For a tuple of permutation matrices $\sigma=(\Pi_a)_{a\in\mathcal A}$ of order $k$, set
\[
 \mathcal L_\sigma=
 \overline{\{(u_w^{-1},\Pi_w):w\in\mathcal A^*\}}
       \subset H\times\mathfrak S_k.
\]
The topology is the product of the compact topology on $H$ and the discrete topology on the finite group $\mathfrak S_k$. This is a compact subgroup. Its definition includes the possible discrepancy between physical relations and hidden-state relations.

\begin{theorem}[An exact stationary minimum]\label{thm:stationaryminimum55}
The value $M_\varepsilon$ is the least positive integer $k$ for which some tuple $\sigma$ of permutations of order $k$, some $\beta_s\in\Delta_k$, and some $d_j(i)\in Y_j$ satisfy
\begin{equation}\label{eq:permutationfeasibility55}
 \norm{\beta_s\Pi d_j-F_{s,j}(h^{-1})}_j\le\varepsilon
 \quad\text{for all }(h,\Pi)\in\mathcal L_\sigma\text{ and all }s,j.
\end{equation}
For fixed $k$ there are at most $(k!)^{|\mathcal A|}$ discrete choices (a crude bound before simultaneous conjugacy reduction), and the remaining feasible set is compact. The optimal stationary error at any fixed width is attained.
\end{theorem}
\begin{proof}
Necessity follows from Theorem~\ref{thm:purification55}, padding by unused labels if its output has fewer than $k$ states, and then taking the closure of word products. Sufficiency is the restriction of \eqref{eq:permutationfeasibility55} to words. The probabilities and legal decoders lie in a finite product of compact sets. Each displayed inequality defines a closed subset of that product. The worst error is continuous in these parameters, uniformly in the compact group variable, so its minimum exists. There are finitely many permutation tuples.
\end{proof}

\begin{theorem}[Finite observable quotient criterion]\label{thm:finitequotient55}
For arbitrary compact $H$ and the stated continuous interface,
\[
 M_\varepsilon<\infty
 \quad\Longleftrightarrow\quad
 r(U)\le\varepsilon\text{ for some open normal subgroup }U\le H.
\]
If such $U$ exists, $|\mathcal S|[H:U]$ deterministic-initialized permutation labels suffice. Conversely, a stationary $k$-label realization supplies such a $U$ with $[H:U]\le k!$.
\end{theorem}
\begin{proof}
Purify a stationary realization and use its joint group $\mathcal L_\sigma$. Put
\[
                  U=\{h\in H:(h,I)\in\mathcal L_\sigma\}.
\]
This is a closed normal subgroup of $H$: conjugate $(h,I)$ by an element above any $g\in H$, using surjectivity of the first projection. For a fixed permutation $\Pi$, a nonempty fiber in $\mathcal L_\sigma$ is a coset $hU$. These at most $k!$ fibers cover $H$, so $U$ has finite index. A closed subgroup of finite index in a compact group is open, since its complement is a finite union of closed cosets.

For each $s,j$ and each fiber, the single legal vector $\beta_s\Pi d_j$ is within $\varepsilon$ of every $F_{s,j}(h^{-1})$ in that fiber. Explicitly, $(hU)^{-1}=Uh^{-1}=h^{-1}U$ by normality, so inversion preserves the ordinary $U$-coset form. Hence $r(U)\le\varepsilon$. Conversely, store $(s,gU)$, update the coset by left multiplication by $u_a$, and decode a minimizing center of $F_{s,j}(gU)$. Compactness supplies the centers. This gives the asserted finite realization.
\end{proof}
For finitely many connected components, every open subgroup contains $K=H^\circ$, and $K$ itself is open normal. Therefore $M_\varepsilon<\infty$ holds exactly when
\[
 \varepsilon\ge\max_{s,j,c\in H/K}\rad_{Y_j}(F_{s,j}(c)),
\]
\emph{without} a return assumption. Executable returns enter the stronger clocked converse, not this stationary statement. The index bound $k!$ is an existence bound on a deterministic quotient, not a claim that the stationary minimum equals its index.

\begin{proposition}[Random initialization strictly saves states]\label{prop:C655}
Let $H=\Z/6\Z$, with an identity letter and a generator, one seed, and binary success probability
\[
 f(n)=\tfrac12+\tfrac14(-1)^n+\tfrac14\cos(2\pi n/3).
\]
At zero error the least deterministic component-quotient width is six, whereas the unrestricted stationary stochastic minimum is five.
\end{proposition}
\begin{proof}
The values at $n=0,\ldots,5$ are
\[
                  (1,\tfrac18,\tfrac58,\tfrac12,\tfrac58,\tfrac18).
\]
This sequence has least period six. A deterministic quotient of the cyclic component action preserving these outputs must therefore have six states. For a five-state realization take the disjoint union of a two-cycle and a three-cycle, both updating the index by $i\mapsto i+1$ modulo their respective lengths, initialize with mass $1/2$ at the zeroth vertex of each, and use decoder rows $(1,0)$ on the two-cycle and $(1,1/4,1/4)$ on the three-cycle. The identity command fixes all five labels. Direct averaging gives exactly $f(n)$, with legal binary probabilities.

If a stationary realization on at most four labels existed, purification would give one whose generator is a permutation of at most four letters. Its order is one of $1,2,3,4$, as follows by listing the possible cycle lengths. Every resulting numerical sequence has a period dividing that order, contradicting the least period six. Thus five is the unrestricted minimum.
\end{proof}

\section{Physical equivariance and the complete boundary}\label{sec:boundary56}
Return now to a compact Hausdorff group $H$ with a finite dense generating alphabet. The output sets may still be norm-compact convex subsets of real Banach spaces. Purification removes transition randomness; a further averaging removes hidden relations that have no physical counterpart.

\begin{theorem}[Physical-equivariant purification]\label{thm:equivariant56}
Every stationary $k$-label realization of error at most $\varepsilon$ admits one on $r\le k$ labels of the form
\begin{equation}\label{eq:equivariant56}
 \alpha_s\rho(h^{-1})d_j,\qquad h\in H,
 \qquad \rho:H\longrightarrow S_r\text{ a continuous homomorphism},
\end{equation}
with the same uniform error. Here permutations are represented by row-action matrices. In particular the command row is $\rho(u_a^{-1})$, all physical group relations hold on the new register, and initialization may remain random.
\end{theorem}
\begin{proof}
Apply Theorem~\ref{thm:purification55}, whose proof uses only finite convex combinations of decoder values. Denote the resulting initialization by $\beta_s$, its decoder by $D_j$, and its compact joint group by $\mathcal L\subset H\times S_k$, with the product topology and discrete permutation factor. Unused labels may be adjoined when necessary. Let $P$ be the second projection and
\[
                 N=\{\pi\in P:(1,\pi)\in\mathcal L\}.
\]
Then $N$ is normal in $P$. The averaging matrix $Q=|N|^{-1}\sum_{n\in N}n$ is stochastic, satisfies $Q^2=Q$, and commutes with $P$. For $(h,\pi)\in\mathcal L$, every $(h,n\pi)$ belongs to $\mathcal L$. Averaging their error inequalities and using convexity of the norm ball gives
\begin{equation}\label{eq:kernelaverage56}
                   \|\beta_sQ\pi D_j-F_{s,j}(h^{-1})\|_j\le\varepsilon.
\end{equation}

Let $O_1,\ldots,O_r$ be the $N$-orbits in the label set, and let $v_i$ be the uniform probability row on $O_i$. The rows of $Q$ are precisely these uniform orbit rows. Write $V$ for the matrix with rows $v_i$. Every $\pi\in P$ maps $N$-orbits bijectively onto $N$-orbits and therefore induces a permutation $\bar\pi$ with
\begin{equation}\label{eq:orbitintertwining56}
                             V\pi=\bar\pi V.
\end{equation}
The induced action of each $n\in N$ is trivial. Two elements of $\mathcal L$ with the same first coordinate differ by an element of $\{1\}\times N$. Hence $(h,\pi)\mapsto\bar\pi$ descends to a homomorphism $\rho:H\to S_r$. It is continuous: the surjection from the compact group $\mathcal L$ onto the Hausdorff group $H$ is a quotient map, and the original orbit action is continuous into a finite discrete group.

Set $\alpha_s(i)=\sum_{x\in O_i}\beta_s(x)$ and $d_j(i)=v_iD_j$. Then $\alpha_sV=\beta_sQ$ and $d_j=VD_j$. Equation~\eqref{eq:orbitintertwining56} gives
\[
                 \alpha_s\rho(h)d_j=\beta_sQ\pi D_j
                 \quad((h,\pi)\in\mathcal L).
\]
Replace $h$ by $h^{-1}$ in \eqref{eq:kernelaverage56}. This proves \eqref{eq:equivariant56}; all decoder values are legal by convexity, and $r\le k$.
\end{proof}

\begin{corollary}[Intrinsic stationary minimum]\label{cor:intrinsicminimum56}
The number $M_\varepsilon$ is the least cardinality of a finite continuous $H$-set for which some probability initialization for each seed and legal decoder for each query approximate the interface within $\varepsilon$. It is independent of the chosen finite dense alphabet. It is unchanged by replacing every target with
\[
                       h\longmapsto F_{s,j}(\varphi(h)q),
\]
where $\varphi$ is a continuous group automorphism and $q\in H$.
\end{corollary}
\begin{proof}
Theorem~\ref{thm:equivariant56} proves necessity, and restriction to command words proves sufficiency. For the last assertion, replace $\rho$ by $\rho\circ\varphi$ and $\alpha_s$ by $\alpha_s\rho(q^{-1})$ in \eqref{eq:equivariant56}. Applying the inverse change of variables proves equality of the minima.
\end{proof}
The finite quotient in Theorem~\ref{thm:finitequotient55} has a more precise description than its factorial bound. With $\mathcal L,P,N$ as above and $U=\{h:(h,I)\in\mathcal L\}$, the correspondence $(h,\pi)\mapsto(hU,\pi N)$ induces $H/U\simeq P/N$. Thus $[H:U]=|P/N|\le k!$. The orbit action in Theorem~\ref{thm:equivariant56} can have an additional kernel; no equality between this index and the minimal label count is asserted.

\begin{theorem}[Complete clocked finite-quotient classification]\label{thm:completeboundary56}
Assume eventual approximate returns on the compact group $H$. At every $\varepsilon\ge0$, including a positive critical radius, the following are equivalent:
\begin{enumerate}
\item $\sup_NW_{N,\varepsilon}<\infty$;
\item $M_\varepsilon<\infty$;
\item $r(U)\le\varepsilon$ for an open normal subgroup $U$ of $H$.
\end{enumerate}
For each $k<M_\varepsilon$ the all-profile occupation bound \eqref{eq:occupation56} holds, and $W_{N,\varepsilon}$ converges to $M_\varepsilon$. Put
\[
 R_0=\max_{s,j,h\in H}\rad_{Y_j}(F_{s,j}(hH^\circ)).
\]
Then $R_0=\inf_{U\lhd H\text{ open}}r(U)$. Above $R_0$ bounded realizations exist; below it every fixed width has bounded occupation. At $\varepsilon=R_0$ bounded clocked realization exists exactly when the infimum is attained by an open normal subgroup.
\end{theorem}
\begin{proof}
The first two conditions are equivalent by Theorem~\ref{thm:stationarization56}. The second and third are equivalent by Theorem~\ref{thm:finitequotient55}. That proof, including constrained center attainment, uses compactness and convexity of $Y_j$, not finite dimensionality.

For completeness, every open subgroup contains $H^\circ$: the image of the connected identity component in a finite discrete quotient is a singleton. Thus $r(U)\ge R_0$. The quotient $H/H^\circ$ is profinite, so its open normal subgroups form a neighborhood basis at the identity. The inverse images $U$ have intersection $H^\circ$. Given a neighborhood $O$ of $H^\circ$ in $H$, compactness supplies one such $U\subset O$: otherwise their intersections with the closed set $H\setminus O$ have the finite intersection property. Uniform continuity of the finite family of targets gives, for every $\eta>0$, a neighborhood $V$ of $1$ such that $\|F_{s,j}(hkv)-F_{s,j}(hk)\|_j<\eta$ for all $h\in H$, $k\in H^\circ$, and $v\in V$. Choose $U\subset H^\circ V$. Every $F_{s,j}(hU)$ is then within $\eta$ of $F_{s,j}(hH^\circ)$, so $r(U)\le R_0+\eta$. This proves the infimum formula and all remaining assertions.
\end{proof}
This result includes the formerly separate zero and positive boundary cases. The finite-rank proof of Theorem~\ref{thm:zero55} is retained below because it identifies an additional exact translate-space obstruction. No passage from an error strictly larger than $R_0$ to equality is used in Theorem~\ref{thm:completeboundary56}.

\subsection{Length phases without a width overhead}
Let $p\ge1$ be an integer and fix a command with product $a$. Let $J$ be the closure of all products of lengths divisible by $p$. The group argument of Lemma~\ref{lem:phasegroup55}, which uses only the block length, shows that $J$ is open normal, $H/J$ is cyclic of order dividing $p$, and all letters belong to $Ja$. Regard $\mathcal A^p$ as a finite alphabet on $J$. Assume that this block alphabet has eventual approximate returns. In particular this holds when the original nonempty exact-return language has greatest common divisor $p$.

For $0\le r<p$ let $M_{r,\varepsilon}$ be the intrinsic stationary minimum on $J$ for
\[
                         f_{r,s,j}(h)=F_{s,j}(ha^r).
\]
\begin{theorem}[Exact limiting width in each phase]\label{thm:phaseminimum56}
Under the preceding assumptions, for every $r$ and $\varepsilon\ge0$,
\begin{equation}\label{eq:phaseminimum56}
 W_{N,\varepsilon}\le M_{r,\varepsilon}\quad(N\equiv r\pmod p),
 \qquad
 \lim_{\substack{N\to\infty\\N\equiv r\ ({\rm mod}\ p)}}W_{N,\varepsilon}
          =M_{r,\varepsilon}.
\end{equation}
Finite limits are eventually attained. If $k<M_{r,\varepsilon}$, the number of positive cuts of width at most $k$ is bounded uniformly over that horizon phase, with no restriction on other widths. Boundedness in the phase is equivalent to the existence of an open normal subgroup $U\le J$ such that
\begin{equation}\label{eq:phasequotient56}
          \max_{s,j,h\in J}\rad_{Y_j}(F_{s,j}(hUa^r))\le\varepsilon.
\end{equation}
\end{theorem}
\begin{proof}
Suppose first that $M_{r,\varepsilon}=K<\infty$ and choose a physical-equivariant stationary realization of $f_r$ on $K$ labels, with representation $\rho:J\to S_K$. For a fixed $N\equiv r\pmod p$, let $g_t$ be the physical product after $t$ commands and set
\begin{equation}\label{eq:gauge56}
 z_t=a^{N-t}g_ta^{-r}\in J,
 \qquad c_{t,b}=a^{N-t}u_ba^{-(N-t+1)}\in J.
\end{equation}
Normality of $J$ and $u_ba^{-1}\in J$ justify the membership assertions. We have $z_0=a^{N-r}$, $z_t=c_{t,b}z_{t-1}$ when the $t$th letter is $b$, and $z_N=g_Na^{-r}$. Initialize with $\alpha_s\rho(z_0^{-1})$, use the clock-dependent permutation $\rho(c_{t,b}^{-1})$, and keep decoder $d_j$. The row product is $\alpha_s\rho(z_N^{-1})$, so its error from $F_{s,j}(g_N)=f_{r,s,j}(z_N)$ is at most $\varepsilon$. This uses the free external clock but no phase register or partial-block storage. It proves the upper bound with exactly $K$ labels.

Fix $k<M_{r,\varepsilon}$ and a cut residue $0\le\ell<p$. Put $b=(r-\ell)\bmod p$, chosen in $\{0,\ldots,p-1\}$. Restrict the input to the fixed prefix $a^\ell$, arbitrary $p$-letter blocks, and fixed suffix $a^b$. Integrate the prefix into initialization and the suffix into the decoder. The block target on $J$ is
\[
 G_{\ell,s,j}(h)=F_{s,j}(a^b h a^\ell)
     =f_{r,s,j}\bigl(\operatorname{Ad}_{a^b}(h)a^{\ell+b-r}\bigr).
\]
Here $a^{\ell+b-r}\in J$ and conjugation by $a^b$ is an automorphism of $J$. Corollary~\ref{cor:intrinsicminimum56} therefore identifies its stationary minimum with $M_{r,\varepsilon}$. Theorem~\ref{thm:stationarization56} on the block alphabet bounds its positive narrow cuts by a constant $L_\ell$. These cuts are precisely the original cuts in residue $\ell$ between the prefix and suffix, apart from the initial block cut. At most  $\ell+b+1\le2p$ further cuts need be charged for this residue, including small horizons shorter than the chosen prefix and suffix. Summing gives the uniform bound $\sum_{\ell=0}^{p-1}L_\ell+2p^2$. As before, a peak at most $k$ would force $N$ below this bound. This proves the limit and occupation assertions. Finally apply Theorem~\ref{thm:finitequotient55} to $f_r$ on $J$ to obtain \eqref{eq:phasequotient56}.
\end{proof}
If the order of $H/J$ is a proper divisor $q$ of $p$, the minima and quotient-radius criteria repeat with period $q$: multiplication by $a^q\in J$ is a right translation of the target. The earlier phase radii have the same repetition. Neither the physical quotient order nor the return-length greatest common divisor is substituted silently for the other.

\subsection{Compact outputs and finite hidden state}
Theorems~\ref{thm:stationarization56}, \ref{thm:equivariant56}, \ref{thm:completeboundary56}, and \ref{thm:phaseminimum56} hold for norm-compact convex Banach-valued outputs. A single-machine purification even needs only convexity: its new decoder values lie in the convex hull of finitely many old values. Compact output sets are used instead in finite-witness extraction, optimal-error attainment, and constrained-center existence. The representative-function and algebraic sections below retain their explicitly finite-dimensional hypotheses.

Finiteness on the hidden side has a different role. On countably many labels the full simplex is not compact in total variation; mass can escape to infinity, matrix ranks need not admit a useful finite minimum, and a recurrent simplex may have infinitely many vertices. On a general compact convex hidden state space affine automorphisms need not permute a finite vertex set: rotation of the disk already acts continuously on infinitely many extreme points. Thus these proofs give neither a finite-label purification theorem for arbitrary operator kernels nor an online observation theorem. Their terminal-output assertion is strictly the one stated in the models.

\begin{remark}[Action and counting conventions]
Probability vectors act on the right of transition matrices. Chronological
row multiplication consequently corresponds to the inverse physical
coordinate $h^{-1}$ used above. The coarse deterministic quotient
construction stores both the seed and the coset; its
$|\mathcal S|[H:U]$ bound is not the unrestricted stochastic minimum.
In the phase theorem the chosen block period $p$, the physical quotient
order $q\mid p$, and the greatest common divisor of actual return
lengths retain their separate meanings.
\end{remark}

\section{The clocked component-radius threshold}\label{sec:intro54}
A finite stochastic register can carry infinitely many probability distributions. Thus an infinite numerical orbit is not, by itself, an obstruction to finite positive realization. A further difficulty arises when the realization may be redesigned at every horizon and its transitions may depend on a free external clock. The obstruction established here survives both freedoms. It is measured by the radius of an observable image inside its allowed output set.

For the clocked theorem in this section, let $H$ be a compact Hausdorff group with finitely many connected components. Put $K=H^\circ$ and $C=H/K$. Then $K$ is an open normal subgroup and $C$ is a finite group. Fix a finite alphabet $\mathcal A$ with elements $u_a\in H$ whose generated group is dense in $H$. Our convention is
\[
 u_{a_1\cdots a_n}=u_{a_n}\cdots u_{a_1}.
\]
The identity $e$ need not be an alphabet element. Instead assume the following executable return condition:
\begin{equation}\label{eq:returns54}
 \text{there is an integer }g\ge0\text{ such that, for every }n\ge g,
 \text{ some word of length }n\text{ has product }e.
\end{equation}
An identity letter gives $g=0$. The condition concerns physical products; the stochastic rows on a return word may perform arbitrary processing.

For nonempty finite seed and query sets $\mathcal S,\mathcal J$, let
\[
 F_{s,j}:H\longrightarrow Y_j
\]
be continuous, where $Y_j$ is a nonempty compact convex subset of a finite-dimensional real normed space $(V_j,\norm{\cdot}_j)$. A query is selected only after the command word. A prescribed joint law of finitely many outputs is treated as one query, not as a collection of marginal constraints. In particular, a categorical interface has
\[
 Y_j=\Delta_{M_j}=\{p\in\R^{M_j}:p_i\ge0,\ \textstyle\sum_i p_i=1\},
 \qquad \norm{x}_j=\tfrac12\sum_i|x_i|.
\]
Then $\norm{p-q}_j$ is total variation. Binary means correspond to $Y=[-1,1]$ with $\norm{x}=|x|/2$.

\begin{definition}[Clocked realization]\label{def:machine54}
At horizon $N$ choose finite nonempty registers $S_t$, $0\le t\le N$, initialization rows $E_s$, row-stochastic command matrices $T_{t,a}:S_{t-1}\to S_t$, and terminal decoder maps $D_j:S_N\to Y_j$. The realized output vector on $w=a_1\cdots a_N$ is
\[
 E_sT_{1,a_1}\cdots T_{N,a_N}D_j.
\]
Its error is the maximum, over all seeds, length-$N$ words, and queries, of its distance from $F_{s,j}(u_w)$. The command width is $\max_t|S_t|$, counting all declared labels, including unused padding labels. Write $W_{N,\varepsilon}$ for its least possible value at error at most $\varepsilon$.
\end{definition}
All retained private information, including randomness kept for future use, is charged to the current label. The rows may depend on $t$ and $N$, but not on an unrecorded input history or on the future query. The external clock and horizon, arbitrary real table descriptions, their construction and lookup, real arithmetic, and exact sampling are uncharged. This is a numerical realization invariant, not uniform workspace or program size. For a finite categorical interface, charging the terminal answer register changes the peak to $\max(W_{N,\varepsilon},\max_j M_j)$. Thus all boundedness and growth results are unchanged; the convention that also charges a sampled binary-answer register has a minimum of two labels. A decoder vector by itself is not an additional sampled register.

For $A\subset Y_j$ define its constrained radius by
\[
 \rad_{Y_j}(A)=\min_{y\in Y_j}\sup_{x\in A}\norm{x-y}_j.
\]
The minimum is attained when $A$ is compact. The critical error is
\begin{equation}\label{eq:radius54}
 \epsc=\max_{s,j,c\in C}\rad_{Y_j}\bigl(F_{s,j}(c)\bigr).
\end{equation}
Here a component is identified with its coset in $H$; its image is compact by continuity. The requirement that the center belong to $Y_j$ is essential for categorical outputs.

\begin{theorem}[Intrinsic radius classification]\label{thm:radius54}
Under \eqref{eq:returns54}, for every $\varepsilon\ge0$ the following are equivalent:
\begin{enumerate}
\item $\sup_N W_{N,\varepsilon}<\infty$;
\item one stationary finite-state realization with permutation command updates has error at most $\varepsilon$ at every horizon;
\item $\varepsilon\ge\epsc$.
\end{enumerate}
At the boundary, at most $|\mathcal S||C|$ command labels suffice. If $\varepsilon<\epsc$, then for every integer $k\ge1$ there is an integer $L_k<\infty$ such that every error-$\varepsilon$ realization satisfies
\begin{equation}\label{eq:occupation54}
 \#\{t\in\{1,\ldots,N\}:|S_t|\le k\}\le L_k.
\end{equation}
All intervening widths are unrestricted. In particular, $W_{N,\varepsilon}\to\infty$.
\end{theorem}
Neither a positive lower bound on state masses nor a reachable-rank, observability, or conditioning assumption occurs in the theorem. The upper realization is deterministic until its terminal decoder. It need not have the minimum width among all stochastic realizations. Theorem~\ref{thm:stationaryminimum55}, rather than a deterministic component partition, characterizes that stationary minimum.

For binary means, \eqref{eq:radius54} reduces to one quarter of the largest componentwise oscillation. For a genuine categorical law it is generally not half its diameter; an analytic three-outcome example below has critical total-variation error $2/3$, whereas every individual event has binary threshold $1/2$. The proof uses finitely many localized vector tests rather than a separate midpoint for each coordinate.

The finite-component assumption does not assert that $H$ is a Lie group. For example, take $H=\prod_{n\ge1}\R/\Z$, choose an element whose every finite coordinate tuple together with $1$ is rationally independent, and include this element and $e$ as commands. The positive orbit is dense and $H$ is connected but has no faithful finite-dimensional representation. The theorem applies to every finite continuous interface on this group. With infinitely many components the component construction need not be finite; Section~\ref{sec:profinite55} gives the finite-observable-quotient criterion and the strict-side clocked theorem in that setting.

There is also a quantitative counterpart. For planar rotations with a dual badly approximable $s$-tuple of angles and signal amplitude $0<\rho<1$, Theorem~\ref{thm:matched54} proves
\[
 W_{N,\varepsilon}=\Theta\bigl(N^{s/(2s+1)}\bigr)
 \quad\text{for every fixed }0\le\varepsilon<\rho/2.
\]
The upper realization is exact even when $\varepsilon>0$. The new lower estimate reaches the entire subcritical interval, by combining harmonic localization with the arithmetic orbit-packing mechanism. This quantitative theorem and the radius classification have distinct hypotheses: algebraic matrix entries alone do not imply the badly approximable angle condition.

\subsection{Executable tests}
Let $\mathcal P=\{u_w:w\in\mathcal A^*\}$. We use standard Haar averaging and the density of matrix coefficients of finite-dimensional unitary representations, in realified form \cite{Folland}. No quantitative mixing assumption for a preassigned random walk is imposed.

\begin{lemma}[Positive words and common lengths]\label{lem:positive54}
The semigroup $\mathcal P$ is dense in $H$, and $\mathcal P\cap K$ is dense in $K$. Every finite subset of $\mathcal P\cap K$ has executable representatives of a common positive length.
\end{lemma}
\begin{proof}
For an element $u$ of a compact group, positive powers return arbitrarily close to $e$ with exponent at least two. Indeed cover the compact closure of its powers by finitely many translates of a sufficiently small neighborhood. Among sufficiently many powers, three lie in one member of the cover; the first and last exponents differ by at least two, and their quotient is arbitrarily close to $e$. This argument also covers finite-order elements by taking a sufficiently large multiple of the order. If $u^r$ approaches $e$ with $r\ge2$, then $u^{r-1}=u^ru^{-1}$ approaches $u^{-1}$ and is still a positive power. Thus $u^{-1}$ is in the closure of positive powers. The closed semigroup $\overline{\mathcal P}$ therefore contains the group generated by the letters, and hence equals $H$. Since $H$ has finitely many components, $K$ is open, and a relatively open subset of $K$ is open in $H$; density gives the second assertion. For finitely many representative lengths $\ell_i$, choose $B\ge\max_i\ell_i+g$, $B\ge1$, and append return words of lengths $B-\ell_i$. The products are unchanged.
\end{proof}

\begin{lemma}[Finite executable gap]\label{lem:gap54}
Let $R:K\to O(E)$ be a continuous representation on a nonzero finite-dimensional Euclidean space, with no nonzero fixed vector. For every $k\ge1$ there are $B\ge1$, a probability law $\nu$ on finitely many executable length-$B$ words with products in $K$, and $d\in(0,1)$ such that
\begin{equation}\label{eq:gap54}
 \E_{w\sim\nu}\max_{1\le i\le k}\ip{v_i}{R(u_w)u}\le1-d
\end{equation}
for all unit vectors $u,v_1,\ldots,v_k\in E$.
\end{lemma}
\begin{proof}
On the compact product of unit spheres consider the minimum
\[
 \gamma=\min_{u,v_1,\ldots,v_k}\int_K
 \left(1-\max_i\ip{v_i}{R(h)u}\right)\,d\mu(h).
\]
Repeated centers are permitted. The integrand is nonnegative. If a minimizing tuple had integral zero, full support of Haar measure would give $R(h)u\in\{v_1,\ldots,v_k\}$ for every $h$. The orbit of $u$ is connected, hence a singleton, contradicting the fixed-vector hypothesis. Thus $\gamma>0$.

Partition the compact image $R(K)$ into finitely many measurable sets of small diameter in the Euclidean operator norm, and replace their Haar masses by atoms at nearby executable products in $K$. Cells of zero Haar mass can be discarded. The function $A\mapsto\max_i\ip{v_i}{Au}$ is uniformly $1$-Lipschitz in that norm. An approximation radius less than $\gamma/2$ therefore preserves a positive defect uniformly in the unit vectors. Lemma~\ref{lem:positive54} gives a common length. Decreasing the defect to a number in $(0,1)$ proves the assertion.
\end{proof}

\begin{lemma}[Conditional-centroid contraction]\label{lem:centroid54}
Let $Z\in E$ be a bounded random physical vector and $S$ the current register. Apply an independent word $w\sim\nu$ from Lemma~\ref{lem:gap54}, and let the actual block kernel generate a terminal state $S'$ having at most $k$ labels. Then
\begin{equation}\label{eq:centroid54}
 \E\norm{\E[R(u_w)Z\mid S']}\le(1-d)\E\norm{\E[Z\mid S]}.
\end{equation}
A fixed intervening return word cannot increase the same conditional norm, regardless of its intermediate or terminal widths.
\end{lemma}
\begin{proof}
Put $z_s=\E[Z\mid S=s]$ and write $T_w(s,i)$ for the actual composite row of the block. Since this row depends on the past only through $S$,
\[
 \Pr(S'=i)\E[R(u_w)Z\mid S'=i]
 =\sum_s\Pr(S=s)\E_w T_w(s,i)R(u_w)z_s.
\]
Choose $v_i$ in the direction of each nonzero terminal centroid, arbitrarily otherwise. Pair, sum in $i$, and use stochasticity to bound the sum by
\[
 \sum_s\Pr(S=s)\E_w\max_i\ip{v_i}{R(u_w)z_s}
 \le(1-d)\sum_s\Pr(S=s)\norm{z_s}.
\]
Zero-probability states contribute zero. No restriction was placed on registers inside the block. On a fixed return word the physical vector is unchanged at its endpoints; the new state is a stochastic image of the old one, so conditional Jensen gives nonincrease. Independence is used only at complete block boundaries.
\end{proof}

\section{Representative functions and localized radii}\label{sec:local54}
A real representative function on $K$ is a matrix coefficient of a finite-dimensional continuous orthogonal representation, or a finite sum of such coefficients. Their algebra contains the constants, is closed under multiplication by tensor products, and is uniformly dense in $C(K,\R)$ by the Peter--Weyl approximation theorem \cite{Folland}. A finite collection of representative functions can be written on one space
\[
 \mathcal V=\bigoplus_{i=1}^a\End(E_i),\qquad
 \Phi(h)=(R_i(h))_{i=1}^a,
\]
where the scalar product is the sum of Frobenius products. Left multiplication defines an orthogonal representation $R$ on $\mathcal V$ and satisfies $\Phi(gh)=R(g)\Phi(h)$. Adjoin a trivial summand when necessary. Let
\[
 P=\int_KR(h)\,d\mu(h),\quad E=(I-P)\mathcal V,\quad
 Z(h)=(I-P)\Phi(h),\quad Q_0=\norm{Z(e)}.
\]
Haar averaging is the orthogonal projection onto fixed vectors. Consequently $Z(gh)=R(g)Z(h)$, the representation on $E$ has no fixed vectors, and $P\Phi(h)=P\Phi(e)$ for $h\in K$.

Choose and fix a coefficient vector $a_p$ for each scalar representative function $p$, and a linear coefficient map $A_G$ for each vector-valued representative function $G$. Their norms need not be intrinsic; fixing any such choices suffices. For an arbitrary joint distribution of $h\in K$ and a finite state $S$, put $Q=\E\norm{\E[Z(h)\mid S]}$ and $\bar p=\int p\,d\mu$. The fixed/moving decomposition gives
\begin{align}
 \E[p(h)\mid S]&=\bar p+\ip{a_p}{\E[Z(h)\mid S]},\label{eq:calibration54}\\
 |\E p-\bar p|&\le\norm{a_p}Q,\qquad
 \norm{\E G-\bar G}\le\norm{A_G}Q.\label{eq:vectorcal54}
\end{align}
Here the norm of $A_G$ is the operator norm from the Euclidean feature space to the indicated output space. If $D:S\to Y$ and $M_Y=\max_{y\in Y}\norm{y}$, then
\begin{equation}\label{eq:decodercal54}
 \norm{\E[p(h)D(S)]-\bar p\E D(S)}\le M_Y\norm{a_p}Q.
\end{equation}
Indeed condition on $S$ in \eqref{eq:calibration54}, multiply by $D(S)$, and apply the triangle inequality. These are moment estimates, not total-variation convergence of a hidden physical distribution.

\begin{lemma}[Finite localization of a radius]\label{lem:radiuslocal54}
Let $f:K\to Y$ be continuous, with $Y$ compact and convex in a finite-dimensional normed space. Suppose $\rad_Y(f(K))>\varepsilon$. There are a vector-valued representative function $G$, a number $\eta>0$, and nonnegative representative functions $p_1,\ldots,p_a$ of Haar mean one such that
\[
 \norm{G-f}_\infty<\eta,\qquad
 r_*:=\min_{y\in Y}\max_i\norm{\bar{p_iG}-y}>\delta:=\varepsilon+\eta.
\]
\end{lemma}
\begin{proof}
Choose $\varepsilon<t<\rad_Y(f(K))$. For each $y\in Y$ some $h\in K$ has $\norm{f(h)-y}>t$. The corresponding open subsets of $Y$ cover $Y$, so a finite subcover yields $h_1,\ldots,h_a$ with
\[
 \min_{y\in Y}\max_i\norm{f(h_i)-y}>t.
\]
The strict inequality follows also at a minimizing $y$, since the maximum is everywhere greater than $t$ and $Y$ is compact.

Fix a sufficiently small positive $\eta$. Approximate the coordinates of $f$ uniformly by representative functions to obtain $\norm{G-f}_\infty<\eta$. For each $i$, choose a nonzero nonnegative continuous bump supported in a neighborhood where $f$ differs from $f(h_i)$ by less than $\eta$. Such a bump exists by normality of the compact Hausdorff space; it has positive Haar integral by full support. Approximate its square root uniformly by a representative function, square the approximation, and normalize its Haar integral. The resulting nonnegative representative function $p_i$ has mean one, and its weighted $f$-mean can be made arbitrarily close to that of the bump. Thus we may ensure
\[
 \norm{\bar{p_iG}-f(h_i)}<3\eta.
\]
The constrained radius changes by at most the maximum perturbation of the finitely many points. Taking $4\eta<t-\varepsilon$ gives $r_*>\varepsilon+\eta$.
\end{proof}

\begin{proposition}[Residual moment from a vector interface]\label{prop:residual54}
Fix the data of Lemma~\ref{lem:radiuslocal54}. Choose one feature space representing all $p_i$ and $p_iG$, and fix their coefficient choices before defining
\[
 L=\max_i\bigl(\norm{A_{p_iG}}+(M_Y+\delta)\norm{a_{p_i}}\bigr),
 \qquad \Delta=r_*-\delta>0.
\]
Suppose a terminal decoder $D(S)\in Y$ obeys
\[
 \norm{\E[D(S)\mid w]-G(h(w))}\le\delta
\]
for every complete deterministic word in a test law. Then $L>0$, $E\ne\{0\}$, $Q_0>0$, and
\begin{equation}\label{eq:residual54}
                         Q\ge\Delta/L.
\end{equation}
\end{proposition}
\begin{proof}
For $p=p_i\ge0$, wordwise correctness and the triangle inequality imply
\[
 \norm{\E[p(h)D(S)]-\E[p(h)G(h)]}\le\delta\E p(h).
\]
This is a conditional mean error, not a pointwise bound on the random decoder. Since $\bar p_i=1$, calibration gives explicitly
\[
                  \E p_i\le1+\norm{a_{p_i}}Q.
\]
Write $y=\E D(S)$, which belongs to $Y$ by convexity. Equations~\eqref{eq:vectorcal54}--\eqref{eq:decodercal54}, together with $\bar p=1$, give
\[
 \norm{y-\bar{p_iG}}\le\delta+
 \bigl(\norm{A_{p_iG}}+(M_Y+\delta)\norm{a_{p_i}}\bigr)Q.
\]
Maximizing in $i$ yields $r_*\le\delta+LQ$. Hence $LQ\ge\Delta>0$. In particular the moving feature and its initial norm cannot vanish: by equivariance, $Q_0=0$ would imply $Z(h)=0$ for all $h$ and therefore $Q=0$.
\end{proof}

\subsection{Proof of Theorem~\ref{thm:radius54}}
Choose $(s,j,c)$ with component radius greater than $\varepsilon$. By density there is an executable prefix $v$ with $u_v\in c$. Put $b=|v|$ and $f(h)=F_{s,j}(hu_v)$ on $K$. Since $Ku_v=c$, this map has the same constrained radius. Apply Lemma~\ref{lem:radiuslocal54} and Proposition~\ref{prop:residual54}.

The moving feature is nonzero independently of any particular test law. Otherwise all represented functions would be constant on $K$; since each $p_i$ has mean one, this would force $p_i=1$ and $G$ constant. Then all $\bar{p_iG}$ would equal a point at distance less than $\eta$ from $Y$, contradicting $r_*>\delta\ge\eta$.

For a fixed width $k$, apply Lemma~\ref{lem:gap54} to the resulting moving feature. Let $B,d$ be its block parameters and set
\[
 \chi=-\log(1-d)>0,\qquad A=\max(1,LQ_0/\Delta).
\]
Start a test word with $v$; between independent length-$B$ test blocks use fixed return words. The accumulated physical product is $hu_v$. Initially $h=e$, so the conditional feature norm after the prefix equals $Q_0$, regardless of the hidden-state distribution. If $J$ block endpoints have width at most $k$, Lemma~\ref{lem:centroid54} gives
\[
 Q_N\le Q_0(1-d)^J.
\]
The machine is correct on every deterministic word of length $N$. Thus any external law supported on such words inherits the conditional mean bound, with error $\delta=\varepsilon+\eta$ for $G$. Proposition~\ref{prop:residual54} implies
\begin{equation}\label{eq:J54}
                         J\chi\le\log A.
\end{equation}
The law depends on the advertised width profile, not on the machine's private state. When $A=1$, \eqref{eq:J54} excludes even one selected narrow block; no logarithmic loss has occurred.

For clarity we give the spacing count with the full return hypothesis. Discard narrow cuts before $b+B+g$ and after $N-g$. From the remaining cuts greedily choose the first and then the next at distance at least $B+g$. Put an independent block immediately before each selected cut. The initial gap, all intervening gaps, and the final suffix have lengths at least $g$, so \eqref{eq:returns54} fills them with executable returns. At most $b+B+2g-1$ cuts were discarded. Each chosen endpoint accounts for at most $B+g$ eligible cuts. Therefore one may take
\begin{equation}\label{eq:Lk54}
 L_k=\left\lceil b+B+2g-1+(B+g)\frac{\log A}{\chi}\right\rceil.
\end{equation}
The constants $B,d$ depend on the fixed moving representation and $k$; $\chi=-\log(1-d)$ depends on that defect. The localization, prefix length $b$, and $A$ depend only on the alphabet, interface, chosen component, and error. Consequently $L_k$ depends on these fixed data, the return certificate $g$, and $k$, but not on $N$ or the machine. When the eligible interval is empty the same estimate follows directly by counting. This proves occupation without restricting any register inside a block or gap. If the entire peak is at most $k$, then $N\le L_k$, which proves divergence directly. Monotonicity in the horizon is not needed when an identity letter is absent.

For the upper bound, use labels $(s,c)\in\mathcal S\times C$, including any labels not reached at a given horizon. Initialize at $(s,K)$ and update by left multiplication
\[
 (s,c)\longmapsto(s,[u_a]c).
\]
Choose a minimizing center of $F_{s,j}(c)$ in $Y_j$ as decoder at $(s,c)$ for query $j$. These are legal decoder values, and the error is at most \eqref{eq:radius54}. All command updates are permutations of one finite state set, independent of clock and horizon. This proves boundary attainment and the equivalence.\qed

\begin{proposition}[A finite return certificate]\label{prop:return54}
Condition~\eqref{eq:returns54} follows if there are finitely many executable identity-product words whose positive lengths have greatest common divisor one.
\end{proposition}
\begin{proof}
Let the lengths be $\ell_1,\ldots,\ell_a$. Their nonnegative sums modulo $\ell_1$ form the subgroup generated by their residues: a subsemigroup of a finite group is a subgroup. The gcd hypothesis makes this the whole residue group. Choose one nonnegative sum $n_r$ for each residue $r$ modulo $\ell_1$, and put $g=\max_r n_r$. Every $n\ge g$ equals its representative $n_r$ plus a nonnegative multiple of $\ell_1$. Concatenating the corresponding return words proves the claim.
\end{proof}
For example, the alphabet $\{R_\theta,R_{-\theta},R_{2\pi/3}\}$, with $\theta/2\pi$ irrational, contains no identity letter but has return lengths two and three. It satisfies \eqref{eq:returns54} with $g=2$. A single irrational rotation letter does not: there is only one word at each horizon, and a horizon-specific constant decoder can store its numerical answer with one command label. Thus synchronization cannot be silently omitted. Nor is a bounded irreversible semigroup enough: for the identity command $I=1$ and contraction command $a=\lambda\in(0,1)$, the mean $\rho\lambda^{\#a(w)}$ has an exact alive/absorbed two-state realization. The group generated with inverses is not bounded. The present theorem is not a compact-semigroup classification.

\section{Exact-return languages and length phases}\label{sec:phases55}
The existence of cofinite returns is not the only useful synchronization regime. In fact an exact-return language, when nonempty, has a rigid eventual form. Here and in this section a return has positive length.

\begin{lemma}[The return period]\label{lem:returnperiod55}
Suppose
\[
 \mathcal R=\{|w|>0:u_w=1\}\ne\varnothing,
 \qquad p=\gcd\mathcal R.
\]
There is an integer $g_p$ such that every multiple $dn$ with $n\ge g_p$ belongs to $\mathcal R$, and no length outside $p\N$ belongs to $\mathcal R$. In particular, every nonempty exact-return language is eventually one arithmetic progression.
\end{lemma}
\begin{proof}
Concatenation makes $\mathcal R$ additive. A finite subset of its elements already has gcd $p$: start with one length and successively decrease the gcd when necessary; a strictly decreasing chain of its positive divisors is finite. Divide these finitely many lengths by $p$. Their gcd is one, so the residue argument of Proposition~\ref{prop:return54} shows that their nonnegative sums contain every sufficiently large integer. Multiplying by $p$ proves the assertion. The exclusion of other lengths is the definition of the gcd.
\end{proof}

Fix a letter with product $a\in H$, and let $J$ be the closure of the products of words whose lengths are multiples of $p$, including the empty word.
\begin{lemma}[The phase subgroup]\label{lem:phasegroup55}
The group $J$ is an open normal subgroup of $H$. The quotient $H/J$ is cyclic of order dividing $p$, generated by $aJ$, and
\begin{equation}\label{eq:phasereachable55}
 \overline{\{u_w:|w|\equiv r\pmod p\}}=a^rJ
 \qquad(0\le r<p).
\end{equation}
If $H/H^\circ$ is finite, then $J^\circ=H^\circ$.
\end{lemma}
\begin{proof}
The compact closure defining $J$ is a group, by the positive-power argument, and is generated densely by the block alphabet $\mathcal A^p$. Also $a^p\in J$. If $x=u_w$ for a word $w$ of length $p$, then the executable word consisting of $p-1$ copies of the chosen letter, followed by $w$, followed by one more copy, has product $axa^{p-1}$ and length $2p$. Thus $axa^{p-1}\in J$, whence $axa^{-1}\in J$. It follows that $aJa^{-1}\subset J$. Applying this inclusion $p$ times and using $a^p\in J$ shows equality. Every letter $b$ has $ba^{p-1}\in J$, so $ba^{-1}\in J$. Thus all letters lie in $Ja$, normalize $J$, and their products of length $n$ lie in $Ja^n$. The finite union $\bigcup_{r=0}^{p-1}Ja^r$ is a closed group containing the letters, hence is $H$. This proves normality, finite index, and openness. Conversely, words consisting of $r$ copies of $a$ followed by arbitrary $p$-blocks have products dense in $Ja^r$, proving \eqref{eq:phasereachable55}. An open subgroup contains the identity component, and its own identity component is therefore the same one.
\end{proof}
The quotient order may be a proper divisor of $p$; the phase radii repeat whenever the physical cosets repeat. The phase index records executable length, not merely the physical component quotient.

Assume now that $H$ has finitely many components and put $K=H^\circ$. For $r\in\{0,\ldots,p-1\}$ define
\begin{equation}\label{eq:phaseradius55}
 \varepsilon_r=\max_{s,j,\ c\in H/K:\ c\subset a^rJ}
                    \rad_{Y_j}(F_{s,j}(c)).
\end{equation}
\begin{theorem}[The phasewise clocked threshold]\label{thm:phases55}
Under the sole synchronization assumption $\mathcal R\ne\varnothing$, for every $r$ and $\varepsilon\ge0$,
\[
 \sup_{N\equiv r\ (\mathrm{mod}\ p)}W_{N,\varepsilon}<\infty
             \quad\Longleftrightarrow\quad
             \varepsilon\ge\varepsilon_r.
\]
Equality is attained by a finite component-permutation realization on the indicated horizons. If $\varepsilon<\varepsilon_r$, then for each $k$ there is a constant $L_{k,r}$, independent of $N$, such that every error-$\varepsilon$ realization at $N\equiv r\pmod p$ satisfies
\[
                 \#\{1\le t\le N:|S_t|\le k\}\le L_{k,r}.
\]
The intervening widths are unrestricted. Consequently the full family is bounded exactly at and above $\max_r\varepsilon_r=\varepsilon_{\mathrm c}$; below it, divergence is asserted phase by phase, not necessarily along every horizon.
\end{theorem}
\begin{proof}
Every product of length $N\equiv r$ lies in $a^rJ$. Store its component and the seed and decode component centers as before. This proves the upper bound, including equality, with at most $|\mathcal S||H/K|$ labels. Centers on other components may be chosen arbitrarily when only this phase is under consideration.

For the converse fix a component $c\subset a^rJ$ and $s,j$ with radius greater than $\varepsilon$. We must count narrow cuts in every residue class, not only at block boundaries of one phase. For each $\ell\in\{0,\ldots,p-1\}$ choose $b_\ell\in\{0,\ldots,p-1\}$ congruent to $r-\ell$ modulo $p$. By \eqref{eq:phasereachable55} and the openness of components there is a fixed word $v_\ell$, of length $p_\ell\equiv\ell\pmod p$, with product in $a^{-b_\ell}c$. Indeed this component is contained in $a^\ell J$. Consider the block experiment on $J$ with alphabet $\mathcal A^p$ and target
\[
        G_\ell(h)=F_{s,j}(a^{b_\ell}h u_{v_\ell}),\qquad h\in J.
\]
It is a legal compact-convex interface. Its image on $K=J^\circ$ is $F_{s,j}(c)$, since $K$ is normal. Lemma~\ref{lem:returnperiod55} gives cofinite exact returns in block length for this alphabet.

Restrict the original machine to words that begin with $v_\ell$, end with $b_\ell$ copies of the chosen letter, and have $M=(N-p_\ell-b_\ell)/p$ arbitrary $p$-blocks in between. When $M\ge0$, absorb the fixed prefix into the initialization and the fixed suffix into the terminal decoder; stochastic averaging keeps that decoder in $Y_j$. This produces an error-$\varepsilon$ clocked block machine for $G_\ell$, with registers at original cuts $p_\ell+dm$, $0\le m\le M$. Theorem~\ref{thm:radius54} bounds its narrow cuts $1\le m\le M$ by a constant depending only on $k$ and the fixed data. The omitted cuts of residue $\ell$ number at most $p_\ell+b_\ell+1$: this also charges the block cut zero, while all counted cuts remain in $1,\ldots,N$. If $M<0$, the whole horizon is bounded by $p_\ell+b_\ell$. Sum these bounds over the $p$ residue classes. This counts all cuts $1,\ldots,N$, independently of all other register widths.
\end{proof}
The constants in this theorem depend on the fixed alphabet, return period and certificate, interface, error, width $k$, and chosen finite prefixes. They do not depend on the horizon or on the machine. The original constants $B,d,\chi,L_k$ in Section~\ref{sec:local54} use $d$ for a contraction defect. The return period here is denoted $p$ to keep the two parameters distinct.

\begin{proposition}[Different phases can have different thresholds]\label{prop:parity55}
There is an alphabet with return period two for which $W_{N,0}=1$ on every even horizon, but $W_{N,\varepsilon}\to\infty$ along odd horizons for every $0\le\varepsilon<\rho/2$.
\end{proposition}
\begin{proof}
Take $H=(\R/2\pi\Z)\times\Z/2\Z$ and letters $(\theta,1),(-\theta,1)$, with $\theta/2\pi$ irrational. The generated group is dense in $H$. An exact return must have zero signed count and hence even length; the two-letter return gives every positive even length. Here $J=(\R/2\pi\Z)\times\{0\}$. For one binary-mean query put $F(\phi,0)=0$ and $F(\phi,1)=\rho\cos\phi$, where $0<\rho\le1$. The even-phase target is identically zero and needs one label. The odd-phase radius in the norm $|\cdot|/2$ is $\rho/2$. Apply Theorem~\ref{thm:phases55}.
\end{proof}

\section{Infinite component groups and boundary attainment}\label{sec:profinite55}
This section retains independent finite-dimensional proofs of the strict-radius and exact translate-rank conclusions. Its output spaces are finite-dimensional. Let $H$ again be an arbitrary compact Hausdorff group, and set $K=H^\circ$. The quotient $H/K$ is profinite. We use two standard compact-group consequences of this fact and of Peter--Weyl theory \cite{Folland}: open normal subgroups form a neighborhood basis in a profinite group, and, for a continuous surjection $\pi:H\to G$ onto a compact Lie group, one has $\pi(K)=G^\circ$. To see the latter assertion, $\pi(K)$ is connected and normal, while $G/\pi(K)$ is a continuous group quotient of the profinite group $H/K$ and is therefore totally disconnected. This forces $G^\circ\subset\pi(K)$; the reverse inclusion is immediate. A continuous finite-dimensional representation of a profinite group has finite image, by the closed subgroup theorem and the no-small-subgroups property of a Lie group \cite{Hall}: an open normal subgroup can be placed inside a no-small-subgroups neighborhood, so its image is trivial.

Define the component radius
\begin{equation}\label{eq:R055}
                 R_0=\max_{s,j,g\in H}\rad_{Y_j}(F_{s,j}(gK)).
\end{equation}
Unlike $r(U)$, this quantity need not itself be realized by any finite quotient.

\begin{lemma}[Approximation by finite quotient radii]\label{lem:infradius55}
With $r(U)$ as in \eqref{eq:quotientradius55},
\[
                 \inf_{U\text{ open normal in }H}r(U)=R_0.
\]
\end{lemma}
\begin{proof}
Every open subgroup contains $K$, so $r(U)\ge R_0$. Fix $\eta>0$. Uniform continuity, for the finite family of interfaces, gives an identity neighborhood $V$ such that
\[
 \norm{F_{s,j}(xv)-F_{s,j}(x)}_j<\eta
       \quad(x\in H,\ v\in V,\ s,j).
\]
There is an open normal subgroup $U$ containing $K$ and contained in $KV$, by the profinite quotient basis. Write each $u\in U$ as $kv$ with $k\in K$, $v\in V$. Then $F_{s,j}(gu)$ lies within $\eta$ of $F_{s,j}(gk)$. A center for the latter component image consequently covers the former coset image with an additional error at most $\eta$. Thus $r(U)\le R_0+\eta$.
\end{proof}

\begin{lemma}[A legal finite-dimensional interface approximation]\label{lem:Lieapprox55}
For every $\eta>0$ there is a continuous homomorphism $\pi:H\to G\subset O(D)$ onto a compact Lie group and continuous maps $\widetilde F_{s,j}:G\to Y_j$ such that
\[
           \max_{s,j}\norm{F_{s,j}-\widetilde F_{s,j}\circ\pi}_\infty<\eta.
\]
Moreover, for every $g\in H$,
\[
 \left|\rad_{Y_j}(F_{s,j}(gK))-
 \rad_{Y_j}(\widetilde F_{s,j}(\pi(g)G^\circ))\right|<\eta.
\]
\end{lemma}
\begin{proof}
Approximate all coordinates of all $F_{s,j}$ by finitely many real representative functions. The direct sum of their orthogonal representations is $\pi$, and the approximants are continuous functions on its compact image $G$. Choose Euclidean inner products on the finitely many output spaces. Project each approximant to the closed convex set $Y_j$ by nearest-point projection. This projection is continuous and nonexpansive in the chosen Euclidean norm. Choose constants $a_j,b_j>0$ with $a_j\|v\|_j\le|v|_2\le b_j\|v\|_j$. An initial error less than $a_j\eta/b_j$ in $\|\cdot\|_j$ then gives a projected error less than $\eta$, by Euclidean nonexpansiveness. These constants are fixed before choosing the approximants. The closed subgroup $G\subset O(D)$ is a compact Lie group. Since $\pi(K)=G^\circ$, the two indicated images have Hausdorff distance less than $\eta$. Their constrained radii therefore differ by less than $\eta$.
\end{proof}

\begin{theorem}[Arbitrary compact groups on both strict sides]\label{thm:profinite55}
For any compact-group interface, $\varepsilon>R_0$ implies a finite stationary permutation realization. At $\varepsilon=R_0$, such a realization exists exactly when $r(U)\le R_0$ for some open normal $U$.

Suppose in addition that executable exact returns have all sufficiently large lengths. If $\varepsilon<R_0$, then for each $k$ there is a horizon-independent upper bound on the number of cuts of width at most $k$, with all intervening widths unrestricted. In particular $W_{N,\varepsilon}\to\infty$. No finite-component assumption is made in either assertion.
\end{theorem}
\begin{proof}
The stationary assertions follow from Lemma~\ref{lem:infradius55} and Theorem~\ref{thm:finitequotient55}. For the lower bound choose a component image of radius greater than $\varepsilon$ and then $\eta>0$ with $2\eta$ smaller than that strict difference. Apply Lemma~\ref{lem:Lieapprox55}. A machine correct for the original interface at error $\varepsilon$ is correct for the approximating interface at error at most $\varepsilon+\eta$. The chosen component radius of that interface is greater than $\varepsilon+\eta$. The projected commands still have cofinite exact returns and generate $G$ densely. Theorem~\ref{thm:radius54} therefore gives the desired occupation bound for the very same registers. This argument never assumes that $K$ is open or that executable products lying exactly in $K$ are dense.
\end{proof}
Theorem~\ref{thm:completeboundary56} now identifies this stationary boundary criterion with clocked boundedness, including $R_0>0$, and under the weaker approximate-return hypothesis. The zero-error boundary also admits the following independent translate-rank argument.

\begin{lemma}[Evaluation rank and translate dimension]\label{lem:evaluation56}
For a real function $f$ on a group, the supremum of the ranks of finite matrices $(f(y_lx_i))_{i,l}$ equals the dimension, possibly infinite, of the span of the functions $x\mapsto f(yx)$. These are left translates with the convention $(L_gf)(x)=f(g^{-1}x)$, after replacing $y$ by $g^{-1}$.
\end{lemma}
\begin{proof}
A finite-dimensional span bounds every evaluation rank. Conversely, if $f_1,\ldots,f_m$ are independent functions, their evaluation vectors $(f_1(x),\ldots,f_m(x))$ span $\R^m$: a proper span would have a nonzero annihilator, giving a linear dependence of the functions. Select $m$ evaluation vectors forming a basis. This gives an invertible evaluation matrix and proves the reverse inequality for every finite $m$.
\end{proof}

\begin{theorem}[The exact clocked boundary]\label{thm:zero55}
Assume cofinite executable exact returns. The following are equivalent:
\begin{enumerate}
\item $\sup_NW_{N,0}<\infty$;
\item all $F_{s,j}$ factor through one finite continuous quotient $H/U$, with $U$ open normal;
\item there is an exact stationary finite permutation realization.
\end{enumerate}
This includes groups with infinitely many connected components.
\end{theorem}
\begin{proof}
The implications from the second assertion to the third and from the third to the first are immediate. Suppose the first holds with bound $k$. Theorem~\ref{thm:profinite55} implies $R_0=0$, so every interface is constant on each $K$-coset.

Fix a real coordinate $f$ of one $F_{s,j}$. For any finite families of executable prefix products $x_i$ and suffix products $y_l$, use cofinite returns to pad all selected prefix words to one common length $A$ and all suffix words to one common length $B$. Choose an exact realization at horizon $A+B$ of width at most $k$. At its cut $A$ the matrix
\[
                            (f(y_lx_i))_{i,l}
\]
factors through the $k$ hidden labels: the left factors are the actual state distributions after the padded prefixes and the right factors are the actual conditional terminal expectations on the padded suffixes. Hence this matrix has rank at most $k$. Padding does not need to preserve hidden states; it only preserves the physical products used in the displayed entries. By density and continuity the same rank bound holds for arbitrary $x_i,y_l\in H$.

It follows that the span of the left translates of $f$ in $C(H,\R)$ has dimension at most $k$. Otherwise $k+1$ linearly independent translates would have an invertible evaluation matrix at some $k+1$ points, contrary to the preceding rank bound. Lemma~\ref{lem:evaluation56} formalizes this evaluation assertion. Equivalently it follows inductively: a nonzero function is nonzero somewhere, and a function outside the span of the preceding ones cannot have all its evaluations forced by theirs.

Take the sum $V$ of these translate spaces over all seed--query coordinates. It is finite-dimensional and invariant under left translation. Translation is continuous in the uniform norm and preserves that norm, so it gives a continuous representation of $H$ on $V$ with compact image; Haar averaging makes it orthogonal. Because $K$ is normal and every original coordinate is constant on $K$-cosets, $K$ acts trivially on every translate and thus on $V$. The representation factors through $H/K$. Its image is finite by the profinite finite-dimensional representation fact above. Its kernel $U$ is open normal. Since every original coordinate belongs to $V$, all interfaces are $U$-invariant and descend to $H/U$.
\end{proof}

\begin{proposition}[A nonattained zero-radius boundary]\label{prop:adic55}
On a compact group with infinitely many components, every positive error can have a bounded stationary realization while the zero-error clocked width is unbounded.
\end{proposition}
\begin{proof}
Let $H=\Z_2$ be the additive group of $2$-adic integers, with commands $0$ and $1$. Their positive products are dense, and the zero command supplies exact returns of every length. If $x=\sum_{n\ge0}x_n2^n$, $x_n\in\{0,1\}$, prescribe the binary success probability
\[
                         f(x)=\sum_{n\ge0}\frac{2x_n}{3^{n+1}}.
\]
Uniform convergence makes $f$ continuous. It is injective: at the first differing digit the contribution $2/3^{n+1}$ exceeds the sum $1/3^{n+1}$ of all possible later differences. Since $H$ is totally disconnected, $K=\{0\}$ and $R_0=0$. Truncation after $m$ digits factors through $\Z/2^m\Z$ and has uniform error at most $\sum_{n=m}^\infty2/3^{n+1}=3^{-m}$. Thus every positive tolerance admits a finite quotient machine. But the injective function $f$ cannot factor through any finite quotient. Theorem~\ref{thm:zero55} excludes bounded exact clocked width, and Theorem~\ref{thm:finitequotient55} excludes a finite exact stationary realization.
\end{proof}

\begin{theorem}[Linear exact width at a profinite boundary]\label{thm:adiclinear55}
For the interface of Proposition~\ref{prop:adic55},
\[
                         W_{N,0}=\Theta(N).
\]
More precisely, if $n$ is a power of two with $2n-1\le N$, then $W_{N,0}\ge n$, while $W_{N,0}\le N+1$. Every exact realization, with unrestricted intermediate widths, has at most $4k$ cuts in $\{1,\ldots,N\}$ of width at most $k$. Nevertheless, for each $\varepsilon>0$, $\sup_NW_{N,\varepsilon}<\infty$.
\end{theorem}
\begin{proof}
For nonnegative integers $z$ write $v_2(z+1)$ for the exponent of two in $z+1$. Incrementing the binary expansion, with $v_2(z+1)$ trailing ones, gives
\begin{equation}\label{eq:adicdifference55}
 b(z):=f(z+1)-f(z)=\frac53\,3^{-v_2(z+1)}-1.
\end{equation}
Let $n=2^m$, and form $B_n=(b(i+j))_{0\le i,j<n}$. Define a function on $\Z/n\Z$ by
\[
 \psi_m(x)=3^{-v_2(x)}\ (x\ne0),\qquad \psi_m(0)=3^{-m}.
\]
Here the valuation of a nonzero residue is independent of its representative. Since $1\le i+j+1\le2n-1$, its valuation equals that used in $\psi_m(i+j+1\bmod n)$, including the case $i+j+1=n$. Thus replacing row $i$ by the old row $(-i-1)\bmod n$ changes its entry at column $j$ to the kernel value at $j-i\bmod n$. This is the circulant with kernel $(5/3)\psi_m-1$.

On $\Z/n\Z$ one has the elementary identity
\[
            \psi_m=1-2\sum_{\ell=1}^m3^{-\ell}
                               \mathbf1_{2^\ell\mid x}.
\]
Use the eigenvector with entries $\exp(2\pi \mathrm i qj/n)$; its eigenvalue is the kernel sum against $\exp(2\pi \mathrm i qx/n)$. The sign is immaterial for subgroup indicators, whose transforms below are real. The Fourier eigenvalue at frequency zero is
\[
 \mu_0=\frac23 n6^{-m}>0.
\]
At frequency $q\ne0$ it is
\[
 \mu_q=-\frac{10n}{3}
           \sum_{\substack{1\le\ell\le m\\ n/2^\ell\mid q}}6^{-\ell}<0.
\]
Indeed a subgroup indicator has Fourier sum $n/2^\ell$ exactly when $n/2^\ell$ divides $q$, and zero otherwise. For $q\ne0$ the summand $\ell=m$ is always present. At $m=0$ the sole eigenvalue is $2/3$. Consequently $B_n$ is invertible for every dyadic $n$.

At an arbitrary command cut $t<N$, take a dyadic $n\le\min(t+1,N-t)$. For each $i<n$ choose a prefix of length $t$ with $i$ increment commands. For each $j<n$ choose two suffixes of length $N-t$ with respectively $j+1$ and $j$ increment commands. Fill all remaining positions with zero commands. The difference of their exact outputs is $b(i+j)$. Hence $B_n$ factors through the actual register $S_t$: one factor consists of the prefix state rows and the other of the differences of the two conditional suffix decoder columns. Invertibility gives $|S_t|\ge n$. The difference column need not be a legal decoder; it is used only for this rank bound on two legal experiments.

Taking $t=n-1$ proves $W_{N,0}\ge n$ whenever $2n-1\le N$. The largest dyadic $n\le(N+1)/2$ is greater than $(N+1)/4$, proving the linear lower order. Storing the number of increment commands uses at most $N+1$ labels and gives the exact upper bound.

For occupation, if $t<N$ and $|S_t|\le k$, the largest dyadic integer not exceeding $\min(t+1,N-t)$ is at most $k$, so $\min(t+1,N-t)<2k$. Such cuts lie within fewer than $2k$ positions of one of the two ends. Including the terminal cut gives at most $4k$ positive cuts. No width at another cut was restricted. The positive-error assertion is the finite-quotient truncation in Proposition~\ref{prop:adic55}.
\end{proof}
This supplies a sharp quantitative law on a genuinely infinite profinite group. Its linear exact-width obstruction and bounded positive-error realizations concern one fixed continuous interface; they are not a uniform approximation statement over all continuous outputs on that group.

\section{Binary, categorical, and observable interfaces}\label{sec:observable53}
\begin{corollary}[Component oscillation for binary means]\label{thm:classification53}
For a finite continuous binary-mean interface $f_{s,j}:H\to[-1,1]$, with the return hypothesis of Theorem~\ref{thm:radius54}, the exact bounded-width threshold in total variation is
\begin{equation}\label{eq:critical53}
 \epsc=\frac14\max_{s,j,c}\left(\max_{h\in c}f_{s,j}(h)-\min_{h\in c}f_{s,j}(h)\right).
\end{equation}
It is attained by a stationary permutation realization. Below it every fixed width has the occupation bound \eqref{eq:occupation54}.
\end{corollary}
\begin{proof}
In the norm $|\cdot|/2$, the radius of an interval is one quarter of its length. Its midpoint lies in $[-1,1]$. Apply Theorem~\ref{thm:radius54}.
\end{proof}
This recovers the scalar component theorem, including all of its nonuniform quantifiers. The vector theorem does not follow by applying this corollary separately to coordinates, because the coordinate centers need not belong to the output simplex.

\begin{proposition}[A three-outcome threshold]\label{prop:categorical54}
On the circle let $\omega=e^{2\pi i/3}$ and, for $j=0,1,2$, set
\[
 p_j(\phi)=\frac19\left|1+e^{i\phi}\omega^{-j}+e^{2i\phi}\omega^{-2j}\right|^2.
\]
For any dense circle command alphabet satisfying the return condition, the categorical law $p=(p_0,p_1,p_2)$ has critical total-variation error $2/3$. Each separately requested binary event $\{j\}$ has threshold $1/2$.
\end{proposition}
\begin{proof}
Orthogonality of the three characters gives $\sum_jp_j=1$; each entry is nonnegative. At $\phi=2\pi j/3$ the vector is the simplex vertex $e_j$. Hence every center $q\in\Delta_3$ has maximum distance at least
\[
 \max_j\TV(e_j,q)=1-\min_jq_j\ge2/3.
\]
The uniform center has distance at most $2/3$ from every point of $\Delta_3$, by convexity and the values at its vertices. Thus the constrained radius is exactly $2/3$. On the other hand each $p_j$ has range $[0,1]$, so its binary mean $2p_j-1$ has oscillation two. Apply Theorem~\ref{thm:radius54} and Corollary~\ref{thm:classification53}. In particular the categorical radius exceeds half the diameter, which here equals $1/2$.
\end{proof}

\begin{corollary}[Joint measure-once unitary laws]\label{cor:quantum54}
Let unitary commands have compact closure $H$ and satisfy the return condition. For a fixed unit seed $\psi_s$ and a finite measurement $E_{j,l}\ge0$, $\sum_lE_{j,l}=I$, prescribe the law
\[
 F_{s,j}(h)_l=\ip{h\psi_s}{E_{j,l}h\psi_s}.
\]
The critical error is its componentwise simplex radius in total variation. For closure $U(d)$ and one orthonormal-basis measurement, this radius is $1-1/d$.
\end{corollary}
\begin{proof}
The displayed map is a continuous simplex-valued polynomial in realified matrix entries. A closed unitary matrix group has finitely many components, so Theorem~\ref{thm:radius54} applies. For $U(d)$, every unit vector is reachable from the seed, and its squared coordinate magnitudes range over the whole simplex. The simplex vertices force radius at least $1-1/d$, attained by the uniform distribution.
\end{proof}
A query selects one complete measurement, not simultaneous answers to incompatible measurements. Several compatible finite outputs, with specified correlations, instead form one joint categorical law. Language recognition at a cutpoint is a different requirement.

For linear binary interfaces on a compact matrix group, write
\[
 f_{s,j}(h)=q_j^Thx_s\in[-1,1],\quad
 V_{\rm r}=\spanop\{hx_s:h\in H,s\in\mathcal S\},
\]
\[
 V_{\rm n}=\{x\in V_{\rm r}:q_j^Thx=0\text{ for all }h\in H,j\in\mathcal J\}.
\]
Both spaces are invariant. The observable space $V_{\rm o}=V_{\rm r}\cap V_{\rm n}^\perp$ is the orthogonal complement of $V_{\rm n}$ \emph{inside} $V_{\rm r}$, canonically realizing the observable quotient.

\begin{theorem}[Intrinsic zero-error criterion]\label{thm:observable53}
For these linear means, bounded exact clocked width is equivalent to finiteness of the generated group's image in $O(V_{\rm o})$, and to trivial action of $K$ on $V_{\rm o}$. An exact stationary permutation realization then exists.
\end{theorem}
\begin{proof}
Corollary~\ref{thm:classification53} identifies bounded exact width with constancy of all interface means on every component. Under that condition, for $k\in K$ and $g,h\in H$, normality gives
\[
 q_j^Th(k-I)gx_s=q_j^T(hkh^{-1})hgx_s-q_j^Thgx_s=0.
\]
Thus $(k-I)V_{\rm r}\subset V_{\rm n}$. On the invariant orthogonal complement $V_{\rm o}$ the difference also lies in $V_{\rm o}$, hence vanishes. The observable action factors through the finite component group. Conversely, if $K$ acts trivially on $V_{\rm o}$, the same difference is invisible and all means are componentwise constant. A finite image of the dense generated group has the same finite image closure from $H$, whose connected subgroup $K$ acts trivially. The zero-dimensional quotient is included.
\end{proof}

\begin{corollary}[The planar threshold]\label{cor:circle53}
Let an orthogonal planar alphabet contain $I$ and an irrational rotation, with any additional orthogonal letters. For signed coordinate seeds $\pm\rho e_1,\pm\rho e_2$, coordinate queries, and $0<\rho\le1$, the threshold is $\epsc=\rho/2$. At and above it one command label with a fair binary decoder suffices. Below it every fixed width has a finite occupation budget.
\end{corollary}
\begin{proof}
The identity component is $SO(2)$, and every specified mean has range $[-\rho,\rho]$ on each component of $H=SO(2)$ or $O(2)$. All midpoint decoders equal zero, so the seed/component register collapses to one label.
\end{proof}
For the noncommuting rational alphabet
\[
 I,\quad R=\begin{pmatrix}3/5&-4/5\\4/5&3/5\end{pmatrix},\quad
 F=\begin{pmatrix}1&0\\0&-1\end{pmatrix},
\]
one has $FRF=R^{-1}\ne R$. The rotation has infinite order: otherwise $(3+4i)/5$ and its inverse would be algebraic integers, whereas their rational sum $6/5$ is not an integer. This example has the sharp threshold above. We do not infer a badly approximable angle from rational matrix entries.

\begin{proposition}[Independent toral coordinates]\label{prop:torus53}
Suppose a component is identified with the full product torus $(\R/2\pi\Z)^m$, surjectively in independent coordinates, and a mean on it equals
\[
 a_0+\sum_{l=1}^m(a_l\cos\theta_l+b_l\sin\theta_l).
\]
Its contribution to \eqref{eq:critical53} is $\tfrac12\sum_l\sqrt{a_l^2+b_l^2}$.
\end{proposition}
\begin{proof}
Each summand has the symmetric range with amplitude $\sqrt{a_l^2+b_l^2}$, and all extrema are simultaneously attainable in the full product. Divide the resulting oscillation by four.
\end{proof}
On a proper subtorus, extrema must instead be taken subject to its relations. Ambient coordinates do not certify independence. Similarly, invisible directions do not create an obstruction: for $H=\{\diag(R_\theta,1)\}$, seed and query $e_3$ give constant mean one, whereas seed $\rho e_1$ and query $e_1$ have threshold $\rho/2$.

\begin{proposition}[Compact similarity]\label{prop:similarity53}
The classification applies to a finite alphabet in $GL(D,\R)$ whose generated group, including inverses, is uniformly bounded, with the executable return hypothesis retained. Conjugating to an invariant Euclidean metric changes neither output probabilities nor stochastic widths.
\end{proposition}
\begin{proof}
Uniform bounds on matrices and inverses make the closure a compact group. The matrix $Q=\int_Hh^Th\,d\mu(h)$ is positive definite and satisfies $A^TQA=Q$. Hence $Q^{1/2}AQ^{-1/2}$ is orthogonal. Transport the continuous output maps through this conjugacy; for a linear mean transport $x_s$ to $Q^{1/2}x_s$ and $q_j$ to $Q^{-1/2}q_j$. Exact identity products and their word lengths are preserved.
\end{proof}

\subsection{A curved categorical image with distinct component centers}
\begin{proposition}\label{prop:curved55}
Let $H=(\R/2\pi\Z)\times\Z/2\Z$, with an identity letter, an irrational rotation in the first factor, and a toggle of the second factor. Put
\[
 q_0=(\tfrac12,\tfrac13,\tfrac16),\qquad
 q_1=(\tfrac16,\tfrac13,\tfrac12),\qquad t=\tfrac1{12},
\]
and prescribe the three-outcome law
\[
 F(\theta,c)_l=q_{c,l}+t\cos(\theta-2\pi l/3),\qquad l=0,1,2.
\]
Each component image is a proper closed curve in the interior of the simplex, with unique constrained total-variation center $q_c$ and radius $t$. The union of both component images has radius $1/4$. Thus the sharp clocked threshold is $1/12$, and the stationary minimum at that threshold is exactly two labels.
\end{proposition}
\begin{proof}
The three cosines sum to zero and every coordinate is at least $1/12$, so $F$ is a legal law in the simplex interior. The first two independent trigonometric coordinates identify a nondegenerate ellipse in its affine plane. For a zero-sum vector $z\in\R^3$, one has $\frac12\norm{z}_1=\max_l|z_l|$. Each coordinate of the curve ranges over $[q_{c,l}-t,q_{c,l}+t]$. Hence for every center $y\in\Delta_3$,
\[
 \sup_\theta\TV(F(\theta,c),y)
       =t+\max_l|q_{c,l}-y_l|.
\]
This proves uniqueness of $q_c$ and radius $t$. For the union, the least possible value of $\max_{c,l}|q_{c,l}-y_l|$ is $1/6$, forced by the two endpoints in the first coordinate and attained by $y=(1/3,1/3,1/3)$. Its radius is therefore $t+1/6=1/4$. A one-label stationary machine has a constant decoder and cannot reach error $t$; two component labels attain it. The clocked threshold follows from Theorem~\ref{thm:radius54}.
\end{proof}

\begin{remark}
The total-variation diameter of a probability simplex is one, so half its diameter is $1/2$; the three-vertex example above has radius $2/3$. In the unitary corollary, $U(d)$ is connected. These facts explain respectively why coordinate midpoints do not solve the categorical problem and why the full-unitary component is the whole group.
\end{remark}

\section{Matching width laws throughout the subcritical interval}\label{sec:rates54}
We now distinguish the dimension of the command arithmetic from the dimension of the physical orbit. Let $s\ge1$ and let the alphabet contain the rotations $R_{2\pi\alpha_1},\ldots,R_{2\pi\alpha_s}$ and $I$. For the matching upper bound the alphabet is exactly these letters, optionally with the reflection $F=\diag(1,-1)$. Use signed coordinate seeds of radius $\rho$ and coordinate queries. Extra letters do not affect the lower bounds, but are not silently included in the upper construction.

For a real number $x$, write $\|x\|_{\R/\Z}$ for distance to the nearest integer. The dual Diophantine condition is
\begin{equation}\label{eq:DA54}
 \|n\cdot\alpha\|_{\R/\Z}\ge b\|n\|_1^{-\tau}
 \qquad(0\ne n\in\Z^s),\qquad b>0,\quad\tau\ge s.
\end{equation}
In particular $1,\alpha_1,\ldots,\alpha_s$ are rationally independent. When $\tau=s$ this is the dual badly approximable condition. All implicit constants below may depend on $\alpha,b,s,\rho,\varepsilon$, but not on $N$ or a stochastic realization.

\begin{theorem}[Full-subcritical matching law]\label{thm:matched54}
Suppose \eqref{eq:DA54} holds with $\tau=s$, and let $0<\rho<1$. For every fixed $0\le\varepsilon<\rho/2$,
\begin{equation}\label{eq:matched54}
                 W_{N,\varepsilon}=\Theta\bigl(N^{s/(2s+1)}\bigr).
\end{equation}
The upper bound is achieved by an exact horizon-dependent realization. At $\varepsilon\ge\rho/2$, one command label suffices. If $\tau\ge s$ is arbitrary, the lower bound $W_{N,\varepsilon}=\Omega(N^{s/(2\tau+1)})$ still holds throughout the same subcritical interval.
\end{theorem}
The exponent for one badly approximable rotation is $1/3$, not logarithmic. The statement includes zero error when the signal has strict amplitude slack $\rho<1$. The threshold theorem continues to include $\rho=1$, but the exact polynomial upper bound here is not asserted at that endpoint. The earlier arithmetic argument obtained this scale in a smaller noise range. The proof below supplies the missing full-noise lower estimate and states the exact upper realization in full.

\begin{proposition}[An explicit family in every command dimension]\label{prop:algebraicangles54}
Let $\xi$ be a real algebraic integer of degree $s+1$. Then
$\alpha=(\xi,\xi^2,\ldots,\xi^s)$ satisfies \eqref{eq:DA54} with $\tau=s$ and some effectively computable $b>0$.
\end{proposition}
\begin{proof}
For $0\ne n\in\Z^s$ put $H=\|n\|_1$, and choose $n_0\in\Z$ so that
$z=n_0+\sum_{l=1}^s n_l\xi^l$ has modulus $\|n\cdot\alpha\|_{\R/\Z}$.
The degree assumption gives $z\ne0$. Write $\xi_1=\xi,\ldots,\xi_{s+1}$ for the conjugates and put $M_j=\max_{1\le l\le s}|\xi_j|^l$. Then
$|n_0|\le H M_1+1/2$ and
\[
 \left|n_0+\sum_{l=1}^s n_l\xi_j^l\right|
 \le H(M_1+M_j+1/2)\qquad(j\ge2).
\]
The norm of the nonzero algebraic integer $z$ is a nonzero integer. Thus
\[
 |z|\ge H^{-s}\prod_{j=2}^{s+1}(M_1+M_j+1/2)^{-1}.
\]
Choose a positive rational $b$ below this computable positive constant and $1/2$. This proves the claim. It concerns algebraic \emph{angles}, not algebraic entries of their rotation matrices.
\end{proof}

\subsection{Harmonic localization}
For the harmonic calculation assume that the alphabet contains $I,R_\theta$ and that one seed-query pair has mean $\rho\cos\phi$ after accumulated angle $\phi$, where $0<\rho\le1$. Other seeds, queries, or commands impose additional constraints and do not invalidate the lower bound. Let $\lambda=e^{i\theta}$ have infinite order. Put $\delta=2\varepsilon<\rho$, and choose an integer
\begin{equation}\label{eq:r53}
 r\ge1,\qquad r>\frac{\delta}{\rho-\delta},\qquad m=r+1.
\end{equation}
For example $r=\lfloor\delta/(\rho-\delta)\rfloor+1$ works. Define
\begin{equation}\label{eq:circleconstants53}
 \Delta_r=2\left(\frac{\rho r}{r+1}-\delta\right)>0,\qquad
 A_r=\frac{2(1+\rho+\delta)(2r+1)m}{\Delta_r}>1.
\end{equation}

\begin{lemma}[Explicit harmonic localization]\label{lem:harmonic53}
The nonnegative trigonometric polynomials
\[
 p_\pm(\phi)=\frac{(1\pm\cos\phi)^r}{2^{-r}\binom{2r}{r}}
\]
have Haar mean one and weighted means
\[
 \int p_\pm(\phi)\rho\cos\phi\,\frac{d\phi}{2\pi}
                         =\pm\frac{\rho r}{r+1}.
\]
For the moving feature vector
\[
 Z_m(\phi)=(\cos\phi,\sin\phi,\ldots,\cos m\phi,\sin m\phi),
\]
any terminal realization with wordwise mean error at most $\delta$ obeys
\begin{equation}\label{eq:harmonicresidual53}
 \E\norm{\E[Z_m(\phi)\mid S]}\ge
 \frac{\Delta_r}{2(1+\rho+\delta)(2r+1)}.
\end{equation}
\end{lemma}
\begin{proof}
The identity $1+\cos\phi=2\cos^2(\phi/2)$, or the constant Laurent coefficient, gives
\[
 C_r:=\int(1+\cos\phi)^r\frac{d\phi}{2\pi}
            =2^{-r}\binom{2r}{r},\qquad
 C_{r+1}/C_r=\frac{2r+1}{r+1}.
\]
Hence the weighted cosine mean is $C_{r+1}/C_r-1=r/(r+1)$. Translation by $\pi$ gives the negative case.

We use the unweighted real coefficient vector in the basis $1,\cos\phi,\sin\phi,\ldots$, without orthonormal Fourier rescaling. The absolute sum of the real cosine coefficients of either $p_\pm$, including the constant coefficient, is
\[
 B_r=\frac{4^r}{\binom{2r}{r}}\le2r+1.
\]
For $p_+$ its coefficients are nonnegative and their sum is $p_+(0)$; for $p_-$ they differ only by alternating signs. The inequality follows because the central binomial coefficient is the largest of the $2r+1$ coefficients summing to $4^r$. Multiplication by $\rho\cos\phi$ has coefficient absolute sum at most $\rho B_r$, since $\cos u\cos v=(\cos(u+v)+\cos(u-v))/2$. Thus the Euclidean norms of the moving coefficient vectors of $p_\pm$ and $p_\pm\rho\cos\phi$ are at most $2r+1$ and $\rho(2r+1)$, respectively. For completeness, let $c=\E D$, $Q=\E\norm{\E[Z_m\mid S]}$. For either weight, the conditional coefficient identity gives $|c-\int p_\pm\rho\cos\phi|\le\delta+(1+\rho+\delta)(2r+1)Q$. Subtracting the two inequalities gives \eqref{eq:harmonicresidual53}. Correctness is averaged word by word before using this identity. The starting feature norm is $\sqrt m\le m$; $A_r$ deliberately uses the latter bound to obtain a rational closed form.
\end{proof}

The integer $r$ in \eqref{eq:r53} is the localization order; it is unrelated to the number $s$ of rotation letters. For fixed subcritical error it is finite and independent of $N$ and $k$.

\begin{theorem}[A full-noise arithmetic occupation bound]\label{thm:arithmetic54}
Assume \eqref{eq:DA54} and $0\le\varepsilon<\rho/2$. Fix the localization parameters $r,m,A_r$ in \eqref{eq:r53}--\eqref{eq:circleconstants53}. For an integer $k\ge1$, put
\[
 L=\left\lceil(2k)^{1/s}\right\rceil,\qquad B=s(L-1).
\]
Every error-$\varepsilon$ realization of peak at most $k$ satisfies
\begin{equation}\label{eq:finiteDA54}
 N<B\left(1+b^{-2}(mB)^{2\tau}\log A_r\right).
\end{equation}
With no restriction on other widths, the number of cuts of width at most $k$ is at most
\begin{equation}\label{eq:DAoccupation54}
 \left\lceil B-1+B b^{-2}(mB)^{2\tau}\log A_r\right\rceil.
\end{equation}
\end{theorem}
\begin{proof}
Use $s$ consecutive groups of $L-1$ command positions. In group $i$, choose $J_i$ uniformly from $\{0,\ldots,L-1\}$, apply $J_i$ copies of the $i$th rotation, and fill the other positions with identities. This defines an external law on length-$B$ words with $L^s\ge2k$ equally likely products. On the harmonic feature use
\[
 R_m(t)=\diag(R_t,R_{2t},\ldots,R_{mt}).
\]
For two distinct count vectors $J,J'$ and any frequency $1\le l\le m$, condition~\eqref{eq:DA54} and $\|l(J-J')\|_1\le mB$ give
\[
 |e^{2\pi i l(J-J')\cdot\alpha}-1|
 =2\bigl|\sin\bigl(\pi l(J-J')\cdot\alpha\bigr)\bigr|
 \ge4b(mB)^{-\tau}=:\sigma.
\]
Here $2\sin(\pi x)\ge4x$ for $0\le x\le1/2$. Also $b\le1/2$, $m\ge2$, and $B\ge1$, so $\sigma\le1$.

For a unit vector $u=(u_1,\ldots,u_m)$ in the frequency blocks, the squared distance between the two represented orbit points is
\[
 \sum_{l=1}^m\norm{u_l}^2
       |e^{2\pi i l(J-J')\cdot\alpha}-1|^2\ge\sigma^2.
\]
An open ball of radius $\sigma/2$ about any unit center contains at most one point. At least half the $L^s$ points therefore have distance at least $\sigma/2$ from all of any $k$ centers. Since the inner-product defect between unit vectors equals half their squared distance, the finite word law has defect at least
\[
                  d_k=\sigma^2/16=b^2(mB)^{-2\tau}.
\]
Apply conditional-centroid contraction to successive blocks and Lemma~\ref{lem:harmonic53} to the terminal query. The starting norm is $\sqrt m\le m$. Thus
\[
 \left\lfloor\frac NB\right\rfloor
 \le\frac{\log A_r}{-\log(1-d_k)}
 \le b^{-2}(mB)^{2\tau}\log A_r,
\]
which proves \eqref{eq:finiteDA54}. The greedy endpoint count, with identity gaps and no prefix, gives \eqref{eq:DAoccupation54}. Internal registers remain unrestricted. Finally $B<s(2k)^{1/s}$, so \eqref{eq:finiteDA54} implies the asserted $\Omega(N^{s/(2\tau+1)})$ lower bound.
\end{proof}

\begin{remark}
Testing a coordinate unit vector in \eqref{eq:DA54} gives $b\le1/2$. In the block construction, $L=\lceil(2k)^{1/s}\rceil\ge2$, hence $B=s(L-1)\ge1$. The integer $r$ is the localization order and $s$ is the number of rotations. For the algebraic-angle example, $1,\xi,\ldots,\xi^s$ are linearly independent over $\Q$ by the degree $s+1$ hypothesis; this ensures that the algebraic integer used in its norm bound is nonzero.
\end{remark}

\subsection{An exact polygon realization}
Write $\eta_\alpha(q)=\max_i\|q\alpha_i\|_{\R/\Z}$. The following elementary construction explains why a common denominator, rather than the number of numerical orbit points, controls the upper width.

\begin{proposition}[Exact enclosure upper bound]\label{prop:polygon54}
Let $0<\rho<1$, $a=(1+\rho)/2$, and $\gamma=\log(a/\rho)>0$. Choose an integer $q_0\ge4$ with $\cos(\pi/q_0)\ge a$. If $q\ge q_0$ and
\begin{equation}\label{eq:enclosure54}
                4\pi^2N\,\eta_\alpha(q)/q^2\le\gamma,
\end{equation}
then the indicated rotation alphabet, with or without $F$, has an exact horizon-$N$ realization with $q$ command labels.
\end{proposition}
\begin{proof}
Let $P_q$ be the regular polygon with vertices $v_j=(\cos(2\pi j/q),\sin(2\pi j/q))$, $0\le j<q$. For each $i$ choose a nearest integer $p_i$ to $q\alpha_i$ and put
\[
 t=\pi/q,\qquad h_i=2\pi|\alpha_i-p_i/q|\le t,\qquad
 \Lambda_i=\frac{\cos(t-h_i)}{\cos t}.
\]
Writing $n_j=(\cos((2j+1)t),\sin((2j+1)t))$, its facet description is
\[
               P_q=\{x:\ip{n_j}{x}\le\cos t\ (0\le j<q)\}.
\]
The facet inequalities of the regular polygon give $R_{2\pi\alpha_i}P_q\subset\Lambda_iP_q$: the nearest facet normal to a rotated vertex has angular difference $t-h_i$. Further,
\[
 \log\Lambda_i
 \le\log(1+\tan(t)h_i)\le\tan(t)h_i
 \le4\pi^2\eta_\alpha(q)/q^2,
\]
since $\cos h+\tan(t)\sin h\le1+\tan(t)h$ and $\tan t\le2t$ for $0\le t\le\pi/4$. Let $\Lambda=\max(1,\Lambda_1,\ldots,\Lambda_s)$. The standard reflection $F=\diag(1,-1)$ sends $v_j$ to $v_{-j\ ({\rm mod}\ q)}$ and preserves the chosen vertex set. The identity and reflection therefore preserve $P_q$, so every command $U$ satisfies $\Lambda^{-1}UP_q\subset P_q$.

For each command and each vertex choose a probability row with barycenter $\Lambda^{-1}Uv_j$ in $P_q$. These are actual common command rows, not independent prefix/suffix factorizations. At most three successors suffice by planar convexity. Initialize the seed $\rho u_s$, $u_s\in\{\pm e_1,\pm e_2\}$, with a row whose barycenter is $au_s$. This is possible since $a\le\cos(\pi/q)$, the polygon's inradius. After a length-$N$ word the mean vertex is
\[
                         a\Lambda^{-N}U_wu_s.
\]
By \eqref{eq:enclosure54}, $\rho\Lambda^N\le a$. Hence the terminal columns
\[
                       d_j(v)=\frac{\rho\Lambda^N}{a}\,e_j^Tv
\]
belong to $[-1,1]$ and return exactly $\rho e_j^TU_wu_s$. The rows can be stationary within a horizon; the terminal scaling and chosen denominator may depend on $N$. No one finite all-horizon exact machine is asserted.
\end{proof}

\begin{lemma}[Comparable simultaneous denominators]\label{lem:denominators54}
Assume \eqref{eq:DA54}. For every integer $Q\ge2$ there is $q$ with $1\le q\le Q^s$, $\eta_\alpha(q)\le1/Q$, and
\[
 q\ge c Q^{s/(\tau+1-s)},\qquad
 c=\left(\frac{b}{s^{\tau+1}}\right)^{s/(\tau+1-s)}.
\]
In particular, for $\tau=s$ one has $cQ^s\le q\le Q^s$.
\end{lemma}
\begin{proof}
Partition the torus into $Q^s$ half-open cubes and apply the pigeonhole principle to $0,\alpha,\ldots,Q^s\alpha$. Their difference gives $q\le Q^s$ and integers $p_i$ with $|q\alpha_i-p_i|\le1/Q$.

Put $M=\lfloor q^{1/s}\rfloor$. The $(M+1)^s>q$ vectors in $\{0,\ldots,M\}^s$ cannot all have distinct residues $n\cdot p$ modulo $q$, even when some nearest integers $p_i$ are negative. Subtract two to obtain a nonzero vector $n$ with $q\mid n\cdot p$ and $H:=\|n\|_1\le s q^{1/s}$. Therefore
\[
 bH^{-\tau}\le\|n\cdot\alpha\|_{\R/\Z}
 \le H/(qQ),\qquad
 q^{(\tau+1-s)/s}\ge bQ/s^{\tau+1}.
\]
The displayed bound follows. This argument supplies the needed quantitative transference directly; no unspecified denominator subsequence is being used.
\end{proof}

\begin{proof}[Proof of Theorem~\ref{thm:matched54}]
The lower estimate follows from Theorem~\ref{thm:arithmetic54}. For $\tau=s$, take the constants $a,\gamma,q_0,c$ above and an integer
\[
 Q\ge\max\left(2,(q_0/c)^{1/s},
                \left(\frac{4\pi^2N}{c^2\gamma}\right)^{1/(2s+1)}\right).
\]
Choosing the least such integer, Lemma~\ref{lem:denominators54} gives $q\le Q^s=O(N^{s/(2s+1)})$, $q\ge q_0$, and
\[
 4\pi^2N\eta_\alpha(q)/q^2
 \le4\pi^2N/(c^2Q^{2s+1})\le\gamma.
\]
Proposition~\ref{prop:polygon54} supplies the exact upper bound. At and above $\rho/2$ use the fair decoder from Corollary~\ref{cor:circle53}. This proves all claims.
\end{proof}
The command count $s$ is an arithmetic parameter, not a replacement for the dimension of the physical circle. The constants deteriorate as $\varepsilon\uparrow\rho/2$ through the localization degree, and as $\rho\uparrow1$ through the enclosure slack. Both dependences are explicit in the finite inequalities. The theorem makes no uniform claim at these limiting parameters.

\begin{corollary}[Unit amplitude at positive error]\label{cor:unitnoise55}
For the exact rotation alphabet of Theorem~\ref{thm:matched54}, with or without its specified reflection, suppose the angle vector is dual badly approximable. At $\rho=1$, for every fixed $0<\varepsilon<1/2$,
\[
                     W_{N,\varepsilon}=\Theta(N^{s/(2s+1)}).
\]
\end{corollary}
\begin{proof}
The harmonic localization and orbit-packing lower bound in Theorem~\ref{thm:arithmetic54} permits $\rho=1$ and every fixed $\varepsilon<1/2$. For the upper bound choose $0<\eta<2\varepsilon$ and put $\rho'=1-\eta\in(0,1)$, decreasing $\eta$ if necessary. Apply the exact polygon construction and comparable-denominator lemma to $\rho'$. On every seed, word, and query its mean differs from the unit-amplitude target by at most $\eta$, so the binary total-variation error is at most $\eta/2<\varepsilon$. Its width has the required order. This is a positive-error upper bound, not an exact realization at unit amplitude.
\end{proof}
The case $\rho=1$, $\varepsilon=0$ is not decided by this corollary. In particular, taking $\eta=0$ would destroy the strictly positive enclosure slack required by Proposition~\ref{prop:polygon54}.

\section{Spherical experiments with a spectral gap}\label{sec:nonabelian56}
The next result concerns a full vector orbit, not an abelian factor of a nonabelian system. Write $\sigma$ for normalized area on $S^{d-1}$, where $d\ge2$. Let a finite orthogonal alphabet contain the identity, and let $\mu$ be a probability law on its letters. Assume that its averaging operator
\[
                  Pf(x)=\E_{R\sim\mu}f(Rx)
\]
satisfies
\begin{equation}\label{eq:spectralgap56}
        \|P^n f\|_{L^2(\sigma)}\le\lambda^n\|f\|_{L^2(\sigma)}
        \quad\left(\int f\,d\sigma=0\right),\qquad 0<\lambda<1.
\end{equation}
The physical target is $F(g)=\rho g x_0$ for a fixed $x_0\in S^{d-1}$, with output set the closed Euclidean unit ball and error measured in Euclidean norm.

\begin{theorem}[A spherical width law up to a logarithm]\label{thm:spherical56}
Assume \eqref{eq:spectralgap56}, $0<\rho<1$, and $0\le\varepsilon<\rho$. There are constants $c,C>0$ such that
\begin{equation}\label{eq:sphericalrate56}
 c\left(\frac{N}{\log(N+2)}\right)^{(d-1)/2}
       \le W_{N,\varepsilon}\le C(N+1)^{(d-1)/2}.
\end{equation}
The upper realization is exact. Every error-$\varepsilon$ realization has at most
\[
                      C k^{2/(d-1)}\log(k+2)
\]
positive cuts of width at most $k$, with other widths unrestricted. The constants may depend on $d,\lambda,\rho,\varepsilon$, but not on the horizon or the realization. At unit amplitude the same bounds hold for every fixed $0<\varepsilon<1$; exact unit amplitude is not included in the upper assertion.
\end{theorem}

\begin{lemma}[Spherical quantization and executable contraction]\label{lem:sphericalgap56}
Under \eqref{eq:spectralgap56}, there are $c_d>0$ and integers $B_k\le C_{d,\lambda}\log(k+2)$ such that
\[
 \E_{w\sim\mu^{B_k}}\max_{i\le k}\langle v_i,u_w u\rangle
                  \le1-c_d k^{-2/(d-1)}
\]
for all unit vectors $u,v_1,\ldots,v_k$. The defect may be decreased so that it lies in $(0,1)$. When $d=3$ one may take defect $1/(2k)$ and
\begin{equation}\label{eq:sphere3constants56}
                  B_k=\left\lceil\frac{\log(64k^2)}{-\log\lambda}\right\rceil.
\end{equation}
\end{lemma}
\begin{proof}
Normalized spherical cap measure obeys
\[
       c_d h^{d-1}\le\sigma\{x:\|x-u\|_2<h\}\le C_d h^{d-1}
                     \qquad(0<h\le1),
\]
uniformly in $u$. This follows by writing area in polar angle as a constant times $\sin^{d-2}\theta\,d\theta$ and using $h=2\sin(\theta/2)$; for $d=2$ the same formula is the arc-length formula. For $f(x)=\max_i\langle v_i,x\rangle$, choose $h=c k^{-1/(d-1)}\le1$ so that the union of the $k$ radius-$h$ caps has area at most $1/2$. Outside them, $1-f(x)\ge h^2/2$. Thus
\begin{equation}\label{eq:haarsphere56}
                         \int f\,d\sigma\le1-a_d k^{-2/(d-1)}.
\end{equation}
The function $f$ is $1$-Lipschitz for chordal distance and bounded by one in absolute value.

Let $K_\delta$ average over the chordal cap of radius $\delta\le1$. It commutes with all rotations, fixes constants, and satisfies $\|K_\delta f-f\|_\infty\le\delta$. The cap density at a center has $L^2$ norm at most $C_d\delta^{-(d-1)/2}$. Hence \eqref{eq:spectralgap56} and Cauchy--Schwarz imply, uniformly in the center,
\[
 \left|P^B f(u)-\int f\,d\sigma\right|
 \le\delta+C'_d\delta^{-(d-1)/2}\lambda^B.
\]
Take $\delta$ to be a fixed sufficiently small multiple of $k^{-2/(d-1)}$ and then $B=O_{d,\lambda}(\log(k+2))$ so that the right side is at most half the defect in \eqref{eq:haarsphere56}. The law of the $B$ successive letters is a finite executable word law, proving the general assertion.

For $d=3$, a chordal cap of radius $h\le2$ has normalized area $h^2/4$. The union bound gives $\Pr(1-f<t)\le kt/2$, so integration for $0\le t\le2/k$ gives $\int(1-f)\,d\sigma\ge1/k$. A radius-$\delta$ cap density has $L^2$ norm $2/\delta$, and $\|f-\int f\|_2\le2$. Thus the mixing error is at most $\delta+4\lambda^B/\delta$. Choose $\delta=1/(4k)$ and \eqref{eq:sphere3constants56}; the error is at most $1/(2k)$.
\end{proof}

\begin{proof}[Proof of Theorem~\ref{thm:spherical56}]
Let $Z=u_wx_0$ and let $D(S)$ be the terminal decoder. For every deterministic word, correctness and $\|Z\|_2=1$ give
\[
           \langle\E[D(S)\mid w],Z\rangle\ge\rho-\varepsilon.
\]
Under any external word law this implies
\begin{equation}\label{eq:sphericalresidual56}
       \rho-\varepsilon\le\E\langle D(S),Z\rangle
             \le\E\|\E[Z\mid S]\|_2,
\end{equation}
since the decoder lies in the unit ball. Lemma~\ref{lem:centroid54}, applied with the executable law of Lemma~\ref{lem:sphericalgap56}, contracts the last conditional norm by $1-a_k$ at each selected narrow endpoint, where $a_k=c_d k^{-2/(d-1)}\in(0,1)$. Initially the norm is one. Fixed identity-letter gaps cannot increase it, even though their stochastic rows may process the register. After $J$ selected endpoints, \eqref{eq:sphericalresidual56} gives
\[
            J[-\log(1-a_k)]\le\log\frac1{\rho-\varepsilon}.
\]
Greedy spacing by $B_k$ as in the occupation proof discards at most $B_k-1$ initial cuts and charges at most $B_k$ eligible cuts to each selected endpoint. Therefore the occupation count is at most
\begin{equation}\label{eq:sphericaloccupation56}
 B_k\left(1+a_k^{-1}\log\frac1{\rho-\varepsilon}\right).
\end{equation}
If every register has width at most $k$, this bounds $N$. Inverting the bound gives the lower order in \eqref{eq:sphericalrate56}. Explicitly, when $k\le(N+1)^{d-1}$ its logarithm is $O_d(\log(N+2))$, and otherwise the asserted lower bound is automatic.

For the upper bound, choose a maximal $\delta$-separated set of vertices on $S^{d-1}$ containing $x_0$, with $0<\delta\le1/2$. The disjoint caps of radius $\delta/2$ show that it has at most $C_d\delta^{-(d-1)}$ vertices. Maximality makes it a $\delta$-net. Its convex hull $Q$ satisfies
\begin{equation}\label{eq:sphericalpolytope56}
                r\mathbb B^d\subset Q\subset\mathbb B^d,
                         \qquad r=1-\delta^2/2,
\end{equation}
because for every unit direction $u$ some vertex $v$ has $\langle u,v\rangle\ge1-\delta^2/2$. The inclusion follows from the supporting half-space characterization of a closed convex body. For each command $R$ and vertex $v_i$, the point $rRv_i$ belongs to $Q$. Choose barycentric coordinates of that point as the $i$th row of one common stochastic command matrix. The identity command is treated in the same way, with the same factor $r$.

Initialize at the vertex $x_0$ and decode vertex $v_i$ as $(\rho/r^N)v_i$. Induction on the row means gives output exactly $\rho u_wx_0$ on every length-$N$ word. The decoder is legal if $r^N\ge\rho$. Choose
\[
                 \delta^2=\min\left(\tfrac14,
                                      \frac{-\log\rho}{N+1}\right).
\]
Since $\log(1-\delta^2/2)\ge-\delta^2$, the required inequality holds. The number of vertices is $O_{d,\rho}((N+1)^{(d-1)/2})$. Both rows and decoder may depend on the horizon, as permitted by the clocked model.

For $\rho=1$ and $\varepsilon>0$, the same lower proof applies. An exact upper realization at amplitude $\rho'=1-\varepsilon/2$ differs from the unit-amplitude target by $\varepsilon/2$ in Euclidean norm, proving the stated upper bound. There is no assertion of exactness at unit amplitude.
\end{proof}

\begin{corollary}[A rational noncommuting alphabet on $\mathrm{SO}(3)$]\label{cor:SO356}
Let the alphabet consist of the identity, the matrices
\[
 R_x=\begin{pmatrix}1&0&0\\0&3/5&-4/5\\0&4/5&3/5\end{pmatrix},
 \qquad
 R_z=\begin{pmatrix}3/5&-4/5&0\\4/5&3/5&0\\0&0&1\end{pmatrix},
\]
and their inverses. For $F(g)=\rho g e_3$, with $0<\rho<1$ and $0\le\varepsilon<\rho$, the bounds are
\[
             c\frac{N}{\log(N+2)}\le W_{N,\varepsilon}\le C(N+1).
\]
At error $\varepsilon\ge\rho$ a single zero decoder suffices. Every subcritical realization has at most $Ck\log(k+2)$ positive cuts of width at most $k$.
\end{corollary}
\begin{proof}
The angle $\theta$ of both rotations has $2\cos\theta=6/5$. It is not a rational multiple of $2\pi$: otherwise $2\cos\theta$ would be a rational algebraic integer, hence an integer. The closures of the two cyclic groups are the full circle subgroups about their distinct axes. Their Lie algebras and bracket span $\mathfrak{so}(3)$, so the generated compact closure is $\mathrm{SO}(3)$. The matrices do not commute.

Under the double cover $\mathrm{SU}(2)\to\mathrm{SO}(3)$, their unit-quaternion lifts are $(2+i)/\sqrt5$ and $(2+k)/\sqrt5$. They have algebraic matrix entries and generate a dense subgroup of $\mathrm{SU}(2)$: its closed image is all of $\mathrm{SO}(3)$, so its Lie algebra has dimension three. The spectral-gap theorem of Bourgain--Gamburd \cite{BG2012} applies. Adding the identity with positive mass makes the symmetric averaging operator have norm strictly less than one on mean-zero $L^2$. The spherical representation is a subrepresentation of the regular representation, so the uniform law on the five displayed letters satisfies \eqref{eq:spectralgap56} on $S^2$. Apply Theorem~\ref{thm:spherical56}. The radius of the full sphere $\rho S^2$ inside the unit ball is $\rho$, attained at zero.
\end{proof}
The spectral gap in this corollary is an imported theorem, not a numerical assertion extracted from finite orbit checks. This retained block estimate has a logarithmic loss. Theorem~\ref{thm:sharpoccupation61} supplies the matching linear lower order under the same spectral-gap hypothesis. It does not compute a numerical gap for this particular alphabet; Theorem~\ref{thm:effectivelps61} gives a separate alphabet with explicit constants.

\section{A uniform orbit criterion for quantitative width}\label{sec:uniform57}
The spherical estimate has an extension in which the dimension of the ambient feature space is not the relevant exponent. Two geometric quantities enter: a uniform small-ball exponent for every possible conditional-centroid direction, and the covering dimension of the particular target orbit. They need not agree. Keeping this distinction is essential for matrix-valued experiments.

Let $G$ be a compact connected Lie group with normalized Haar measure $m_G$ and a bi-invariant Riemannian distance. Let $R:G\to O(E)$ be a continuous orthogonal representation on a finite-dimensional real Euclidean space. Fix $z_0\in E$ of norm one, put
\[
 \mathcal O=R(G)z_0,\qquad Q=\operatorname{conv}(\mathcal O),
 \qquad q=\dim\mathcal O,
\]
and assume that $\mathcal O$ spans $E$ and that $E$ has no nonzero fixed vector. Thus $q\ge1$. A finite alphabet in $G$ contains the identity and carries a probability law $\mu$ for which
\begin{equation}\label{eq:groupgap57}
 \|P^nf\|_2\le\lambda^n\|f\|_2\quad\left(\int_Gf\,dm_G=0\right),
 \qquad Pf(g)=\E_{a\sim\mu}f(u_ag),\quad 0<\lambda<1.
\end{equation}
Letters outside the support of $\mu$ are allowed. All command products and stochastic rows retain the conventions of Definition~\ref{def:machine54}.

The uniform geometric hypothesis is that, for some $p>0$, $A<\infty$ and $h_0>0$,
\begin{equation}\label{eq:smallball57}
 \sup_{u,v\in S(E)}m_G\{g:\|R(g)u-v\|<h\}\le Ah^p
                      \qquad(0<h\le h_0).
\end{equation}
It is imposed on all unit vectors $u$, not only on $\mathcal O$. A conditional mean of points in $\mathcal O$, after normalization, need not belong to $\mathcal O$.

\begin{lemma}[From orbit small balls to executable contraction]\label{lem:orbitgap57}
Under \eqref{eq:groupgap57}--\eqref{eq:smallball57}, constants $c,C>0$ and integers $B_k\le C\log(k+2)$ can be chosen so that
\begin{equation}\label{eq:orbitgap57}
 \E_{w\sim\mu^{B_k}}\max_{1\le i\le k}
       \langle v_i,R(u_w)u\rangle\le1-c k^{-2/p}
\end{equation}
for every $u,v_1,\ldots,v_k\in S(E)$. The constant $c$ can be decreased so that the defect lies in $(0,1)$ for every $k\ge1$.
\end{lemma}
\begin{proof}
Choose $h=c_0k^{-1/p}\le\min(h_0,1)$ with $A c_0^p\le1/2$. Outside the union of the $k$ open balls in \eqref{eq:smallball57},
\[
 1-\max_i\langle v_i,R(g)u\rangle\ge h^2/2.
\]
The union has measure at most $1/2$. Consequently the Haar mean of
$f(g)=\max_i\langle v_i,R(g)u\rangle$ is at most $1-a_k$, where
$a_k=c_0^2k^{-2/p}/4$. Repeated centers are allowed. All such functions are bounded by one in absolute value and have a common Lipschitz constant $L$: differentiating the finite-dimensional representation bounds the derivative uniformly on unit vectors.

Write $D=\dim G$, distinct from the cap exponent $p$ and orbit dimension $q$. The ball calculation uses the fixed bi-invariant Riemannian metric on $G$. In exponential coordinates at the identity, Haar measure has a smooth positive density. Hence a metric ball of radius $\delta\le\delta_0$ has measure at least $c_G\delta^D$. Average $P^Bf$ over this ball. Left translations are isometries, so $P^Bf$ has Lipschitz constant at most $L$. Cauchy--Schwarz and \eqref{eq:groupgap57} give
\[
 |P^Bf(1)-\bar f|
 \le L\delta+2c_G^{-1/2}\delta^{-D/2}\lambda^B.
\]
Take $\delta$ to be a sufficiently small fixed multiple of $a_k$, also bounded by $\delta_0$. Then choose $B_k\ge1$ to make the second term at most $a_k/4$. The first term is at most $a_k/4$ as well. Since $\log(a_k^{-1})=O(\log(k+2))$, this requires $B_k\le C\log(k+2)$. Products under $\mu^{B_k}$ are genuine words of one common length. Their reverse chronological order has the same law because the letters are independent with the same law. Thus $P^{B_k}f(1)$ is the expectation in \eqref{eq:orbitgap57}. This proves the lemma with defect $a_k/2$.
\end{proof}

\begin{lemma}[Quadratic inner approximation of an orbit hull]\label{lem:orbithull57}
There are constants $C_1,C_2,\delta_0>0$ such that, for $0<\delta\le\delta_0$, points $z_1,\ldots,z_K\in\mathcal O$, including $z_0$, can be chosen with
\begin{equation}\label{eq:orbithull57}
 K\le C_1\delta^{-q},\qquad
 (1-C_2\delta^2)Q\subseteq\operatorname{conv}\{z_1,\ldots,z_K\}\subseteq Q,
 \qquad 0<1-C_2\delta^2<1.
\end{equation}
\end{lemma}
\begin{proof}
The orbit is a compact smooth embedded homogeneous manifold. Finitely many coordinate charts with bounded derivatives and inverse derivatives give an intrinsic $\delta$-net of size at most $C\delta^{-q}$; adjoining $z_0$ changes only $C$. Use the induced Riemannian metric on the orbit. If the functional $x\mapsto\langle b,x\rangle$ attains its maximum at $x_b\in\mathcal O$, its derivative along the orbit vanishes there. The second derivative along unit-speed orbit geodesics has absolute value at most $C\|b\|$, uniformly in their starting points and directions, by compactness of the unit tangent bundle and smoothness of the embedding. A net point within intrinsic distance $\delta$ of $x_b$ therefore satisfies
\[
 \langle b,z_i\rangle\ge h_Q(b)-C\|b\|\delta^2,
 \qquad h_Q(b)=\max_{x\in Q}\langle b,x\rangle.
\]
Haar averaging gives $\int R(g)z_0\,dm_G=0$. Moreover $0$ lies in the interior of $Q$ in $E$. Otherwise a nonzero supporting functional would be nonnegative on the orbit and have Haar mean zero; continuity and full Haar support would make it vanish on the orbit, contrary to its spanning $E$. Thus $r\mathbb B_E\subset Q$ for some $r>0$, and $h_Q(b)\ge r\|b\|$. The support inequality becomes
$\max_i\langle b,z_i\rangle\ge(1-C\delta^2/r)h_Q(b)$.
The support-function criterion for inclusion of compact convex sets proves \eqref{eq:orbithull57}. Decrease $\delta_0$ to obtain the last bounds.
\end{proof}

\begin{theorem}[Geometric width and occupation]\label{thm:uniformorbit57}
Assume \eqref{eq:groupgap57}--\eqref{eq:smallball57}. Take one seed and one query with target $F(g)=\rho R(g)z_0$, output set $Q$, and Euclidean error. For $0<\rho\le1$ and $0\le\varepsilon<\rho$, put $\zeta=\rho-\varepsilon\in(0,1]$. Every error-$\varepsilon$ realization satisfies
\begin{equation}\label{eq:orbitoccupation57}
 \#\{1\le t\le N:|S_t|\le k\}
 \le C\log(k+2)\left(1+k^{2/p}\log\frac1\zeta\right),
\end{equation}
where $C$ depends on the geometric and spectral-gap data but not on $\rho,\varepsilon,N,k$ or the realization. In particular, at fixed parameters,
\begin{equation}\label{eq:orbitrates57}
 W_{N,\varepsilon}\ge c\left(\frac{N}{\log(N+2)}\right)^{p/2}.
\end{equation}
For $0<\rho<1$ there is an exact common-row realization with
\begin{equation}\label{eq:orbitupper57}
 W_{N,0}\le C\left(1+\frac{N+1}{-\log\rho}\right)^{q/2}.
\end{equation}
For $\rho=1$ and fixed $0<\varepsilon<1$, the upper order is $O_\varepsilon((N+1)^{q/2})$. When $p=q$, the lower and upper exponents agree, with the displayed logarithmic loss in the lower bound. The constrained error threshold is exactly $\rho$, attained by the constant zero decoder.
\end{theorem}
\begin{proof}
Put $Z=R(u_w)z_0$ and let $D(S)\in Q\subseteq\mathbb B_E$ be the decoder. For each deterministic word, the error bound gives
$\langle\E[D(S)\mid w],Z\rangle\ge\rho-\varepsilon$.
Under any external word law, conditioning on the register yields
\[
 \zeta\le\E\langle D(S),Z\rangle
          \le\E\|\E[Z\mid S]\|.
\]
The conditional-centroid Lemma~\ref{lem:centroid54} applies to the law of Lemma~\ref{lem:orbitgap57}. The uniformity over every direction in \eqref{eq:smallball57} is precisely what justifies this application. It contracts the conditional norm by $1-a$ at a selected endpoint of width at most $k$, where $a=c k^{-2/p}\in(0,1)$. The initial norm is one. Fixed identity-letter fillers cannot increase it, regardless of their rows or intervening widths. After $J$ selected endpoints,
\[
 J[-\log(1-a)]\le\log(1/\zeta).
\]
Greedily select narrow cuts at separation at least $B_k$, omitting the first $B_k-1$ cuts. Each selected endpoint accounts for at most $B_k$ cuts. The count is at most $B_k(1+a^{-1}\log(1/\zeta))$, proving \eqref{eq:orbitoccupation57}, also when $\zeta=1$. In that case even one selected endpoint is impossible. More precisely one may take the integer bound
\[
 L_k=B_k\left(1+\left\lceil c^{-1}k^{2/p}\log(1/\zeta)\right\rceil\right).
\]
For a full width bound $k$, either $k\ge N^{p/2}$, or
$\log(k+2)\le C_p\log(N+2)$. In the latter case the occupation inequality
gives $N\le C_\zeta\log(N+2)k^{2/p}$. These alternatives prove
\eqref{eq:orbitrates57}.

For the upper bound take the net in Lemma~\ref{lem:orbithull57}, and write $r=1-C_2\delta^2$. The orbit hull is invariant under every command. Hence $rR(u_a)z_i$ is in the convex hull of the net points. For each of the finitely many command--vertex pairs, choose barycentric coordinates independently; no measurable selection is involved. Use them as row $i$ of a stochastic matrix $T_a$, the same at every epoch within this horizon. Initialize at $z_0$. The conditional row means then evolve as $r^tR(u_w)z_0$. Decode label $i$ by $(\rho/r^N)z_i$, which belongs to $Q$ if $r^N\ge\rho$, since $0\in Q$. Choosing
\[
 \delta^2=\min\left(\delta_0^2,\frac1{2C_2},
                          \frac{-\log\rho}{2C_2(N+1)}\right)
\]
gives this inequality by $\log(1-x)\ge-2x$ for $0\le x\le1/2$. The net size gives \eqref{eq:orbitupper57}. At unit amplitude and positive error use the exact construction for amplitude $1-\varepsilon/2$; its discrepancy from the desired target is $\varepsilon/2$.

Finally, Haar mean zero and $\|R(g)z_0\|=1$ imply
\[
 \int_G\|\rho R(g)z_0-y\|^2\,dm_G=\rho^2+\|y\|^2.
\]
Every constrained center therefore has worst error at least $\rho$, while $y=0$ attains it. Below it \eqref{eq:orbitrates57} rules out bounded width. None of the constructions requires transition matrices to satisfy the physical group relations, and extra finitely many commands do not change the upper construction.
\end{proof}
The geometric ingredients in Lemma~\ref{lem:orbithull57} are the usual local geometry of compact homogeneous spaces; related covering estimates are developed in \cite{Szarek57}. The additional realization statement is the common stochastic-row construction and its combination with a contraction uniform over conditional-centroid directions. A small-ball estimate restricted to the target orbit alone does not imply Theorem~\ref{thm:uniformorbit57}.

\begin{remark}[Dependence of constants]
The cap and executable-gap constants depend only on the fixed group with
its bi-invariant metric, the representation, the uniform cap data
$(A,p,h_0)$, and the declared $L^2$ gap $\lambda$. The hull constants also
depend on the target orbit, its second derivatives, and the inradius of
$Q$ in its span $E$. The occupation bound further depends on
$\zeta=\rho-\varepsilon$. None depends on the horizon, hidden transition
rows, or masses of individual labels. The orbit is the smooth embedded
quotient $G/G_{z_0}$ and $q$ is its real dimension. The logarithmic factor
is the loss of this mixing-block argument.
Theorem~\ref{thm:sharpoccupation61} removes it by a different
entropy--transport argument; all statements and proofs of the block
estimate remain valid.
\end{remark}

\section{Spectral entropy and sharp occupation bounds}\label{sec:entropy61}
The pointwise mixing block in Lemma~\ref{lem:orbitgap57} costs a
logarithm at each selected cut. We now replace that block argument by an
entropy estimate along the selected cuts themselves. The auxiliary
entropy is attached to a smoothed law of \emph{centroid directions}; it
is not the Shannon entropy of the hidden register. The normalization by
centroid mass will be essential.

Let $E$ have real dimension $m\ge2$, let $\Sigma=\{u\in E:\norm u=1\}$,
and let $\sigma$ be normalized round measure on $\Sigma$. Write
$d_\Sigma$ for its geodesic distance and $W_2$ for the corresponding
quadratic Wasserstein distance. All logarithms in this section are
natural. Let $G\subseteq O(E)$ be compact, with normalized Haar measure
$m_G$, and let $\lambda$ be a probability on the command letters. On
$L^2(\Sigma,\sigma)$ define
\[
 Pf(x)=\sum_a\lambda_a f(R(u_a)^{-1}x),\qquad
 \Pi f(x)=\int_G f(R(g)^{-1}x)\,dm_G(g).
\]
Here $R$ denotes the given action; writing $G\subseteq O(E)$ identifies
its image when convenient. The projection $\Pi$ need not have
one-dimensional range. Our spectral hypothesis is
\begin{equation}\label{eq:actiongap61}
 \norm{Pf-\Pi f}_2\le\kappa\norm{f-\Pi f}_2,
                         \qquad 0\le\kappa<1.
\end{equation}
This is a bound on the whole function space, not merely on the
finite-dimensional representation $E$. The group bound
\eqref{eq:groupgap57} implies \eqref{eq:actiongap61}, with the same
$\kappa$: decompose the unitary representation on $L^2(\Sigma)$ into
irreducibles; each nontrivial irreducible occurs in the regular
representation, where the assumed bound holds (use the adjoint
operator if needed for the inverse convention). The trivial summands
are exactly the range of $\Pi$ \cite{Folland}.

\begin{remark}[The full spectral hypothesis]\label{rem:fullgap61}
A contraction on $E$ alone does not imply \eqref{eq:actiongap61}.
For finitely many circle rotations, simultaneous rational approximation
provides arbitrarily high characters on which the averaging eigenvalue
is arbitrarily close to one, although its first harmonic may contract.
Such an alphabet does not satisfy the full spectral hypothesis. This
is why the present theorem does not replace the different Diophantine
width laws by a spurious spherical exponent.
\end{remark}

\subsection{A compactly supported entropy potential}
For $0<h\le1$ set
\begin{equation}\label{eq:kernel61}
 K_h(u,x)=Z_h^{-1}\left(1-\frac{\norm{x-u}^2}{h^2}\right)_+^2,
 \qquad Z_h=\int_\Sigma\left(1-\frac{\norm{x-u}^2}{h^2}\right)_+^2d\sigma(x).
\end{equation}
The normalizing constant is independent of $u$. If $\nu$ is a
probability on $\Sigma$, write $K_h\nu(x)=\int K_h(u,x)d\nu(u)$.
For a density $f$ put
\[
 \mathcal H(f)=\int f\log f\,d\sigma,\qquad
 \mathcal I(f)=\int_{\{f>0\}}\frac{\norm{\nabla f}^2}{f}\,d\sigma,
 \qquad 0\log0=0.
\]

\begin{lemma}[Entropy continuity for spherical mixtures]\label{lem:smoothentropy61}
There are constants $D_m\ge1$ and $B_m<\infty$, depending only on $m$,
such that every finitely supported probability $\nu$ satisfies
\begin{equation}\label{eq:kernelbounds61}
 0\le\mathcal H(K_h\nu)\le L_h:=\log(D_mh^{-(m-1)}),\qquad
 \mathcal I(K_h\nu)\le B_mh^{-2}.
\end{equation}
For any two finitely supported probabilities $\nu,\eta$,
\begin{equation}\label{eq:enttransport61}
 |\mathcal H(K_h\nu)-\mathcal H(K_h\eta)|
                 \le \frac{\sqrt{B_m}}h W_2(\nu,\eta).
\end{equation}
For $m=3$ one can take $D_3=12$ and $B_3=24$.
\end{lemma}
\begin{proof}
Put $s=\norm{x-u}^2$. Polar coordinates give a density proportional to
$s^{(m-3)/2}(1-s/4)^{(m-3)/2}$ on $0<s<4$. On $0<s\le1$ its second
factor is bounded above and below by positive dimension-dependent
constants. Substitution $s=h^2v$ therefore gives
$Z_h\asymp_m h^{m-1}$. Also $\norm{\nabla s}^2=4s-s^2$, so
\[
 \mathcal I(K_h(u,\cdot))=
 \frac4{Z_hh^4}\int_{s<h^2}(4s-s^2)\,d\sigma
                            \le B_mh^{-2}.
\]
Increasing $D_m$ gives $K_h\nu\le D_mh^{-(m-1)}$.
Jensen's inequality gives the lower entropy bound, and the pointwise
upper bound gives the upper one. Cauchy--Schwarz in a finite mixture
shows that its Fisher information is at most the mixture of the
individual Fisher informations.

Choose an optimal finite coupling $(\pi_{ij})$ of $\nu$ and $\eta$.
Join its paired centers by shortest constant-speed spherical geodesics.
Each path is $u_{ij}(t)=\exp(tA_{ij})u_i$, where $A_{ij}$ is
skew-symmetric, acts on a two-dimensional plane, and has operator norm
$d_\Sigma(u_i,v_j)$. An antipodal pair is treated by choosing any such
plane. Define
\[
 f_t=\sum_{i,j}\pi_{ij}K_h(u_{ij}(t),\cdot),\qquad
 J_t(x)=\sum_{i,j}\pi_{ij}K_h(u_{ij}(t),x)A_{ij}x.
\]
Rotation invariance gives $\partial_t f_t+\operatorname{div}J_t=0$.
Pointwise Cauchy--Schwarz and integration yield
\[
 \int_{\{f_t>0\}}\frac{\norm{J_t}^2}{f_t}\,d\sigma
 \le\sum_{i,j}\pi_{ij}d_\Sigma(u_i,v_j)^2=W_2(\nu,\eta)^2.
\]
The preceding Fisher bound holds at every $t$. To justify entropy
differentiation even where $f_t$ vanishes, replace $f_t$ by
$(f_t+\delta)/(1+\delta)$ and $J_t$ by $J_t/(1+\delta)$.
The kernel is $C^1$, the sphere is compact, and integration by parts
followed by Cauchy--Schwarz gives
\[
 \left|\frac d{dt}\mathcal H\left(\frac{f_t+\delta}{1+\delta}\right)\right|
 \le\frac1{1+\delta}\sqrt{\mathcal I(f_t)}
       \left(\int\frac{\norm{J_t}^2}{f_t}\,d\sigma\right)^{1/2}
 \le\frac{\sqrt{B_m}}h W_2(\nu,\eta).
\]
Here quotients at common zeros are zero; the displayed regularized
identity can equivalently be obtained by smoothing the $C^1$ kernel.
Integrate in $t$ and let $\delta\downarrow0$. Uniform boundedness and
continuity of $x\log x$ at zero justify the limit and prove
\eqref{eq:enttransport61}.

On $\Sigma=S^2$, the variable $s$ has constant density $1/4$.
Direct integration gives
\[
 Z_h=h^2/12,\qquad
 \mathcal I(K_h(u,\cdot))=24/h^2-4.
\]
These imply the asserted explicit constants.
\end{proof}
The relation between entropy continuity and quadratic transport is
classical; compare Polyanskiy--Wu \cite[Proposition~1]{PW61} for regular
Euclidean densities. The proof above is included because its densities
have compact support and zeros, and because the ambient sphere and the
finite-mixture path, rather than Euclidean regularity, are used here.

\begin{lemma}[Spectral entropy loss on a thin support]\label{lem:entropydefect61}
Suppose \eqref{eq:actiongap61} holds. Let $f$ be a bounded probability
density supported in a measurable set $A\subseteq\Sigma$ for which
\[
                         \Pi\mathbf1_A\le\theta<1.
\]
Then
\begin{equation}\label{eq:entropydefect61}
 \mathcal H(f)-\mathcal H(Pf)\ge(1-\kappa^2)(1-\theta).
\end{equation}
\end{lemma}
\begin{proof}
Orbital Cauchy--Schwarz gives
$(\Pi\sqrt f)^2\le(\Pi\mathbf1_A)(\Pi f)$ and hence
$\norm{\Pi\sqrt f}_2^2\le\theta$. The invariant and orthogonal
components are preserved by $P$. Since $\norm{\sqrt f}_2=1$,
\[
 \norm{P\sqrt f}_2^2
 \le\norm{\Pi\sqrt f}_2^2+\kappa^2\norm{(I-\Pi)\sqrt f}_2^2
 \le1-(1-\kappa^2)(1-\theta).
\]
For nonnegative numbers $z_a$ with $z=\sum_a\lambda_a z_a$, the
relative-entropy/Hellinger inequality gives
\[
 \sum_a\lambda_a z_a\log z_a-z\log z
 \ge\sum_a\lambda_a(\sqrt{z_a}-\sqrt z)^2.
\]
For completeness this follows by summing
$t\log(t/s)-t+s\ge(\sqrt t-\sqrt s)^2$; continuity handles zeros.
Apply it pointwise to $z_a=f(R(u_a)^{-1}x)$ and integrate. Invariance
of $\sigma$ and Cauchy--Schwarz give
\[
 \mathcal H(f)-\mathcal H(Pf)
 \ge2\left(1-\int\sqrt{Pf}\,P\sqrt f\,d\sigma\right)
 \ge2(1-\norm{P\sqrt f}_2).
\]
The inequality $2(1-\sqrt{1-d})\ge d$ for $0\le d\le1$ completes the
proof. No logarithmic Sobolev inequality or pointwise mixing estimate
has been assumed.
\end{proof}

\subsection{Transport paid for by loss of centroid mass}
For a unit physical random vector $X$ and a finite register $S$, put
\begin{equation}\label{eq:directionlaw61}
 c_s=\E[X\mid S=s],\quad
 \alpha=\sum_s\Pr(S=s)\norm{c_s},\quad
 \nu=\frac1\alpha\sum_{c_s\ne0}\Pr(S=s)\norm{c_s}\,
                       \delta_{c_s/\norm{c_s}}.
\end{equation}
Only positive-mass labels contribute. The directional law is defined
when $\alpha>0$; its weights are not, in general, the probabilities of
the hidden labels.

\begin{lemma}[Centroid-loss transport]\label{lem:centroidtransport61}
Apply an independent letter $a\sim\lambda$, followed by any fixed
physical identity-product word, using the actual stochastic rows of
the machine. Let $\alpha',\nu'$ be the quantities at the endpoint and
suppose $\alpha'\ge\zeta>0$. If $\lambda*\nu$ denotes the mixture of
the rotated laws $R(u_a)_*\nu$, then $\delta=\alpha-\alpha'\ge0$ and
\begin{equation}\label{eq:centroidtransport61}
 W_2(\lambda*\nu,\nu')^2\le\frac{3\pi^2}{\zeta}\delta.
\end{equation}
There is no bound on any intervening register.
\end{lemma}
\begin{proof}
Write $T_a(s,i)$ for the actual composite row, including the fixed
filler. Conditional independence of the row from the unrecorded past
is precisely the machine definition. Thus
\[
 p'_i c'_i=\sum_{s,a}p_s\lambda_aT_a(s,i)R(u_a)c_s.
\]
Set $w_{sai}=p_s\lambda_aT_a(s,i)\norm{c_s}$,
$q_i=\sum_{s,a}w_{sai}$, and $m_i=p'_i\norm{c'_i}$.
For $m_i>0$ put $v_i=c'_i/\norm{c'_i}$; otherwise choose an arbitrary
unit vector. With $u_{sa}=R(u_a)c_s/\norm{c_s}$ for nonzero terms,
\begin{equation}\label{eq:flowidentity61}
 \sum_{s,a}w_{sai}\norm{u_{sa}-v_i}^2=2(q_i-m_i).
\end{equation}
In particular $m_i\le q_i$, $\sum_iq_i=\alpha$,
$\sum_im_i=\alpha'$, and $\delta\ge0$.

Normalize the flows by $\alpha$. They couple $\lambda*\nu$ to
$\widehat\nu=\sum_i(q_i/\alpha)\delta_{v_i}$. Chordal and geodesic
distance satisfy $d_\Sigma(u,v)\le(\pi/2)\norm{u-v}$, so this
coupling has geodesic quadratic cost at most $\pi^2\delta/(2\alpha)$.
If $\delta>0$, the excess masses give a probability $\chi$ with
\[
 \widehat\nu=(\alpha'/\alpha)\nu'+(\delta/\alpha)\chi.
\]
Keeping the common part and transporting the remaining part across the
sphere costs at most $\pi^2\delta/\alpha$. The triangle inequality
for $W_2$ and $(a+b)^2\le2a^2+2b^2$ give
$W_2(\lambda*\nu,\nu')^2\le3\pi^2\delta/\alpha$.
Since $\alpha\ge\alpha'\ge\zeta$, this proves the result.
The case $\delta=0$ follows from the same couplings without an excess
measure. Zero centroids and zero-probability labels require no division.
\end{proof}

\subsection{The sharp lower law}
Assume, for some $A\ge1$, $p>0$, and $h_0>0$, the all-direction cap
bound
\begin{equation}\label{eq:allcaps61}
 \sup_{u,v\in\Sigma}m_G\{g:\norm{R(g)u-v}<h\}
                          \le Ah^p\quad(0<h\le h_0).
\end{equation}
This positive-power cap bound excludes every fixed unit vector. Hence
$E^G=\{0\}$ and the Haar mean of each orbit is zero; in particular
$0\in\operatorname{conv}R(G)z_0$. The constants are fixed with the
representation. For $k\ge1$ define
\begin{equation}\label{eq:parameters61}
 c_0=\min\{1,h_0,(2A)^{-1/p}\},\quad h_k=c_0k^{-1/p},\quad
 \gamma=(1-\kappa^2)/2,\quad L_k^{\rm ent}=\log(D_mh_k^{-(m-1)}).
\end{equation}

\begin{theorem}[Occupation without logarithmic loss]\label{thm:sharpoccupation61}
Let the command alphabet contain a physical identity letter, and assume
\eqref{eq:actiongap61} and \eqref{eq:allcaps61}. Fix a unit seed $z_0$,
let $Q=\operatorname{conv}R(G)z_0$, and suppose the target is
$\rho R(u_w)z_0$, with $0<\rho\le1$ and legal decoder values in $Q$.
Error is Euclidean error in $E$. For $0\le\varepsilon<\rho$ write
$\zeta=\rho-\varepsilon$. Every horizon-$N$ realization, including
one redesigned at that horizon, satisfies
\begin{equation}\label{eq:sharpoccupation61}
 \#\{1\le t\le N:|S_t|\le k\}\le
 \left\lceil 1+\frac{2L_k^{\rm ent}}\gamma+
       \frac{3\pi^2B_m}{\gamma^2h_k^2}\frac{1-\zeta}{\zeta}\right\rceil.
\end{equation}
In particular, for fixed subcritical parameters there are constants
$C_\zeta,c_\zeta>0$ such that
\begin{equation}\label{eq:sharpwidth61}
 \#\{1\le t\le N:|S_t|\le k\}\le C_\zeta k^{2/p},
                 \qquad W_{N,\varepsilon}\ge c_\zeta N^{p/2}.
\end{equation}
The intervening widths are unrestricted. The constants depend only on
$m,A,p,h_0,\kappa,\zeta$ and the fixed round normalization.
\end{theorem}
\begin{proof}
If there are $J\ge2$ positive cuts of width at most $k$, list them as
$t_1<\cdots<t_J$. Use a fixed prefix up to $t_1$. In each subsequent
gap put one independent letter with law $\lambda$, then identity
letters until the next selected cut. Finish by a fixed suffix. This is
an external distribution on actual words of length $N$, chosen from the
advertised width profile, not from the machine's private state. It does
not require the machine's rows on an identity letter to be identities.

Let $X$ be the current unit physical vector. Write $\alpha_j,\nu_j$
for \eqref{eq:directionlaw61} at $t_j$. Wordwise correctness and
$\norm{D_i}\le1$ imply at the terminal cut
\[
 \zeta\le\E\langle X_N,D(S_N)\rangle
       =\sum_i\Pr(S_N=i)\langle\E[X_N\mid S_N=i],D_i\rangle
       \le\alpha_N.
\]
The first inequality averages the conditional \emph{mean} error on
each deterministic word. It makes no assertion about a single sampled
label's error. Conditional Jensen, or the flow calculation above,
shows that centroid mass cannot increase along any gap or fixed
suffix. Thus $1\ge\alpha_j\ge\zeta$, and
\[
 \delta_j:=\alpha_j-\alpha_{j+1}\ge0,\qquad
                         \sum_{j<J}\delta_j\le1-\zeta.
\]

Fix $h=h_k$ and put $f_j=K_h\nu_j$. Each $f_j$ is supported on at most
$k$ open chordal caps of radius $h$. The cap assumption gives
$\Pi\mathbf1_{\{f_j>0\}}\le Akh^p\le1/2$.
Lemma~\ref{lem:entropydefect61} therefore gives
$\mathcal H(Pf_j)\le\mathcal H(f_j)-\gamma$.
Equivariance of the kernel says $Pf_j=K_h(\lambda*\nu_j)$.
Lemmas~\ref{lem:smoothentropy61} and~\ref{lem:centroidtransport61} imply
\[
 \mathcal H(f_{j+1})\le\mathcal H(f_j)-\gamma
              +\frac{\pi\sqrt{3B_m}}{h\sqrt\zeta}\sqrt{\delta_j}.
\]
Writing $M=J-1$, summing, using $0\le\mathcal H(f_j)\le L_h$, and
applying Cauchy--Schwarz only over these $M$ gaps, gives
\begin{equation}\label{eq:entropybudget61}
 \gamma M\le L_h+
       \frac{\pi\sqrt{3B_m}}h\sqrt{M\frac{1-\zeta}{\zeta}}.
\end{equation}
Young's inequality bounds the second term by $\gamma M/2+
3\pi^2B_m(1-\zeta)/(2\gamma h^2\zeta)$. This proves
\eqref{eq:sharpoccupation61}; $J\le1$ is immediate. The logarithm
in $L_h=O(1+\log k)$ is additive, not multiplied by $k^{2/p}$.
Since $1+\log k=O_p(k^{2/p})$, the first bound in
\eqref{eq:sharpwidth61} follows. If the whole profile has width at
most $k$, its $N$ positive cuts are all counted, proving the second.
\end{proof}

\begin{remark}[Endpoint information and synchronization]\label{rem:endbudget61}
The unrelaxed inequality \eqref{eq:entropybudget61} also holds.
When $\zeta=1$, it gives $J\le1+L_k^{\rm ent}/\gamma$, consistent
with exponential exact width for extreme outputs. For fixed positive
error it gives the polynomial law, but does not classify every joint
limit $N\to\infty$, $\varepsilon\downarrow0$.

If the identity letter is absent but physical identity-product words
exist at every length at least $g$, greedily retain narrow cuts at
spacing at least $g+1$. Each retained cut accounts for at most $g+1$
original narrow cuts. One independent command followed by a return of
the remaining length fits each retained gap. The prefix and suffix
may be any fixed words of their required lengths. The same proof
multiplies the occupation bound by at most $g+1$. Approximate returns
with accuracy-dependent thresholds are not included in this quantitative
extension without additional estimates.
\end{remark}

\begin{corollary}[Matching power laws]\label{cor:matchedorbits61}
In the setting of Theorem~\ref{thm:uniformorbit57}, replace its converse
by \eqref{eq:sharpwidth61}. In particular, if the all-direction exponent
$p$ equals the real dimension $q$ of the target orbit, then
\[
 W_{N,\varepsilon}=\Theta((N+1)^{q/2})
\]
for every fixed $0<\rho<1$, $0\le\varepsilon<\rho$, and also for
$\rho=1$, $0<\varepsilon<1$. For a general fixed-point-free compact
connected Lie representation one may take $p=p_*$ from
Theorem~\ref{thm:minorbit58}; the lower and upper exponents are then
$p_*/2$ and $q/2$, respectively, with no logarithmic loss in the lower
bound.
\end{corollary}
\begin{proof}
Apply Theorem~\ref{thm:sharpoccupation61}. The upper common-row
construction in Theorem~\ref{thm:uniformorbit57} is exact for
$\rho<1$ and has $O((N+1)^{q/2})$ labels. At $\rho=1$ and positive
error use its amplitude-slack version. The two bounds match when
$p=q$. Theorem~\ref{thm:minorbit58} supplies the remaining assertion.
\end{proof}
The endpoint $\rho=1,\varepsilon=0$ is not an assertion of the matching
polynomial corollary. No exact leading constant or asymptotic ratio is
claimed. The spectral constant is effective only when a certified value
is supplied; a fully numerical example is given below.

\begin{remark}[Return-free extension]
Theorem~\ref{thm:driftoccupation62} below removes the identity-letter
and return assumptions from the occupation bound. It uses covariance
under the actual deterministic physical fillers, rather than treating
approximate returns as exact. The present synchronized theorem is its
special case; all smoothing and transport estimates remain unchanged.
\end{remark}

\section{Entropy transport without a return condition}\label{sec:drift62}
The entropy potential is invariant under a fixed physical rotation. This
allows arbitrary executable words between the spectral steps; they need
not return the physical system to its initial state. We first make the
covariance explicit and then retain the dependence on the error instead
of absorbing it into a fixed-error constant.

Keep the sphere $\Sigma=S(E)$, round probability measure, kernel $K_h$,
and constants $D_m,B_m$ from Lemma~\ref{lem:smoothentropy61}. For an
orthogonal map $B$, write $B_*\nu$ for transport of a measure and
$U_Bf(x)=f(B^{-1}x)$ for transport of a density. Directly from the
radial kernel and invariance of round measure,
\begin{equation}\label{eq:covariance62}
 K_h(B_*\nu)=U_B(K_h\nu),\qquad
 H(U_Bf)=H(f),\qquad
 W_2(B_*\nu,B_*\eta)=W_2(\nu,\eta).
\end{equation}
No commutation between $B$ and the averaging operator is asserted.

\begin{lemma}[Transport after a deterministic physical drift]
\label{lem:drifttransport62}
In Lemma~\ref{lem:centroidtransport61}, replace the identity-product
filler by any fixed executable word, with physical orthogonal action
$B$. If the resulting centroid masses satisfy $\alpha'\ge\zeta>0$,
then $\alpha'\le\alpha$ and
\begin{equation}\label{eq:drifttransport62}
 W_2\bigl(B_*(\lambda*\nu),\nu'\bigr)^2
             \le 3\pi^2\frac{\alpha-\alpha'}{\zeta}.
\end{equation}
The filler may have any length and its internal registers may have
arbitrary widths. A fixed physical word, with no random spectral step,
also cannot increase centroid mass.
\end{lemma}
\begin{proof}
Let $T_a(s,i)$ be the actual composite row for the random letter followed
by this fixed word. Conditional independence at the incoming register
and orthogonality give the exact identity
\[
 p'_i c'_i=\sum_{s,a}p_s\lambda_a T_a(s,i)BR(u_a)c_s.
\]
Use the weights $w_{s,a,i}=p_s\lambda_a T_a(s,i)\|c_s\|$ and unit
vectors $BR(u_a)c_s/\|c_s\|$ in the proof of
Lemma~\ref{lem:centroidtransport61}. Their weighted sum is the displayed
centroid, so its squared-distance identity and both couplings are
unchanged. They give \eqref{eq:drifttransport62}. Terms with zero mass
or zero centroid are omitted, exactly as in that lemma. For a fixed
word, conditional Jensen and the isometry $B$ give
$\sum_i p'_i\|c'_i\|\le\sum_s p_s\|Bc_s\|=\alpha$.
\end{proof}

Let $G$ be a compact connected Lie group with an orthogonal
representation $R:G\to O(E)$, $m=\dim E\ge2$. Fix a unit vector
$z_0$ whose orbit spans $E$, and suppose $E^G=\{0\}$. Put
$Q=\operatorname{conv}(R(G)z_0)$. The finite nonempty command alphabet
need not contain the identity or any exact return. Assume that a
probability law $\lambda$ on the \emph{actual letters} satisfies the
full action bound
\begin{equation}\label{eq:fullgap62}
 \|Pf-\Pi f\|_2\le\kappa\|f-\Pi f\|_2,\qquad
 Pf(x)=\sum_a\lambda_a f(R(u_a)^{-1}x),\qquad \kappa<1.
\end{equation}
Here $\Pi f(x)=\int_Gf(R(g)^{-1}x)\,dg$ projects onto all invariant
functions, not merely the constants. Suppose also that
\begin{equation}\label{eq:caps62}
 \sup_{u,v\in\Sigma}\int_G
   \boldsymbol1_{\{\|R(g)u-v\|<h\}}\,dg\le Ah^p
                         \quad(0<h\le h_0),
\end{equation}
where $p>0$. These are global action and all-direction hypotheses;
contraction on $E$ alone is not a substitute for \eqref{eq:fullgap62}.

\begin{theorem}[Return-free sharp occupation and error budget]
\label{thm:driftoccupation62}
Under \eqref{eq:fullgap62}--\eqref{eq:caps62}, consider the target
$F(g)=\rho R(g)z_0$, legal decoders in $Q$, and Euclidean error
$0\le\varepsilon<\rho\le1$. Set
\[
 \zeta=\rho-\varepsilon,\quad
 c_0=\min\{1,h_0,(2A)^{-1/p}\},\quad
 h=c_0 k^{-1/p},\quad
 \gamma=(1-\kappa^2)/2,\quad L_h=\log(D_mh^{-(m-1)}).
\]
For any error-$\varepsilon$ clocked realization, let
$J=\#\{1\le t\le N:|S_t|\le k\}$. If $J\ge1$ and $M=J-1$, then
\begin{equation}\label{eq:unrelaxed62}
 \gamma M\le L_h+
 \frac{\pi\sqrt{3B_m}}{h}
                   \sqrt{M\frac{1-\zeta}{\zeta}}.
\end{equation}
Consequently, as a real bound before taking an integer ceiling,
\begin{equation}\label{eq:driftoccupation62}
 J\le 1+\frac{2L_h}{\gamma}
          +\frac{3\pi^2B_m}{\gamma^2h^2}\frac{1-\zeta}{\zeta}.
\end{equation}
No condition is imposed on return words or on intervening register
widths. Constants other than $\zeta$ depend only on
$m,A,p,h_0,\kappa$ and the fixed spherical metric. In particular,
for fixed $\zeta>0$, $J=O_\zeta(k^{2/p})$ and
$W_{N,\varepsilon}=\Omega_\zeta(N^{p/2})$.
\end{theorem}
\begin{proof}
If $J\le1$, the asserted real bound is immediate. Otherwise list the
narrow cuts as $t_1<\cdots<t_J$. Choose an arbitrary deterministic
prefix of length $t_1$. In each gap put first one independent letter
of law $\lambda$, and then a fixed word of length $t_{j+1}-t_j-1$.
Such a word exists at every length because the alphabet is nonempty:
repeat any one of its letters. Finish with any fixed suffix of length
$N-t_J$. This external word law depends on the advertised width profile,
not on the machine's private state. Write $B_j$ for the physical action
of the fixed filler in gap $j$. The physical state at every cut remains
one unit orbit point; no equality of the $B_j$ to the identity is needed.

Let $\alpha_j,\nu_j$ be its conditional-centroid mass and directional
law at cut $t_j$. Terminal correctness gives $\alpha_N\ge\zeta$ by
exactly the inner-product argument in Theorem~\ref{thm:sharpoccupation61}.
This uses wordwise correctness of the mean decoder, not correctness of
each label. Lemma~\ref{lem:drifttransport62} and the fixed-suffix
inequality imply
\[
 1\ge\alpha_1\ge\cdots\ge\alpha_J\ge\zeta,
 \qquad\sum_{j=1}^{M}(\alpha_j-\alpha_{j+1})\le1-\zeta.
\]
There are at most $k$ directions at each selected cut. Thus
$f_j=K_h\nu_j$ is supported on at most $k$ open chordal caps, and
\eqref{eq:caps62} gives $\Pi\boldsymbol1_{\operatorname{supp}_+f_j}
\le kAh^p\le1/2$. The kernel vanishes at each cap boundary.
Lemma~\ref{lem:entropydefect61} and \eqref{eq:covariance62} therefore give
\[
 H\bigl(U_{B_j}P f_j\bigr)=H(P f_j)\le H(f_j)-\gamma.
\]
By covariance, $U_{B_j}P f_j=K_h(B_{j*}(\lambda*\nu_j))$.
Entropy continuity and Lemma~\ref{lem:drifttransport62}, with
$\delta_j=\alpha_j-\alpha_{j+1}$, yield
\[
 H(f_{j+1})\le H(f_j)-\gamma+
           \frac{\pi\sqrt{3B_m}}{h\sqrt\zeta}\sqrt{\delta_j}.
\]
Sum, use $0\le H(f_j)\le L_h$, and apply Cauchy--Schwarz to the
$M$ nonnegative losses. This proves \eqref{eq:unrelaxed62}.
The inequality $ab\le(a^2+b^2)/2$, with
$a=\sqrt{\gamma M}$, gives \eqref{eq:driftoccupation62}.
For a full peak bound $k$ all $N$ positive cuts are counted. Since
$\log(k+1)=O(k^{2/p})$, the last conclusions follow. The constants
are independent of the filler lengths, their actions, and all actual
hidden tables.
\end{proof}

\begin{remark}[What has, and has not, been removed]\label{rem:gapboundary62}
The return hypothesis has disappeared from the quantitative converse,
not from the general terminal stationarization theorem. The full action
gap remains indispensable to this argument. A single irrational rotation
letter has a unitary Koopman operator and cannot satisfy
\eqref{eq:fullgap62}; its one-label clocked realization is therefore no
counterexample. No entropy is assigned to the hidden register here:
$H$ is relative entropy of an auxiliary smoothed directional law against
normalized round measure. The classical entropy--transport principle
\cite{PW61} supplies neither the Markov row identity nor the width
quantifiers. The logarithmic initial range $L_h$ is retained in
\eqref{eq:unrelaxed62}; it is not paid again at every narrow cut.
\end{remark}

\begin{corollary}[A fixed executable spectral block]\label{cor:blockgap62}
The same conclusions hold if \eqref{eq:fullgap62} is supplied by a
probability law on executable words of one fixed length $b\ge1$, rather
than individual letters. In \eqref{eq:driftoccupation62} multiply the
right side by $b$.
\end{corollary}
\begin{proof}
Greedily select narrow cuts at separation at least $b$. Each selected
cut accounts for at most $b$ original cuts, including the last incomplete
group of cuts. Apply the proof above with one independent spectral word
at the start of every selected gap and an arbitrary fixed filler for its
remaining length. There is no restriction before the first selected cut
or after the last. The alphabet of spectral blocks is finite even when
their internal hidden widths are unbounded.
\end{proof}

\begin{theorem}[A joint amplitude--accuracy law]\label{thm:jointlaw62}
Assume the preceding geometric and action-gap hypotheses. Let
$q=\dim R(G)z_0$, and write $W_N(\rho,\varepsilon)$ to display the
amplitude dependence. Fix $\zeta_0>0$. Uniformly for
$0\le\varepsilon<\rho\le1$ with $\rho-\varepsilon\ge\zeta_0$ and
$\tau=1-\rho+\varepsilon>0$, constants $C_i$ depending only on the
fixed data and $\zeta_0$ give
\begin{align}
 N&\le C_0+C_1\log k+C_2\tau k^{2/p},
                       &&k=W_N(\rho,\varepsilon),\label{eq:jointlower62}\\
 W_N(\rho,\varepsilon)&\le C_3\left(1+\frac{N+1}{\tau}\right)^{q/2}.
                                                        \label{eq:jointupper62}
\end{align}
If $p=q$ and a parameter sequence satisfies
$\rho_N-\varepsilon_N\ge\zeta_0$ and
$\log(1/\tau_N)=o(N)$, then
\begin{equation}\label{eq:jointmatched62}
 W_N(\rho_N,\varepsilon_N)
            =\Theta\left((N/\tau_N)^{p/2}\right)\quad(N\longrightarrow\infty).
\end{equation}
In particular, at pure amplitude and any positive subexponentially
vanishing Euclidean tolerance, the matched law is
$\Theta((N/\varepsilon_N)^{p/2})$. It includes polynomial and
stretched-exponential accuracy sequences.
\end{theorem}
\begin{proof}
In \eqref{eq:driftoccupation62}, $1-\zeta=\tau$ and
$\zeta^{-1}\le\zeta_0^{-1}$. Also
$L_h=\log(D_mc_0^{-(m-1)})+(m-1)\log k/p$. Counting all positive cuts
proves \eqref{eq:jointlower62}.

For the upper bound set $\beta=\rho-\varepsilon/2$. Then
$0<\beta<1$, $1-\beta\ge\tau/2$, and the target change has norm at
most $\varepsilon/2$. The proof of the upper construction in
Theorem~\ref{thm:uniformorbit57} uses only the invariant orbit hull and
Lemma~\ref{lem:orbithull57}; it does not use its identity-letter or
spectral assumptions. Applied at amplitude $\beta$, it gives an exact
common-row machine with at most
$C(1+(N+1)/(-\log\beta))^{q/2}$ labels. The inequality
$-\log\beta\ge1-\beta\ge\tau/2$ proves \eqref{eq:jointupper62}.
When $\varepsilon=0$, this is the same exact construction at
$\beta=\rho<1$; no limiting construction at $\beta=1$ is used.

Let $k_N$ be the minimum width. The uniform upper bound gives
$\log k_N=O(1+\log(N+1)+\log(1/\tau_N))=o(N)$.
Consequently the first two terms of \eqref{eq:jointlower62} sum to at
most $N/2$ for all sufficiently large $N$. It follows that
$k_N\ge(N/(2C_2\tau_N))^{p/2}$. The matching upper bound when $p=q$
proves \eqref{eq:jointmatched62}. The comparison constants depend only
on the fixed data and $\zeta_0$; the starting horizon can depend on the
parameter sequence.
\end{proof}

The least-orbit theorem supplies $p=p_*$ for each fixed-point-free
compact connected Lie representation. Thus the joint law applies when
$p_*=q$. The point $\tau=0$ is precisely pure amplitude and zero error;
there \eqref{eq:unrelaxed62} instead gives $\gamma(J-1)\le L_h$.
Neither the joint theorem nor its proof determines optimal constants
or the entire exponentially small transition window.

\section{Positive matrix realizations in every dimension}\label{sec:matrix57}
Let $d\ge2$ and let $\mathcal H_d^0$ be the real Euclidean space of traceless Hermitian $d\times d$ matrices, with scalar product $\langle A,B\rangle=\operatorname{tr}(AB)$ and Frobenius norm $\|\cdot\|_{\rm F}$. Let
\[
 \mathcal D_d=\{D=D^*\succeq0:\operatorname{tr}D=1\},\quad
 c_d=I/d,\quad r_d=\sqrt{1-1/d}.
\]
This section concerns numerical outputs in $\mathcal D_d$. The requirement that every decoder value be positive with trace one is part of the model; it is not imposed merely on the final average.

\begin{lemma}[Small balls on the complex Grassmannian]\label{lem:grassmann57}
Let $P$ be a fixed rank-$m$ orthogonal projection in $\C^d$, with $1\le m<d$. For Haar $U\in\mathrm{SU}(d)$ there is $C_{d,m}<\infty$ such that
\[
 \Pr\{\|UPU^*-Q\|_{\rm F}<h\}\le C_{d,m}h^{2m(d-m)}\quad(0<h\le1)
\]
for every rank-$m$ projection $Q$. For rank one, writing
$\operatorname{dist}_{\rm pr}(P,Q)=\sqrt{1-\operatorname{tr}(PQ)}$, one has $\operatorname{dist}_{\rm pr}(P,Q)=\norm{P-Q}_{\rm F}/\sqrt2$ and exactly
\begin{equation}\label{eq:projectivecap57}
 \Pr\{\operatorname{dist}_{\rm pr}(UPU^*,Q)<h\}=h^{2(d-1)}
                         \qquad(0\le h\le1).
\end{equation}
\end{lemma}
\begin{proof}
Conjugation is transitive, so fix $Q=\operatorname{diag}(I_m,0)$. A projection sufficiently close to $Q$ has range the graph of a unique matrix $Z\in\C^{(d-m)\times m}$, with
\[
 P(Z)=\begin{pmatrix}I\\Z\end{pmatrix}(I+Z^*Z)^{-1}(I\quad Z^*).
\]
This notation means the product of the three displayed matrices. The map and its inverse are smooth near zero, and its derivative at zero is $Z\mapsto\left(\begin{smallmatrix}0&Z^*\\Z&0\end{smallmatrix}\right)$, of Frobenius norm $\sqrt2\|Z\|_{\rm F}$. The invariant probability measure is normalized Riemannian volume on this compact homogeneous manifold; in this chart its density is smooth and bounded on a sufficiently small closed neighborhood. The inverse chart puts a Frobenius $h$-ball inside a Euclidean $Ch$-ball in real dimension $2m(d-m)$. This proves the bound for small $h$; increasing the constant covers the remaining interval. Transitivity makes the constant independent of $Q$.

For rank one choose $Q=e_1e_1^*$. A uniform unit vector can be obtained by normalizing independent standard complex Gaussian coordinates. If $X$ is the squared modulus of the first coordinate before normalization and $Y$ is the sum of the other squared moduli, then $X$ is exponential of mean one and $Y$ is an independent sum of $d-1$ such variables. For $0<z<1$,
\[
 \Pr\left\{\frac{X}{X+Y}>1-z\right\}
 =\E\exp\left(-\frac{1-z}{z}Y\right)=z^{d-1}.
\]
Put $z=h^2$. The endpoint cases follow by continuity. This proves \eqref{eq:projectivecap57}, and also shows that $\operatorname{dist}_{\rm pr}$ is the Frobenius chordal distance divided by $\sqrt2$.
\end{proof}

\begin{lemma}[A uniform conjugacy-orbit exponent]\label{lem:conjugacysmall57}
There is $C_d<\infty$ such that
\begin{equation}\label{eq:conjugacycap57}
 \sup_{\substack{A,B\in\mathcal H_d^0\\\|A\|_{\rm F}=\|B\|_{\rm F}=1}}
 m_G\{U:\|UAU^*-B\|_{\rm F}<h\}
                 \le C_d h^{2(d-1)}\qquad(0<h\le1).
\end{equation}
The exponent cannot be increased uniformly over $A$ and $B$.
\end{lemma}
\begin{proof}
Write the eigenvalues of $A$ as $a_1\ge\cdots\ge a_d$. Since their sum is zero and their squared sum is one, $a_1-a_d\ge d^{-1/2}$. Some adjacent gap therefore satisfies
\begin{equation}\label{eq:separatedcut57}
 a_m-a_{m+1}\ge\gamma_d:=\frac1{\sqrt d(d-1)}.
\end{equation}
Let $P_A$ be the spectral projection onto the first $m$ eigenvectors. It is well defined despite possible multiplicities within either block.

Two conjugates $A',A''$ of $A$, with eigenbases $x_i,y_j$ in this order, obey the exact identity
\begin{equation}\label{eq:spectralidentity57}
 \|A'-A''\|_{\rm F}^2
       =\sum_{i,j}(a_i-a_j)^2|\langle x_i,y_j\rangle|^2.
\end{equation}
Expanding both traces proves it. The cross-block terms and \eqref{eq:separatedcut57} imply
\begin{equation}\label{eq:projectorbound57}
 \|P_{A'}-P_{A''}\|_{\rm F}
                    \le\gamma_d^{-1}\|A'-A''\|_{\rm F},
\end{equation}
since the squared norm on the left is exactly the sum of the two cross-block overlap sums.

If the set in \eqref{eq:conjugacycap57} is nonempty, choose one member $U_0$. Every other member has $\|UAU^*-U_0AU_0^*\|_{\rm F}<2h$. Thus its rank-$m$ spectral projection lies within $2h/\gamma_d$ of $U_0P_AU_0^*$. The Haar law of $UP_AU^*$ is the invariant Grassmannian law. Lemma~\ref{lem:grassmann57} bounds the measure by
$C_{d,m}(2h/\gamma_d)^{2m(d-m)}$ for $h\le\gamma_d/2$. Since $m(d-m)\ge d-1$, and there are finitely many $m$, this gives \eqref{eq:conjugacycap57} with a constant independent of all eigenvalue multiplicities. Larger $h$ are covered by increasing $C_d$.

For sharpness take $A=B=(P-c_d)/r_d$ with $P$ rank one. By \eqref{eq:projectivecap57}, for sufficiently small $h$ the probability is
\[
                    (r_d^2h^2/2)^{d-1}.
\]
It cannot be bounded by a constant times $h^{p'}$ as $h\downarrow0$ for any $p'>2(d-1)$.
\end{proof}
The spectral split, rather than an assumption that all centroids remain rank one, supplies the uniformity needed by the stochastic contraction. The conjugation representation has ambient dimension $d^2-1$, but its worst small-ball exponent is $2(d-1)$.

\begin{theorem}[Unitary matrix-orbit width]\label{thm:matrixwidth57}
Let a finite alphabet in $\mathrm{SU}(d)$ contain the identity and support a probability law satisfying \eqref{eq:groupgap57}. Fix a rank-one projection $P_0$. For $0<a\le1$, set
\begin{equation}\label{eq:matrixexperiment57}
 F_a(U)=c_d+a(UP_0U^*-c_d),\qquad Y=\mathcal D_d,
\end{equation}
with Frobenius error. Its exact bounded-width threshold is $a r_d$, attained by the one-label decoder $c_d$. For $0\le\varepsilon<a r_d$, put $\zeta=a-\varepsilon/r_d$. Every realization satisfies
\begin{equation}\label{eq:matrixoccupation57}
 \#\{1\le t\le N:|S_t|\le k\}
 \le C\left(1+\log(k+1)+\frac{1-\zeta}{\zeta}k^{1/(d-1)}\right).
\end{equation}
For fixed $0<a<1$ and $0\le\varepsilon<a r_d$,
\begin{equation}\label{eq:matrixrate57}
 cN^{d-1}
            \le W_{N,\varepsilon}\le C(N+1)^{d-1},
\end{equation}
and the upper realization can be exact. The same orders hold for $a=1$ and every fixed $0<\varepsilon<r_d$.

More explicitly, for $a<1$ and $\varepsilon=0$ set $a'=a$; for $\varepsilon>0$ set $a'=a-\varepsilon/(2r_d)$. Whenever $a'<1$, an error at most $\varepsilon/2$ realization has
\begin{equation}\label{eq:matrixupper57}
 K\le C_d\left(1+\frac{N+1}{-\log a'}\right)^{d-1}
\end{equation}
labels. In the zero-error case this realization is exact. In \eqref{eq:matrixoccupation57} the constant depends only on $d$ and the spectral-gap data, not on the amplitude or tolerance. Extra finitely many unitary commands are allowed in both bounds.
\end{theorem}
\begin{proof}
Center at $c_d$ and divide all outputs by $r_d$. The resulting legal set is
\[
 Q=\operatorname{conv}\{(P-c_d)/r_d:P\text{ rank one}\}\subseteq\mathbb B_{\mathcal H_d^0}.
\]
The last inclusion follows from $\operatorname{tr}D^2\le1$ for $D\in\mathcal D_d$. The orbit spans $\mathcal H_d^0$: if a traceless Hermitian $A$ has $\operatorname{tr}(AP)=0$ for every rank-one $P$, then $v^*Av=0$ for every $v$, hence $A=0$. Conjugation-fixed matrices are scalar, so there are no fixed vectors in this traceless space. The orbit is complex projective space, of real dimension $2(d-1)$. Lemma~\ref{lem:conjugacysmall57} and Theorem~\ref{thm:sharpoccupation61}, with $p=q=2(d-1)$, prove the lower and occupation bounds, as well as the radius statement after multiplying the norm by $r_d$.

We give a direct inner approximation to make positivity of the upper rows transparent. Choose a maximal $\delta$-separated family of rank-one projections $P_i$ for $\operatorname{dist}_{\rm pr}$, containing $P_0$. Its $\delta/2$ balls are disjoint, so \eqref{eq:projectivecap57} gives
\[
                         K\le(2/\delta)^{2(d-1)}.
\]
For any traceless Hermitian $A$, put $M=\lambda_{\max}(A)$. When $A\ne0$, $M>0$ and $\lambda_{\min}(A)\ge-(d-1)M$. Approximate a top eigenline by $P_i$ at distance at most $\delta$. Spectral decomposition then gives
\begin{equation}\label{eq:matrixsupport57}
 \operatorname{tr}(AP_i)\ge M-(M-\lambda_{\min}(A))\delta^2
                         \ge(1-d\delta^2)M.
\end{equation}
The trivial $A=0$ case is automatic. Comparing support functions in $\mathcal H_d^0$, with $K_0=\mathcal D_d-c_d$, yields
\[
 (1-d\delta^2)K_0\subseteq\operatorname{conv}\{P_i-c_d\}\subseteq K_0.
\]
Put $r=1-d\delta^2>0$. For each actual command $U_b$ choose one row of barycentric coordinates satisfying
\begin{equation}\label{eq:matrixrow57}
 \sum_jT_b(i,j)(P_j-c_d)=rU_b(P_i-c_d)U_b^*.
\end{equation}
Rows are nonnegative and sum to one by construction. Initialize at $P_0$ and decode label $i$ as
\[
                  D_i=c_d+(a'/r^N)(P_i-c_d).
\]
This is in $\mathcal D_d$ when $r^N\ge a'$: it is a convex combination of $c_d$ and $P_i$. Choosing
\[
 \delta^2=\min\left(\frac1{4d},\frac{-\log a'}{2d(N+1)}\right)
\]
ensures this inequality. Induction in \eqref{eq:matrixrow57} gives mean output exactly $F_{a'}(u_w)$. The target discrepancy is $|a-a'|r_d=\varepsilon/2$ when $\varepsilon>0$. The cardinal bound proves \eqref{eq:matrixupper57}. The case $a=1,\varepsilon=0$ is deliberately excluded from this dilation construction and is treated exactly in the next section.
\end{proof}
The family contains the qubit sphere as $d=2$, but for $d\ge3$ the pure-state orbit is a proper curved subset of the unit sphere in $\mathcal H_d^0$. Thus \eqref{eq:matrixrate57} does not follow by substituting the ambient dimension into the spherical theorem. The entropy argument removes the logarithmic loss. No optimal leading constant, exact asymptotic ratio, or dimension-uniform constant is claimed.

\begin{corollary}[Algebraic generating alphabets]\label{cor:algebraicunitary57}
Theorem~\ref{thm:matrixwidth57} applies to the identity together with any finite symmetric algebraic-entry set generating a dense subgroup of $\mathrm{SU}(d)$, equipped with a lazy symmetric law. Such a finite alphabet exists for every $d\ge2$.
\end{corollary}
\begin{proof}
The spectral gap is the imported theorem of Bourgain--Gamburd \cite{BG2012}. Laziness makes its one-sided gap an absolute $L^2$ contraction on mean-zero functions. Existence can be seen by placing two algebraic irrational-angle $\mathrm{SU}(2)$ rotations in every adjacent pair of coordinate directions. For example, use the quaternion lifts in the explicit adjacent-block construction below within each pair. Each pair generates its full $\mathrm{SU}(2)$ subgroup. Their Lie algebras contain the adjacent off-diagonal matrix units and their brackets, which generate $\mathfrak{su}(d)$. The connected compact closure is consequently $\mathrm{SU}(d)$. All entries remain algebraic. The spectral-gap constant is not computed by the finite regression checks.
\end{proof}

\subsection{A finite whole-law encoding}
Let $q_0=d^2-1$ and choose a Frobenius-orthonormal basis $H_1,\ldots,H_{q_0}$ of $\mathcal H_d^0$. Since $\|H_l\|_{\rm op}\le1$, the $2q_0$ matrices
\[
 E_{l,\pm}=\frac{I\pm H_l/2}{2q_0}
\]
are positive and sum to $I$. Put $\Phi(D)_{l,\pm}=\operatorname{tr}(E_{l,\pm}D)$ and $Y_\Phi=\Phi(\mathcal D_d)$. Every entry is at least $1/(4q_0)$. The map is injective, because
\begin{equation}\label{eq:tomography57}
 D-c_d=\sum_l2q_0\bigl(\Phi(D)_{l,+}-\Phi(D)_{l,-}\bigr)H_l.
\end{equation}
On the vector space parallel to the affine span of $Y_\Phi$, use the norm of the reconstructed matrix on the right. Then $\Phi$ is an affine isometry, and replacing each legal matrix decoder by its image preserves every clocked or stationary width exactly. This is one whole categorical law with a compact convex legality constraint, not separate binary event tolerances. For ordinary total variation the explicit norm comparison is
\[
 \frac{\|D-D'\|_{\rm F}}{4q_0}
 \le\TV(\Phi(D),\Phi(D'))
 \le\frac{\|D-D'\|_{\rm F}}{4\sqrt{q_0}}.
\]
It follows directly by comparing the $\ell^1$ and $\ell^2$ norms of the basis coefficients. Enlarging the legal output set from $Y_\Phi$ to the entire simplex is a different realization problem. In particular the extreme-output argument below cannot be transferred to that enlarged set: its target vectors are interior points of the simplex.

\begin{remark}[An explicit algebraic alphabet in every dimension]
In each adjacent coordinate plane $(j,j+1)$ embed the two matrices
\[
 \frac15\begin{pmatrix}-3+4i&0\\0&-3-4i\end{pmatrix},\qquad
 \frac15\begin{pmatrix}-3&4\\-4&-3\end{pmatrix},
 \qquad 1\le j<d,
\]
and act identically on all other coordinates. Include their inverses and
the identity. Each pair has infinite-order eigenvalues, since its trace
$-6/5$ is not an integer, and its two torus closures generate the adjacent
$SU(2)$. The adjacent skew-Hermitian matrix units and their brackets
span $\mathfrak{su}(d)$. The closed generated subgroup is therefore
$SU(d)$. Bourgain--Gamburd is applied to this actual finite algebraic
set, not merely to the existence of an unspecified set. The resulting
spectral-gap constant is not numerically evaluated here.
\end{remark}
The Frobenius chordal distance between rank-one projections is $\sqrt2$
times the projective distance used in the net construction. In the
conjugacy-cap comparison, choosing one member of a nonempty target ball
avoids imposing any spectral condition on its arbitrary center. All
adjacent-gap and cap constants here depend on $d$. The amplitude
$ a'=a-\varepsilon/(2r_d)$ in the positive-error upper bound deliberately
spends only half the available error.

\section{Extreme outputs and the exact endpoint}\label{sec:extreme57}
The strict amplitude slack in the polytope construction is not a technical
convenience that can simply be removed. At an extreme output, convex
averaging has no room to hide a different physical state. This observation
gives an exact finite-horizon formula, independently of spectral gap or
synchronization.

\begin{theorem}[The extreme-orbit formula]\label{thm:extreme57}
Let a group act by affine bijections on a compact convex set $Y$. Let
$x_0\in Y$ have an orbit consisting of extreme points of $Y$, and let a
finite nonempty command alphabet act through this group. There is one seed
and one terminal query, with target $u_w x_0$. Put
\[
 \mathcal O_t=\{u_wx_0:|w|=t\},\qquad n_t=|\mathcal O_t|.
\]
Every exact clocked stochastic realization at horizon $N$ satisfies
\begin{equation}\label{eq:extremeprofile57}
                         |S_t|\ge n_t\quad(0\le t\le N).
\end{equation}
One deterministic realization attains all these inequalities simultaneously.
Consequently
\begin{equation}\label{eq:extremewidth57}
                         W_{N,0}=n_N.
\end{equation}
Neither an identity letter nor an exact return is required.
\end{theorem}
\begin{proof}
Fix $t$ and one suffix $v$ of length $N-t$. Such a suffix exists because the
alphabet is nonempty; at $t=N$ use the empty word. Let $D_v(i)\in Y$ be the
conditional expected terminal decoder after this fixed suffix, when the
cut-$t$ label is $i$. Convexity guarantees that it lies in $Y$. For every
prefix $w$ of length $t$, exact correctness says
\[
             u_vu_wx_0=\sum_i p_w(i)D_v(i),
\]
where $p_w$ is the actual distribution at that cut. Extremality forces
$D_v(i)=u_vu_wx_0$ whenever $p_w(i)>0$. Since the action of $u_v$ is
injective, prefixes with different physical outputs have disjoint
nonempty supports in $S_t$. This proves \eqref{eq:extremeprofile57}, including
the terminal cut.

Conversely, at cut $t$ store the current point of $\mathcal O_t$. A command
sends this point to its image in $\mathcal O_{t+1}$; the terminal decoder is
the stored point. The initialization is deterministic. Finally, any one
fixed command gives an injection $\mathcal O_t\to\mathcal O_{t+1}$, so $n_t$
is nondecreasing. The peak of this realization is $n_N$, proving
\eqref{eq:extremewidth57}.
\end{proof}
The same conclusion holds for an orbit on a Euclidean unit sphere with
legal output set its convex hull: every orbit point is extreme by strict
convexity of the ambient Euclidean ball. For density matrices, a rank-one
projection is extreme. Indeed, if $P=\sum_i p_iD_i$ with $D_i\ge0$ and
$\tr D_i=1$, then $\tr((I-P)D_i)=0$ for every positive weight. Positivity
forces $D_i$ to be supported on the range of $P$, and hence $D_i=P$.


The cut argument fixes one suffix, whose affine action is injective.
The cardinalities are nondecreasing because the nonempty alphabet
contains a letter whose injective action maps each reachable set into
the next one. No sampled-answer floor is included in this width.

\section{Effective algebraic-circle bounds}\label{sec:arithmetic53}
For algebraic unit-circle eigenvalues of infinite order, an explicit separation estimate replaces the Diophantine hypothesis of the preceding section. The parameters $r,m,A_r$ below are those of \eqref{eq:r53}--\eqref{eq:circleconstants53}; in this section $r$ denotes the localization order, not the number of command rotations. The alphabet contains $I,R_\theta$, and one interface mean is $\rho\cos\phi$.

\begin{lemma}[Separated powers give a block defect]\label{lem:separated53}
Suppose $|\lambda^n-1|\ge\sigma$ for every integer $1\le n\le2mk$, where $0<\sigma\le1$. Write $R_m(t)=\diag(R_t,R_{2t},\ldots,R_{mt})$ for the representation acting on $Z_m$. On this representation, the uniform law on the $2k$ words with products $R_\theta^j$, $0\le j<2k$, each padded to length $2k$, has defect at least $\sigma^2/16$ in \eqref{eq:gap54}.
\end{lemma}
\begin{proof}
Write a unit vector as $u=(u_1,\ldots,u_m)$ in its two-dimensional frequency blocks. For distinct indices $i,j\in\{0,\ldots,2k-1\}$,
\[
 \norm{R_m(i\theta)u-R_m(j\theta)u}^2
 =\sum_{l=1}^m\norm{u_l}^2|\lambda^{l(i-j)}-1|^2\ge\sigma^2.
\]
A ball of radius strictly less than $\sigma/2$ about any of $k$ centers contains at most one of these $2k$ points. Thus at least $k$ points have distance at least $\sigma/2$ from all centers (the same conclusion follows by taking the open balls of radius $\sigma/2$). For unit centers the defect is one half of the squared distance to the closest center. Averaging gives at least $(1/2)(\sigma^2/8)=\sigma^2/16$.
\end{proof}

\begin{theorem}[Full-noise algebraic-circle bound]\label{thm:algebraic53}
Assume $\lambda$ is algebraic and not a root of unity. Let its primitive irreducible integer polynomial be
\[
 P(z)=a z^d+a_{d-1}z^{d-1}+\cdots+a_0,\qquad a>0.
\]
Set $H_P=\max(a,|a_{d-1}|,\ldots,|a_0|)$ and
\[
                     B_P=a(1+H_P/a)^{d-1}>1.
\]
With $r,m,A_r$ as in \eqref{eq:r53}--\eqref{eq:circleconstants53}, every error-$\varepsilon$ realization of peak at most $k$ satisfies
\begin{equation}\label{eq:algwidth53}
 N<2k\left(1+16\,2^{2(d-1)}B_P^{4mk}\log A_r\right).
\end{equation}
Its number of command cuts $t\in\{1,\ldots,N\}$ with width at most $k$, even if other widths are unrestricted, is at most
\begin{equation}\label{eq:algoccupation53}
 2k-1+32k\,2^{2(d-1)}B_P^{4mk}\log A_r.
\end{equation}
Consequently, for fixed algebraic data, $\rho$, and $\varepsilon<\rho/2$,
\begin{equation}\label{eq:algrate53}
 W_{N,\varepsilon}\ge
 \frac{\log N-O(\log\log N)}{4m\log B_P}.
\end{equation}
All constants in the finite inequalities are explicit. If $\lambda=(a_1+ib_1)/c$ with integers $a_1,b_1$, $c>0$, $a_1^2+b_1^2=c^2$, and $\lambda$ not a root of unity, then \eqref{eq:algwidth53}--\eqref{eq:algrate53} hold with the factor $2^{2(d-1)}$ deleted and $B_P$ replaced by $c$.
\end{theorem}
\begin{proof}
List the conjugates with $\lambda_1=\lambda$. They are distinct by irreducibility in characteristic zero. None is a root of unity: otherwise irreducibility would make $P$ a cyclotomic polynomial up to its primitive sign. The resultant
\[
 \operatorname{Res}(P,z^n-1)=a^n\prod_{j=1}^d(\lambda_j^n-1)
\]
is therefore a nonzero integer. Each conjugate has modulus at most $R_P=1+H_P/a$. For completeness, if $|z|>R_P$, then
\[
 \sum_{j=0}^{d-1}|a_j||z|^j
 <\frac{H_P|z|^d}{|z|-1}<a|z|^d,
\]
so $P(z)\ne0$. Using $|z^n-1|\le2\max(1,|z|)^n$ for each other conjugate, the resultant bound gives
\begin{equation}\label{eq:resultant53}
 |\lambda^n-1|\ge
 2^{-(d-1)}\left(a\prod_{j=2}^d\max(1,|\lambda_j|)\right)^{-n}
 \ge 2^{-(d-1)}B_P^{-n}.
\end{equation}
For $1\le n\le2mk$, take $\sigma=2^{-(d-1)}B_P^{-2mk}$ in Lemma~\ref{lem:separated53}. The block defect is at least
\[
                  d_k=\frac{2^{-2(d-1)}B_P^{-4mk}}{16}.
\]
Repeated conditional contraction and Lemma~\ref{lem:harmonic53} imply
\[
 \left\lfloor\frac{N}{2k}\right\rfloor
 \le \frac{\log A_r}{-\log(1-d_k)}
 \le 16\,2^{2(d-1)}B_P^{4mk}\log A_r.
\]
This proves \eqref{eq:algwidth53}. The greedy occupation argument in the proof of Theorem~\ref{thm:radius54}, with $b=0$ and block length $2k$, proves \eqref{eq:algoccupation53}. Taking logarithms gives \eqref{eq:algrate53}: fix $C>1/(4m\log B_P)$. If $k\ge C\log N$ the claim is immediate; otherwise $\log k\le\log C+\log\log N$, which can be absorbed in the logarithm of \eqref{eq:algwidth53}. The implicit constant may depend on $P,\rho,\varepsilon$, not on $N$.

In the rational-circle case, $(a_1+ib_1)^n-c^n$ is a nonzero Gaussian integer, of modulus at least one. Thus $|\lambda^n-1|\ge c^{-n}$ directly. Repeat the proof with $\sigma=c^{-2mk}$. Here $c>1$ automatically. A nonminimal supplied denominator is still valid but weakens the bound.
\end{proof}
The threshold is sharp; these width rates are not asserted sharp. The factor $m$ records the increased localization degree as the error approaches the boundary. At smaller errors, the retained coordinate-calibration estimates can have better constants. One should use the larger of applicable lower bounds, rather than claiming numerical improvement at every parameter value.

\begin{proposition}[An explicit nonrational algebraic example]\label{prop:salem53}
Let $t=(1-\sqrt{13})/2$ and
\[
 \lambda=\frac{t+i\sqrt{4-t^2}}2.
\]
Then $|\lambda|=1$, $\lambda$ is not a root of unity, and one can take $d=4$ and $B_P=8$ in Theorem~\ref{thm:algebraic53}.
\end{proposition}
\begin{proof}
The number $t$ lies strictly between $-2$ and $2$ and solves $t^2-t-3=0$. Substitution of $t=z+z^{-1}$ gives
\[
                       P(z)=z^4-z^3-z^2-z+1.
\]
Its reduction modulo two is $z^4+z^3+z^2+z+1$. It has no linear factor over $\mathbb F_2$ and is not divisible by the only monic irreducible quadratic $z^2+z+1$; hence it is irreducible. Gauss's lemma proves irreducibility over $\mathbb Q$. The conjugate value $(1+\sqrt{13})/2>2$ gives a positive real conjugate of $\lambda$ greater than one. A root of unity cannot have such a conjugate. Finally $a=H_P=1$ and $d=4$, so $B_P=2^3=8$.
\end{proof}
For a reproducible rational calibration, take $\rho=1/10$ and $\varepsilon=1/25$. Then $\delta=2/25$, $r=5$, $m=6$, $\Delta_r=1/150$, and $A_r=23364$. Thus the rational-circle alphabet above obeys
\[
                 N<2k\bigl(1+16\cdot5^{24k}\log23364\bigr).
\]
This is a concrete bound inside the interval between the former small-error obstruction and the sharp threshold $1/20$. The exact coefficient and resultant checks accompanying the article verify the displayed arithmetic, not the universal theorem by finite experimentation.


\section{Finite encodings and boundary quotients}\label{sec:effective54}
The general continuous theorem is not an algorithm from arbitrary real data. We now give an explicit finite input class for which its conclusions are effective. A rational orthogonal input consists of a finite alphabet in $O(D,\Q)$ containing $I$, finite seed and query sets, and rational polynomials $F_{s,j,l}$ in the matrix entries specifying categorical probabilities. The promise $F_{s,j}(H)\subset\Delta_{M_j}$ can itself be checked once the compact closure is computed. Real algebraic output numbers are encoded by polynomial equations and isolating data.

We use two established algorithmic inputs: computation of a finite defining description of the Zariski closure of a finitely generated rational matrix group \cite{DJK,NPSW}, and real quantifier elimination with effective semialgebraic connected components \cite{BPR}. The former is a terminating closure algorithm, not merely an enumeration of valid polynomial identities with an unknown stopping point. Compact orthogonal groups are real algebraic, so the real points of the computed closure in $O(D)$ give the Euclidean closure $H$ \cite{BJKP}. The complex algebraic components must not be substituted for the real connected components used here.

\begin{theorem}[Effective rational classification]\label{thm:effective54}
For a rational orthogonal input, the following are computable:
\begin{enumerate}
\item a semialgebraic description of every connected component of $H$, the finite group $C$, and the component of each command;
\item the critical total-variation error $\epsc$ as an exact real algebraic number, and an attaining stationary component-permutation machine with algebraic decoders;
\item for each algebraic $0\le\varepsilon<\epsc$ and integer $k\ge1$, an integer occupation constant $L_k$ valid for every error-$\varepsilon$ clocked realization;
\item for every fixed $N$ and algebraic $\varepsilon\ge0$, the least width $W_{N,\varepsilon}$ and an algebraic realization attaining it.
\end{enumerate}
These are terminating computability statements. No polynomial running-time bound, practical implementation of general closure/quantifier elimination, or finite-random-bit cost bound is asserted.
\end{theorem}
\begin{proof}
Compute $H$ using the cited closure algorithm, then its real connected components and algebraic sample points. The component containing $I$ is $K$. Products of sample points and semialgebraic membership tests determine the component multiplication table and the command images. In particular $K$ is now an effectively given compact semialgebraic group.

For a component $c$ and a query, the condition that $q$ be a radius-$t$ center is the real first-order formula
\begin{equation}\label{eq:QEcenter54}
 q\in\Delta_{M_j},\qquad
 \forall h\in c:\quad
 \sum_{l=1}^{M_j}|F_{s,j,l}(h)-q_l|\le2t.
\end{equation}
Absolute values have a finite semialgebraic description. Quantifier elimination gives the set of admissible $(q,t)$. Its least $t$ exists by compactness and is real algebraic; algebraic sampling supplies an attaining $q$. Maximize over the finite interface and components. This proves (1)--(2), including equality at the critical error.

We detail effectivity of the occupation proof, since computing the threshold alone would not supply it. Choose a component with radius greater than the given $\varepsilon$, and enumerate positive words until one lies in that component. Its rational product gives the prefix used in the proof of Theorem~\ref{thm:radius54}. The resulting map $f(h)=F_{s,j}(hu_v)$ is polynomial, so it can be used directly in place of $G$, with no approximation error.

For any finite tensor degree, all required polynomial functions are coefficients of the explicitly given tensor representation. Its fixed subspace is
\[
 \{z:\ \forall h\in K,
                 (R(h)-I)z=0\}.
\]
This linear subspace is effectively semialgebraic. Real algebraic sampling, successively outside the span of the vectors already chosen, gives an algebraic basis in finitely many steps. Its Gram matrix gives the orthogonal projection $P$ exactly. Thus every polynomial Haar moment is computable algebraically from $P\Phi(e)$; numerical integration of Haar measure is unnecessary.

Enumerate finite lists of squares of rational polynomials on $K$, discarding zero Haar integrals and normalizing the others. For a list $p_i$, compute $b_i=\int p_if\,d\mu$ and its constrained simplex radius by quantifier elimination. Rational polynomials are uniformly dense among real polynomial restrictions on a compact matrix group; the localization proof therefore guarantees that some enumerated list has radius strictly greater than $\varepsilon$. Strict algebraic comparison detects that list. Fix a positive rational $\Delta$ below the radius difference and a positive rational upper bound $L$ for all coefficient constants in Proposition~\ref{prop:residual54}. These give the certified residual $Q\ge\Delta/L$.

Next enumerate finite rational probability laws on executable words of a common length whose products lie in $K$, together with rational $d\in(0,1)$. The assertion \eqref{eq:gap54}, quantified over the finitely many unit vectors in the algebraic moving space, is semialgebraic: each maximum of finitely many inner products can be expressed by a finite disjunction. It is therefore decidable. Lemma~\ref{lem:gap54} and sufficiently small rational perturbations of its atom weights ensure that this search terminates. Padding by the identity letter makes every finite choice executable at one common length.

The resulting $B,d,L,\Delta,Q_0$ determine \eqref{eq:Lk54}. To avoid exact comparisons involving logarithms, take a rational $A'\ge\max(1,LQ_0/\Delta)$ and use
\[
 \frac{\log A}{-\log(1-d)}\le\frac{A'}d.
\]
An integer above $b+B-1+BA'/d$ is a computable occupation bound. This proves (3).

For (4), fix a candidate width $k$ and use $k$ labels at every cut, padding smaller registers with unused labels. Stochastic rows, initialization, and categorical decoder probabilities are finitely many variables in compact simplexes. For each of the finitely many length-$N$ words, the output vector is polynomial in these variables, and the target vector is rational. The total-variation inequalities are semialgebraic. Quantifier elimination therefore decides feasibility and supplies an algebraic point when feasible. A deterministic register storing the seed and entire prefix is a finite exact realization, so increasing $k$ must terminate. This proves the least-width claim for each finite horizon, without claiming a finite-horizon cutoff for all stationary or nonuniform realization questions.
\end{proof}
The distinction between an effective theorem and an implemented solver matters. The accompanying finite checks verify displayed exact examples and proof constants; they do not execute the general group-closure and quantifier-elimination construction. The finite-bit compiler for rational finite groups remains in the complete predecessor supplement with its separate hypotheses and costs. It is not used to reprice arbitrary algebraic decoders as cost-free finite-bit procedures.

\subsection{The least component-quotient realization}
The upper bound $|\mathcal S||C|$ can often be reduced. Here is an exact finite characterization of that reduction, with its scope stated explicitly. Put $X=\mathcal S\times C$. A \emph{component-quotient realization} has a deterministic state of the form $\pi(s,c)$, where $\pi:X\to Q$ is a surjection; its command update is induced by $(s,c)\mapsto(s,[u_a]c)$. It is stationary and has no additional state recording a representative word within a component.

Call a partition $\mathcal B$ of $X$ equivariant if
\[
 (s,c)\sim(s',c')\quad\Longrightarrow\quad
 (s,dc)\sim(s',dc')\qquad(d\in C).
\]
For a block $B\in\mathcal B$ and query $j$, set
\[
 A_{B,j}=\bigcup_{(s,c)\in B}F_{s,j}(c).
\]
\begin{theorem}[Finite quotient minimum]\label{thm:quotient54}
The least number of command labels in an error-$\varepsilon$ component-quotient realization is
\begin{equation}\label{eq:quotient54}
 \min\left\{|\mathcal B|:\mathcal B\text{ equivariant and }
         \rad_{Y_j}(A_{B,j})\le\varepsilon\text{ for every }B,j\right\},
\end{equation}
with value $+\infty$ if there is no admissible partition. For rational orthogonal categorical inputs this minimum and an attaining quotient are computable. If $\epsc=0$ and $\varepsilon=0$, the coarsest admissible partition is the future-output equivalence
\[
 (s,c)\sim(s',c')\quad\Longleftrightarrow\quad
 F_{s,j}(dc)=F_{s',j}(dc')\quad\text{for every }d\in C,j,
\]
where componentwise constant maps are identified with their values.
\end{theorem}
\begin{proof}
An equivariant partition induces well-defined command permutations on its blocks. Choosing a minimizing center for each $A_{B,j}$ gives exactly the stated error. Conversely, the fibers of $\pi$ in any component-quotient realization form an equivariant partition. Its decoder must approximate every executable product in each corresponding component. Such products are dense in that component, so continuity extends the error bound to $A_{B,j}$. Its decoder is therefore a permitted center. This proves \eqref{eq:quotient54}. The command generators suffice to check equivariance, since their images generate the finite group $C$.

There are finitely many partitions of $X$. For polynomial categorical data their radius tests are of the form \eqref{eq:QEcenter54}, with a finite disjunction over the relevant seed/component pairs, so the minimum is computable. In the exact componentwise constant case, equality of all future outputs is an equivariant equivalence relation. Every admissible partition refines it, and its own blocks have singleton output images for each query. It is therefore the unique coarsest admissible partition.
\end{proof}
The minimization in \eqref{eq:quotient54} is a deterministic quotient minimum, distinct from the unrestricted stationary minimum in Theorem~\ref{thm:stationaryminimum55} and from the unrestricted nonuniform minimum at the boundary. Randomized encodings can identify numerical behaviors without identifying component pairs. What is classified here is exactly the finite deterministic quotient construction that attains the radius threshold. Theorem~\ref{thm:effective54}(4) separately computes the unrestricted minimum at each fixed horizon.

\subsection{Stationary minimization by permutation closures}
\begin{theorem}[Effective unrestricted stationary minimum]\label{thm:effectivemin55}
For rational orthogonal commands generating $H$, rational polynomial categorical outputs, and algebraic $\varepsilon\ge0$, one can compute $M_\varepsilon$ and an attaining algebraic machine when it is finite. An identity letter is not required. At each proposed width $k$, at most $(k!)^{|\mathcal A|}$ augmented orthogonal group closures and finite real quantifier-elimination problems suffice.
\end{theorem}
\begin{proof}
First compute $H$, its real connected components, and their constrained radii as in Theorem~\ref{thm:effective54}(1)--(2); those steps do not use an identity letter. The group $H\subset O(D)$ has finitely many components. Theorem~\ref{thm:finitequotient55} shows that $M_\varepsilon=\infty$ below their maximum radius. Otherwise $M_\varepsilon\le |\mathcal S||H/H^\circ|$.

For $k$ between one and this upper bound enumerate all command-permutation tuples $\sigma$. The matrices
\[
                         \diag(u_a^{-1},\Pi_a)\in O(D+k,\Q)
\]
generate a compact closure which is exactly the joint group $\mathcal L_\sigma$ of Theorem~\ref{thm:stationaryminimum55}. Compute it by the same algebraic closure procedure. In \eqref{eq:permutationfeasibility55}, the unknown seed rows and decoder rows are constrained to simplexes. The mean output is bilinear in these rows, the group variable is restricted by computed defining equations whose real zero set inside the augmented orthogonal group is exactly $\mathcal L_\sigma$, not merely by some equations vanishing on it, and $h^{-1}=h^T$ in the physical block. Consequently all target coordinates are polynomial in the universally quantified matrix entries. Total variation is semialgebraic. Real quantifier elimination decides the resulting existential--universal formula and supplies algebraic rows whenever it is feasible. The first feasible $k$ is the minimum by purification. The finite upper bound proves termination of the entire search.
\end{proof}
This is a structural reduction of stationary minimization, not merely an enumeration of arbitrary stochastic transition parameters. It does not equate stationary and finite-horizon minima. Theorem~\ref{thm:stationarization56} identifies their limit under its additional return hypothesis. The two-cycle/three-cycle example in Proposition~\ref{prop:C655} already prevents replacing the permutation-feasibility test by deterministic component partitions.

\subsection{Algorithmic cost and arithmetic representation}
There are three distinct layers of effectivity. First, the group-closure procedure produces explicit defining equations; real-component decomposition and algebraic sampling then identify the component table. Second, radius computation and the stationary permutation tests are finite quantifier-elimination tasks on those equations. At fixed $k$, the stationary seed and output variables number $k(|\mathcal S|+\sum_jM_j)$ before normalization equalities are eliminated; the group variables can be taken among the $(D+k)^2$ augmented matrix entries. Thus the displayed finite search is independent of a horizon parameter. Its factorial number of discrete cases and the cost of the real-algebraic subproblems should not be mistaken for an efficient algorithm. No bound for the full closure-to-machine process in terms of only the original bit length is extracted here.

Third, occupation certificates in Theorem~\ref{thm:effective54}(3) add searches over localizer degrees and executable word laws. Their termination follows from strict analytic margins; the present proof supplies no a priori useful bound on the successful degree or word length. Fixed-horizon minimization, by contrast, has finitely many word constraints at each $N$ but their number is exponential in $N$ before repetitions are identified. These stages should be priced separately rather than hidden behind the word ``effective.''

All algebraic Haar moments in the localizer search are represented in a real algebraic number field with isolating data, using the fixed-subspace projector and its algebraic Gram inverse. Signs and strict comparisons are exact real-algebraic comparisons; a quadrature tolerance is not a certificate. Every rational command word has a rational matrix product and hence a rational target vector for rational polynomial readouts. The orthogonal closure fact used above is precisely that the Euclidean closure of a subgroup of $O(D)$ is the real zero set, inside $O(D)$, of its vanishing polynomial ideal, as used in \cite{BJKP}; replacing real connected components by complex algebraic ones would give the wrong quotient.

\section{Finite descriptions and decision complexity}\label{sec:complexity56}
The qualitative reduction is insensitive to the encoding of real numbers. Complexity statements are not. We first specify a finite input model and then give an explicit finite-precision counterpart to label width.

\begin{theorem}[Finite-group stationary minimization is $\exists\R$-complete]\label{thm:ETR56}
The input consists of a finite group $G$ given by its full multiplication table, finite seed and query sets, rational categorical vectors $F_{s,j}(g)$ for every $g\in G$, a nonnegative rational tolerance $\varepsilon$, and an integer $k\ge1$ in binary. Deciding whether $M_\varepsilon\le k$ is complete, under polynomial-time many-one reductions, for the existential theory of the reals. Hardness already holds for the trivial group, one query, and zero error. In particular the rational orthogonal polynomial-interface stationary problem is $\exists\R$-hard even with constant readouts and trivial dynamics.
\end{theorem}
\begin{proof}
Let $n=|G|$. If $k\ge n|\mathcal S|$, a deterministic seed--group register realizes every target exactly, so the answer is yes. Otherwise $k$ is bounded by a polynomial in the explicit input length, even though its encoding is binary.

For membership, Theorem~\ref{thm:equivariant56} permits permutation matrices $P_g$ forming a representation of $G$. Introduce real variables for their entries, constrain each entry $x$ by $x(x-1)=0$, impose every row and column sum equal to one, and impose
\[
                         P_1=I,
                 \qquad P_gP_h=P_{gh}\quad(g,h\in G).
\]
Introduce simplex variables $\alpha_s$ and categorical decoder rows $d_j(i)$. For every $s,j,g$ and output coordinate $c$, use a nonnegative auxiliary variable $z_{s,j,g,c}$ with
\[
 z_{s,j,g,c}\ge\pm\bigl((\alpha_sP_{g^{-1}}d_j)_c-F_{s,j}(g)_c\bigr),
 \qquad\sum_c z_{s,j,g,c}\le2\varepsilon.
\]
This is an existential real formula with rational coefficients, degree at most three, and polynomial size. Padding a representation by fixed unused labels handles an optimum strictly smaller than $k$. Theorem~\ref{thm:equivariant56} proves equivalence with the original unrestricted stationary problem. There is no group-closure computation in this explicit finite-group input model.

For hardness, let $A$ be a rational nonnegative matrix and let $k$ be a proposed nonnegative rank. Delete zero rows; the all-zero case is decided directly. Normalize each remaining row by its positive row sum to obtain a rational row-stochastic matrix $F$. This transformation is polynomial-time and preserves nonnegative rank. Interpret the rows as seeds and the columns as the categories of one query on the trivial group. A zero-error realization gives a stochastic factorization $F=\alpha d$ of inner dimension at most $k$, so the nonnegative rank is at most $k$.

Conversely, from any nonnegative factorization $F=UV$ delete a zero row of $V$ together with its column of $U$, and let $s_i=\sum_cV_{ic}>0$. Put $d_{ic}=V_{ic}/s_i$ and $\alpha_{bi}=U_{bi}s_i$. Then $d$ is row-stochastic and $\sum_i\alpha_{bi}=\sum_cF_{bc}=1$, so $\alpha$ is row-stochastic. Identity command rows yield an exact stationary realization. Thus $M_0$ equals the real nonnegative rank of $F$. Shitov's universality theorem, specifically its complexity consequence \cite[Theorem~2]{Shitov}, makes this decision problem $\exists\R$-complete. The last assertion follows by using the one-dimensional identity matrix as the only physical command and constant rational polynomial readouts.
\end{proof}
The upper bound is for an explicitly tabulated finite group. The hardness transfers to rational matrix input, but a polynomial-size existential description of an arbitrary compact matrix closure has not been supplied. Theorem~\ref{thm:effectivemin55} remains the separate termination result for that broader input class. The complexity obstruction above already occurs in randomized initialization and decoding; it is not caused solely by closure algorithms.

\subsection{A finite-table resource counterpart}
For a categorical interface, a dyadic table of precision $b$ has each probability in $2^{-b}\Z$. Besides peak label width $k$, charge its binary table length, all retained working registers, and the number of fair random bits used. A horizon-$N$ clocked implementation can store initialization, $N$ command tables, and terminal decoders. Its probability-table length is
\[
 O\left(b\left(|\mathcal S|k+N|\mathcal A|k^2+
                                k\sum_jM_j\right)\right).
\]
A uniform implementation is one finite algorithm that produces or accesses these tables from the encoded input, horizon, and accuracy. This requirement is additional to the original nonuniform model. An arbitrary real table, without an effective encoding, is not such an algorithm.

\begin{lemma}[Dyadic rounding of a probability row]\label{lem:dyadic56}
A probability row $p$ with $m$ entries has a dyadic row $\widehat p$ of precision $b$ such that
\[
                   \TV(p,\widehat p)\le(m-1)2^{-b}.
\]
One exact draw from $\widehat p$ uses at most $b$ independent fair bits.
\end{lemma}
\begin{proof}
Round the first $m-1$ entries down to multiples of $2^{-b}$ and put the remaining mass in the last entry. This is a probability row. Its total variation error equals the mass transferred to the last entry, which is at most $(m-1)2^{-b}$. For sampling, draw a uniform integer in $\{0,\ldots,2^b-1\}$ using $b$ bits and compare it with the integer cumulative row sums. Degenerate rows can use fewer bits.
\end{proof}

\begin{theorem}[Precision and random-bit budgets]\label{thm:bits56}
A categorical horizon-$N$ realization of peak width $k$ and error at most $\varepsilon$ admits a dyadic realization on the same registers whose error is at most
\begin{equation}\label{eq:clockround56}
       \varepsilon+\bigl((N+1)(k-1)+M_*-1\bigr)2^{-b},
                         \qquad M_* =\max_j M_j.
\end{equation}
Sampling initialization, all transitions, and one terminal answer requires at most $(N+2)b$ fair bits.

For a stationary compact-group realization of width $k$, the corresponding error bound can instead be made independent of $N$:
\begin{equation}\label{eq:stationaryround56}
                    \varepsilon+(k+M_*-2)2^{-b}.
\end{equation}
The stationary dyadic realization has at most $k$ labels, permutation command tables, table length
\[
 O\left(|\mathcal A|k\log(k+1)+
                   b\left(|\mathcal S|k+k\sum_jM_j\right)\right),
\]
and uses at most $2b$ fair bits in total for one experiment of any length. Given exact algebraic encodings of the selected initialization and decoder rows, the rounding and sampling are effective finite procedures.
\end{theorem}
\begin{proof}
Apply Lemma~\ref{lem:dyadic56} to every row. Couple the initial draws maximally and then, as long as the two labels agree, maximally couple their next transition draws. The probability of any disagreement is at most the sum of the initialization and $N$ row errors, hence at most $(N+1)(k-1)2^{-b}$. Conditioned on matching terminal labels, the terminal categorical rows have total variation error at most $(M_*-1)2^{-b}$. The coupling inequality gives \eqref{eq:clockround56}, uniformly over words and queries. The sample count is the number of draws times $b$.

For the stationary claim, first use Theorem~\ref{thm:purification55}. Its transition matrices are permutations and are already dyadic at every precision. Round only initialization and terminal categorical rows. The initialization discrepancy cannot increase under a permutation, so the same coupling gives \eqref{eq:stationaryround56} with no factor $N$. There are no random transition draws, leaving only initialization and one terminal answer. A permutation is stored as $k$ integer images, giving the displayed description bound. Exact comparison of algebraic numbers with dyadic rationals computes the necessary floors and handles equality, after which integer cumulative sampling is immediate.
\end{proof}
For a prescribed extra error $\eta>0$, choosing $b$ with the added term in \eqref{eq:clockround56} or \eqref{eq:stationaryround56} at most $\eta$ gives an explicit precision budget. The persistent label uses $\lceil\log_2k\rceil$ bits; an implementation of the sampler also retains a $b$-bit random integer and bounded indexing registers while sampling. These temporary registers are not disguised as free persistent randomness. The description and random-bit bounds do not assert a polynomial time bound for constructing an algebraic optimum. They also do not turn an exact irrational realization into an exact finite-bit one: the extra tolerance is essential in that case.

\section{Positive finite-precision matrix decoders}\label{sec:matrixprecision57}
The label-width bounds leave real table descriptions uncharged. There is
nevertheless a finite-description counterpart for each supplied encoded
realization. Entrywise rounding of a density matrix is unsafe: it need not
preserve positivity. A Gram-factor construction avoids this problem.

\begin{lemma}[Rational positive rounding]\label{lem:positiveprecision57}
Let $D\in\mathcal D_d$ and let $L=D^{1/2}$. Approximate the real and imaginary
parts of every entry of $L$ by dyadic multiples of $2^{-b}$, each within
$2^{-b}$, to obtain $Q$. Set $\delta=\sqrt2d\,2^{-b}$. If $\delta\le1/4$, then
\begin{equation}\label{eq:gramround57}
 \widehat D=\frac{QQ^*}{\tr(QQ^*)}\in\mathcal D_d\cap
            M_d(\Q+i\Q),\qquad
 \norm{\widehat D-D}_{\mathrm F}\le8\sqrt2d\,2^{-b}.
\end{equation}
The construction is effective when $D$ is given by exact algebraic entries.
Its rational output has $O(d^2(b+\log(d+1)))$ bits for fixed encoding
conventions.
\end{lemma}
\begin{proof}
Write $Q=L+E$. Then $\norm E_{\mathrm F}\le\delta$ and
$\norm L_{\mathrm F}^2=\tr D=1$. Hence $Q\ne0$,
$\tr(QQ^*)\ge(1-\delta)^2$, and
\[
 \norm{QQ^*-D}_{\mathrm F}\le(2+\delta)\delta,
 \qquad |\tr(QQ^*)-1|\le(2+\delta)\delta.
\]
The first estimate follows from the two cross terms and $EE^*$, and the
second from $\norm Q_{\mathrm F}^2-\norm L_{\mathrm F}^2$. Since
$\norm D_{\mathrm F}\le1$, division by $\tr(QQ^*)$ gives an error at most
$2(2+\delta)\delta/(1-\delta)^2\le8\delta$. Positivity and trace one are
exact, not numerical tolerances. After multiplying $Q$ by $2^b$, the
normalization is a ratio of integer Gram entries to a positive integer
sum of squares. Their bit sizes are $O(b+\log(d+1))$. The positive square
root of an algebraic positive matrix has algebraic entries and can be
specified and approximated by real-algebraic arithmetic. This proves the
effectivity assertion without an efficiency claim for obtaining $D$.
\end{proof}

\begin{theorem}[A supplied realization with finite tables and fair bits]
\label{thm:matrixbits57}
Suppose an error-$\varepsilon$ horizon-$N$ realization for a matrix
interface uses at most $k$ labels, and its initialization, command rows,
and decoders are supplied with exact algebraic encodings. Let
$\sqrt2d\,2^{-b}\le1/4$. There is a rational finite-table implementation,
with the same persistent label widths, whose Frobenius error is at most
\begin{equation}\label{eq:matrixbits57}
 \varepsilon+\sqrt2\bigl((N+1)(k-1)+8d\bigr)2^{-b}.
\end{equation}
Each hidden draw uses $b$ independent fair bits, so at most $(N+1)b$ such
bits suffice. For a supplied stationary permutation realization the
additional error is at most
\begin{equation}\label{eq:stationarymatrixbits57}
                  \sqrt2\bigl(k-1+8d\bigr)2^{-b},
\end{equation}
and only $b$ fair bits are needed for initialization, independently of the
horizon. The numerical matrix decoder itself is not a quantum preparation
or a sampled terminal answer register.
\end{theorem}
\begin{proof}
Apply Lemma~\ref{lem:dyadic56} to each initialization and command row after
padding to $k$ labels. Each row changes by at most $(k-1)2^{-b}$ in total
variation. Induction, or a maximal coupling at every actual draw, bounds
the terminal distribution discrepancy by $(N+1)(k-1)2^{-b}$; this upper
bound may harmlessly exceed one. Replace each decoder by the legal
rational matrix in Lemma~\ref{lem:positiveprecision57}. The Frobenius
diameter of $\mathcal D_d$ is $\sqrt2$, so changing the distribution incurs
at most $\sqrt2$ times its total-variation discrepancy. The independent
decoder error is at most $8\sqrt2d\,2^{-b}$. This proves
\eqref{eq:matrixbits57}. A dyadic probability row of denominator $2^b$ is
sampled by one uniform $b$-bit integer and its cumulative integer table.

For a permutation realization the command rows already have exact
integer entries and need not be rounded or sampled. Only the
initialization distribution and decoder are changed. The same estimate
therefore gives \eqref{eq:stationarymatrixbits57} and the horizon-independent
fair-bit bound.
\end{proof}
The row tables and their cumulative sums require
$O((N|\mathcal A|k^2+|\mathcal S|k)b)$ bits in the general case, in addition
to $O(kd^2(b+\log(d+1)))$ bits for one matrix decoder. Temporary random
integers and table indices require their own $O(b+\log(k+1))$ workspace;
this workspace is not hidden in the persistent-label invariant. Multiple
queries multiply the decoder storage by their number. To implement a
sampled measurement one must additionally specify and charge its sampler.
The theorem compiles a supplied encoded realization; it neither computes
a width minimizer nor assigns finite encodings to arbitrary unspecified
real tables. Applying it after the exact inner-polytope construction uses
an explicit extra error budget, not an exact finite-precision claim at an
extreme zero-error endpoint.

Here an exact algebraic entry is specified by defining polynomials and
real isolating intervals (for both parts of a complex entry), with
compatible field operations. The positive square root is the unique
positive semidefinite solution of $L^2=D$, and its isolating data can be
obtained by real-algebraic algorithms. These representation and output
bit-length statements are not running-time bounds for finding the
original realization or its square root. The constructed rational qubit
compiler is a separate, more explicit algorithm.

\section{Relation to numerical automata and positive realization}\label{sec:comparison54}
The input/output object and the order of quantifiers are decisive in comparing realization theorems. Classical probabilistic automata assign numerical acceptance probabilities to words, while stochastic-language results also study threshold recognition \cite{Rabin,Paz}. A numerical equivalence of two fixed machines is not the assertion that for every horizon there exists a new bounded-width approximate machine. Conversely, language equivalence need not preserve the numerical probabilities.

For weighted automata over a field, the finite-Hankel-rank characterization gives finite linear representations and identifies their minimum dimension \cite{BR}. It does not enforce nonnegative rows, stochastic normalization, or simplex-valued decoders. Our common-row contraction concerns these constraints simultaneously. The invariant observable quotient in Theorem~\ref{thm:observable53} uses the same reachable/invisible linear reduction; its additional conclusion is a finite-group criterion for exact positive realization with the full horizon-dependent quantifier.

The classical positive-realization problem asks for a positive system representing specified input/output data, and distinguishes positive dimension from unrestricted linear dimension \cite{BF04}. Invariant cones and polyhedral sections are part of that theory, not new notions introduced here. Proposition~\ref{prop:polygon54} is an explicit common invariant-enclosure construction. The new matched statement pairs it with a full-subcritical lower estimate against all clocked stochastic machines. Neither one cutwise nonnegative factorization nor a finite-rank Hankel matrix supplies that lower estimate.

Brodsky and Pippenger identify bounded-error measure-once quantum languages as group languages \cite[Theorem 3.3]{BP}. Their probabilistic simulation result preserves a cutpoint language with a potentially changed cutpoint \cite[Theorem 3.6]{BP}; it is not uniform additive preservation of every acceptance probability. Our categorical unitary corollary fixes an entire terminal measurement law in total variation and allows horizon-dependent stochastic approximants. The distinction is already visible in the $d$-outcome threshold $1-1/d$.

The closure theorem and strict-threshold decision method of Blondel, Jeandel, Koiran, and Portier \cite{BJKP}, and the effective Zariski-closure algorithms \cite{DJK,NPSW}, are algorithmic premises of Section~\ref{sec:effective54}. Computing the closure is not a new claim of this article. The additional outputs are a constrained minimax noise threshold, an attaining categorical decoder, and a terminating construction of an all-machine occupation constant. Recent Lie-algebra-based closure computations \cite{deGraaf} give an alternative representation of the algebraic closure, but their complex connected components still require real-component analysis for the present application.

Peter--Weyl approximation, Haar averaging, and the representative-function algebra are classical compact-group and almost-periodic-function machinery \cite{Folland}. Uniform approximation by a finite sum of matrix coefficients gives a finite \emph{linear} representation of an output function. It does not give positive stochastic rows or the repeated information loss quantified by \eqref{eq:occupation54}. Similarly, Steinberg's minimal compact topological automaton and enveloping-monoid results \cite[Theorems 2.2 and 2.6]{Steinberg} concern recognition of a language. His compact topological-monoid criterion \cite[Proposition 2.8]{Steinberg} uses clopen recognition, not a continuous simplex-valued response with a prescribed error norm. The finite-quotient radius formula is specific to groups. The stationarization theorem itself now holds for irreversible topological monoids with eventual approximate returns (Theorem~\ref{thm:stationarization56}).

Distribution-based bisimulation metrics provide robust comparisons for specified probabilistic automata \cite{FSZ}. They do not by themselves impose a minimum over all horizon-specific stochastic realizations of an external group experiment. The finite future-output equivalence in Theorem~\ref{thm:quotient54} is deliberately a component-quotient statement, rather than a claim that approximate equivalence always has a unique minimal positive stochastic realization.

\begin{center}\small
\begin{tabular}{p{.27\textwidth}p{.31\textwidth}p{.32\textwidth}}
\toprule
Comparison object & Error or equivalence & Machine quantifier\\
\midrule
Compact stochastic groups \cite{Flor} & Idempotent and permutation structure; no external error specification & One matrix group or semigroup kernel\\
Advice regular languages \cite{FP} & Neutral-letter recognition & One-scan advice automata; Boolean output\\
Weighted series \cite{BR} & Exact scalar coefficients; linear rank & One linear representation\\
Group-language simulation \cite{BP} & Cutpoint language; not additive probability error & One automaton\\
Positive realization \cite{BF04} & Exact specified numerical response & One positive system\\
Topological recognition \cite{Steinberg} & Clopen language acceptance & One recognizing action\\
Bisimulation metrics \cite{FSZ} & Distributional comparison of given systems & Two specified automata\\
Theorem~\ref{thm:stationarization56} & Same closed numerical error and label bound & Arbitrarily many narrow cuts; one extracted stationary machine\\
Theorem~\ref{thm:radius54} & Worst-word norm error; TV for a whole categorical law & A new clocked stochastic machine at every horizon\\
Theorem~\ref{thm:matched54} & Binary TV, every fixed $\varepsilon<\rho/2$ & Same nonuniform lower model; exact upper construction\\
\bottomrule
\end{tabular}
\end{center}
The contribution combines stationary purification and its exact finite-quotient criterion with a sharp constrained clocked radius, phasewise arbitrary-width occupation, and, for the specified Diophantine family, a matching full-subcritical growth law. The elementary radius center, the arithmetic pigeonhole argument, the closure algorithm, and the harmonic-analysis ingredients are not asserted individually new. The accompanying theorem-level audit records the closest comparisons and their changed hypotheses; it is not a claim of exhaustive priority clearance.


\subsection{The classical matrix-semigroup boundary}
Flor's description of nonnegative matrix groups \cite[Theorems 2 and 3]{Flor} identifies their idempotent geometry and the finiteness of bounded groups; his account explicitly credits Brown and Schwarz for earlier compact and stochastic cases. We do not claim this matrix structure as new. Theorem~\ref{thm:purification55} uses a minimal-rank idempotent in a joint semigroup, not just in the matrix projection: its physical coordinate must be the group identity. Sandwiching each command then preserves all uniform error inequalities by closure, even when the original stochastic rows violate every nontrivial physical group relation. The resulting initialization may be random, a point essential to exact state minimization.

The stationary theorem has no return hypothesis. The clocked theorem has a different quantifier and needs synchronization; the empty-return example and the phase-dependent example show why the distinction cannot be removed by terminology. The exact-return-length reduction is an additive-semigroup argument, whereas the profinite zero-error converse uses finite Hankel rank and continuous finite-dimensional representations. The new stationarization theorem removes the reversibility requirement for the equality of limiting width and stationary width, but does not supply a finite-quotient radius formula for an arbitrary irreversible monoid.

The finite-quotient existence theorem does not say that $[H:U]$ is the optimal number of stochastic labels. The actual minimum is the permutation-feasibility invariant in Theorem~\ref{thm:stationaryminimum55}; the finite cyclic example separates them. For infinitely many components, the infimum of finite-quotient radii can fail to be attained, even at zero. Finally, the exactly matched quantitative law retains its specified Diophantine alphabet. Theorem~\ref{thm:spherical56} gives the retained block estimate, strengthened to a matching power law by Theorem~\ref{thm:sharpoccupation61}, for spherical representations with a spectral gap; Corollary~\ref{cor:SO356} gives a concrete rational noncommuting instance. The exact unit-amplitude endpoint is not covered by these upper constructions. These distinctions delimit the proved statements without weakening any of the retained conclusions.

\subsection{What the reductions add}
The algebraic ingredients of Lemmas~\ref{lem:idempotent55} and \ref{lem:simplex55} are the classical minimum-rank kernel mechanism, the group corner of a compact matrix semigroup, and the recurrent simplex of a stochastic idempotent. Flor's Theorem~2 describes nonnegative idempotents; his Theorem~3 and preceding Lemma~3 give the permutation structure of bounded nonnegative groups \cite{Flor}. The paper explicitly attributes earlier cases to Brown and Schwarz. The present proof adds the physical coordinate to the closure, chooses the idempotent over its identity, and tracks the closed error balls through sandwiching. Theorem~\ref{thm:equivariant56} further averages the finite kernel above the physical identity and descends to an actual continuous $H$-action. These are specified realization reductions using classical structure, not a new classification of matrix-semigroup kernels.

For deterministic Boolean recognition, neutral-letter removal of advice is an established phenomenon. Fijalkow--Paperman \cite[Theorem~10]{FP} prove the Crane Beach property for advice regular languages by padding distinguishing words and suffixes, and their account relates one-scan programs to earlier Ramsey arguments. Their Myhill--Nerode argument in fact preserves the deterministic distinguishability bound $K$; preservation of a state bound alone is not claimed as new here. A stochastic $k$-label register has a continuum of state distributions, so this discrete class count does not apply. Theorem~\ref{thm:stationarization56} instead colors approximate tuples of stochastic composite kernels, extracts one common tuple, and absorbs a final filler matrix into the decoder. It preserves the same numerical tolerance at equality and the same $k$, and bounds all narrow-cut occurrences with unrestricted intermediate widths. Eventual approximate returns further replace a literal neutral letter. The finite Ramsey theorem itself is classical \cite{Ramsey}.

The spherical lower estimate uses a classical spectral gap, cap measure, and smoothing; its upper estimate uses a polytope enclosing a contracted sphere. The contribution is their realization-theoretic combination with the all-profile conditional-centroid bound. The rational $\mathrm{SO}(3)$ example invokes \cite{BG2012} as a deep external input. Similarly, the lower-complexity input in Theorem~\ref{thm:ETR56} is Shitov's nonnegative-factorization universality theorem \cite{Shitov}; the stochastic normalization and finite-group existential upper reduction are given explicitly. Neither external theorem is claimed as part of our proof technology's novelty.

\subsection{Companion results}
The companion supplement gives additional results on compatible reachable sections, one-surplus geometry, stable sections, finite-horizon certificates, word distortion, two-state duality, finite-bit compilation, and global residual hierarchies. The scalar two-extreme argument and earlier quantitative constants are also retained there. The present main proofs are self-contained apart from the classical results explicitly cited; none of these companion derivations is used to conceal a missing step in a principal theorem.

\subsection{Orbit geometry, positive matrix outputs, and recognition}
The homogeneous-space covering estimates and the smooth orbit geometry used in Section~\ref{sec:uniform57} are classical; see, for example, Szarek~\cite{Szarek57}. The spectral-gap input is Bourgain--Gamburd~\cite{BG2012}. Dense algebraic free generators in Section~\ref{sec:extreme57} come from Breuillard--Gelander~\cite[Theorem~1.1]{BreuillardGelander57}; the later repair and supplementary proof details are recorded in~\cite{BreuillardGelanderCorrection57}. None of these results is claimed as a new group-theoretic theorem here. The realization argument requires, in addition, a contraction uniform over every conditional-centroid direction, one common positive row for each command and label, and a terminal decoder in the original legal matrix set.

Quantum finite automata already exhibit strong state-complexity separations in language recognition; Ambainis--Freivalds~\cite{AF57} is a relevant early comparison. Theorem~\ref{thm:freeendpoint57} instead asks a classical stochastic machine to reproduce one entire legal density-matrix vector at every input word, permits horizon-specific redesign and a free clock, and compares exact and fixed-error width for the same interface. Its exact lower bound uses extreme points of the decoder set; it does not apply unchanged to a single acceptance probability or to an enlarged unconstrained output simplex. A finite tomography map preserves width only when its legal image and reconstructed norm are retained. This comparison identifies different quantifiers and resources, not an exhaustive priority claim about all quantum or probabilistic automata formalisms.

The spectral-entropy argument removes the logarithm from the nonabelian lower bound. Equality of lower and upper powers requires equality of least-orbit and target-orbit dimensions; no universally matching law for every representation is asserted. General online outputs, adaptive feedback, and unrestricted compact semigroup radius formulas require additional arguments. The same-width stationarization result for topological monoids remains available under its actual approximate-return hypothesis; it does not provide that missing radius formula.

\section{The least orbit dimension and uniform small balls}\label{sec:minorbits58}
The geometric hypothesis needed for a stochastic converse concerns every
conditional-centroid direction, not only the prescribed target orbit.
For a fixed finite-dimensional orthogonal representation of a compact
connected Lie group, this hypothesis follows from a rank invariant.
Throughout this section Haar measure is normalized, and the Lie algebra
has a fixed invariant inner product. No constants uniform in the dimension
or in the representation are asserted.

\begin{theorem}[Uniform caps from the least orbit dimension]
\label{thm:minorbit58}
Let $G$ be a compact connected Lie group and $R:G\to O(E)$ a continuous
finite-dimensional representation with $E^G=\{0\}$ and $E\ne\{0\}$.
Write $D=\dim G$ and choose a basis $X_1,\ldots,X_D$ of its Lie algebra,
using the same symbols for their infinitesimal actions on $E$. Set
\begin{equation}\label{eq:minrank58}
 p_* =\min_{\norm u=1}\operatorname{rank}
       [X_1u\ \cdots\ X_Du]
     =\min_{\norm u=1}\dim R(G)u.
\end{equation}
Then $p_*\ge1$, and there is $C<\infty$ such that
\begin{equation}\label{eq:mincap58}
 \sup_{\norm u=1,\ v\in E}
 \mu\{g:\norm{R(g)u-v}<h\}\le C h^{p_*}
                 \qquad(0<h\le1).
\end{equation}
The exponent is optimal as a uniform exponent: it cannot be replaced by
any larger exponent with a constant independent of $u,v,h$.
\end{theorem}
\begin{proof}
The tangent space to the orbit at $u$ is the span of the $X_i u$.
An orbit of dimension zero is a singleton, since $G$ is connected; its
point would be fixed. This proves $p_*\ge1$. Put $r=p_*$ and fix
orthonormal coordinates on $E$.

For each unit $u$, some $r$-row, $r$-column minor of
$[X_1u\ \cdots\ X_Du]$ is nonzero. The maximum of the absolute values
of these finitely many minors is a continuous positive function on the
unit sphere. Its minimum is therefore positive. The entries are uniformly
bounded. Consequently there is $\sigma>0$ such that for every unit $u$
one can choose $r$ input indices and $r$ output coordinates whose derivative
matrix $A_u$ has least singular value at least $\sigma$.

For each of the finitely many permutations of the chosen Lie algebra
basis consider the product chart
\[
 \phi(x,y)=\exp(x_1X_{i_1})\cdots\exp(x_rX_{i_r})
            \exp(y_1X_{i_{r+1}})\cdots\exp(y_{D-r}X_{i_D}).
\]
All these charts are diffeomorphisms on one sufficiently small cube about
zero. Their images contain a common identity neighborhood $W$, and their
Haar densities are bounded above by a common constant. By shrinking the
cube once more, uniform continuity of the derivatives on the unit sphere
and the finite list of charts ensures the following. For every unit $u$,
choose the minor above and let $\pi_u:E\to\R^r$ be its coordinate
projection. For all $(x,y)$ in the cube,
\begin{equation}\label{eq:submersion58}
 \left\|\partial_x\bigl(\pi_u R(\phi(x,y))u\bigr)-A_u\right\|
                                                    \le\sigma/2.
\end{equation}
For fixed $y$, the map $f_{u,y}(x)=\pi_u R(\phi(x,y))u$ is injective.
Indeed integration of the derivative on the segment between $x$ and
$x'$ gives
\[
 \norm{f_{u,y}(x)-f_{u,y}(x')}
       \ge\norm{A_u(x-x')}-\tfrac12\sigma\norm{x-x'}
       \ge\tfrac12\sigma\norm{x-x'}.
\]
Every singular value of its derivative is at least $\sigma/2$.
The change-of-variables formula therefore bounds the $r$-dimensional
volume of the inverse image of a ball of radius $h$ by
$\omega_r(2/\sigma)^r h^r$. A ball in $E$ projects into a ball of the
same radius in these coordinates. Integrating in $y$ and using the bounded
Haar density proves the required estimate on this chart, uniformly in
$u$ and in the center $v$.

Choose finitely many right translates $Wg_j$ covering $G$. On such a
translate apply the preceding argument to the unit vector $R(g_j)u$.
For its selected minor, use the corresponding permuted chart; that chart
contains $W$. Haar invariance and a finite sum give \eqref{eq:mincap58}
for sufficiently small $h$. Enlarging $C$ covers the remaining
$h\le1$. This proof does not require a spectral decomposition of the
arbitrary center $v$.

Finally choose $u_*$ whose orbit has dimension $r$. The orbit is the
smooth embedded compact homogeneous space $G/G_{u_*}$. The pushforward
of Haar measure is its invariant probability measure, with positive
smooth density in local coordinates. At $u_*$, sufficiently small
ambient chordal balls have measure bounded below by $c h^r$ for some
$c>0$. Thus no larger uniform exponent is possible.
\end{proof}

\begin{proposition}[An algebraic rank certificate]\label{prop:rankcompute58}
Suppose the infinitesimal matrices $X_i$ are supplied with exact real
algebraic entries, and are known to specify the representation in
Theorem~\ref{thm:minorbit58}. Then $p_*$ is computable by real quantifier
elimination. A positive algebraic lower bound for the largest
$p_*\times p_*$ minor on the unit sphere is computable as well. If a unit
seed $z_0$ is algebraic, its orbit dimension $q$ is the rank of
$[X_1z_0\ \cdots\ X_Dz_0]$.
\end{proposition}
\begin{proof}
For each integer $r$ ask whether a unit vector exists on which all
$(r+1)\times(r+1)$ minors vanish. These are finitely many polynomial
conditions over real algebraic numbers; the first affirmative answer is
$p_*$. With $r=p_*$, the sum of squares of all $r$-minors has a positive
minimum on the unit sphere, by the first paragraph of the preceding
proof. Quantifier elimination computes that minimum as a real algebraic
number with isolating data. If there are $M$ minors, their largest
absolute value is at least $\sqrt{\gamma/M}$, where $\gamma$ is that minimum. Algebraic rank at the supplied seed gives the last assertion.
\end{proof}
The proposition computes the rank invariant from \emph{supplied}
infinitesimal matrices. It neither certifies that arbitrary input matrices
integrate to a compact group nor computes a random-walk spectral gap.

\begin{corollary}[Uniform width exponents without a separate cap hypothesis]
\label{cor:allrepresentations58}
In the orbit realization model of Theorems~\ref{thm:uniformorbit57} and~\ref{thm:sharpoccupation61}, assume
the declared spectral gap, but replace its small-ball hypothesis by
$E^G=\{0\}$. Let $p_*$ be \eqref{eq:minrank58}, and let $q$ be the real
dimension of the unit target orbit spanning $E$. For fixed
$0<\rho\le1$ and $0\le\varepsilon<\rho$,
\[
 W_{N,\varepsilon}\ge cN^{p_*/2}.
\]
For $0<\rho<1$ the exact upper bound is $C(N+1)^{q/2}$; for $\rho=1$
it holds at each fixed positive error. Every fixed width has the
occupation bound of Theorems~\ref{thm:uniformorbit57} and~\ref{thm:sharpoccupation61}. Error is measured in
the Euclidean norm defining the legal orbit hull.
\end{corollary}
\begin{proof}
Apply Theorem~\ref{thm:minorbit58} to obtain its uniform cap estimate with
exponent $p_*$. The lower and upper arguments of
Theorems~\ref{thm:uniformorbit57} and~\ref{thm:sharpoccupation61} then apply without alteration.
\end{proof}
The group dimension $D$, least orbit dimension $p_*$, and target orbit
dimension $q$ need not coincide. For the standard sphere representation
$p_*=q=d-1$. For conjugation on traceless Hermitian $d\times d$ matrices,
an orbit with eigenvalue multiplicities $m_1,\ldots,m_b$ has dimension
$d^2-\sum_i m_i^2$. A nonzero traceless matrix has $b\ge2$, and the
largest possible sum is $(d-1)^2+1$. Thus $p_*=2(d-1)$, attained by a
centered rank-one projection. The explicit Grassmannian estimates below
also provide this exponent; they are retained as a direct geometric
proof with stated dimension-dependent constants.

The converse now uses Theorem~\ref{thm:sharpoccupation61}, rather than
paying the pointwise-mixing block length at each narrow cut. Thus no
logarithm separates the bounds when $p_*=q$. A supplied algebraic
infinitesimal representation makes $p_*$ computable; this statement
does not compute a spectral gap or assert dimension-uniform constants.

\section{A fixed rational experiment with an exponential exact profile}
\label{sec:rationalendpoint58}
We now give fixed, explicit matrices and a rational pure seed. The
freeness and trivial-stabilizer assertions have elementary finite-field
proofs; no generic choice of generators or transcendental seed is used.
The spectral-gap input is needed only for the noisy lower bound.

\begin{lemma}[A quaternion free triple]\label{lem:freequaternion58}
In the unit quaternions put
\[
 a=(1+2\mathbf i)/\sqrt5,\qquad
 b=(1+2\mathbf j)/\sqrt5,\qquad
 c=(1+2\mathbf k)/\sqrt5.
\]
Every nonempty freely reduced word in $a^{\pm1},b^{\pm1},c^{\pm1}$ is
nonscalar. In particular these elements freely generate a group of rank
three, also after passage to the projective unitary group.
\end{lemma}
\begin{proof}
Use the Hamilton relations $\mathbf i^2=\mathbf j^2=\mathbf k^2=-1$,
$\mathbf i\mathbf j=\mathbf k=-\mathbf j\mathbf i$. Modulo five these
relations are respected by
\[
 \mathbf i\mapsto\begin{pmatrix}2&0\\0&3\end{pmatrix},\quad
 \mathbf j\mapsto\begin{pmatrix}0&1\\4&0\end{pmatrix},\quad
 \mathbf k\mapsto\begin{pmatrix}0&2\\2&0\end{pmatrix}.
\]
The six numerators $1\pm2\mathbf i,1\pm2\mathbf j,1\pm2\mathbf k$
therefore have respective images
\begin{equation}\label{eq:sixmatrices58}
 \begin{gathered}
 \begin{pmatrix}0&0\\0&2\end{pmatrix},\quad
 \begin{pmatrix}2&0\\0&0\end{pmatrix},\quad
 \begin{pmatrix}1&2\\3&1\end{pmatrix},\\
 \begin{pmatrix}1&3\\2&1\end{pmatrix},\quad
 \begin{pmatrix}1&4\\4&1\end{pmatrix},\quad
 \begin{pmatrix}1&1\\1&1\end{pmatrix}.
 \end{gathered}
\end{equation}
They have rank one and their image lines are the six distinct points of
$\mathbb P^1(\mathbb F_5)$. The kernel line of each matrix is the image
line of the matrix with the opposite sign. Hence two consecutive
matrices have zero product exactly when their letters are inverses.
Writing each rank-one matrix as a column times a row shows that a product
of them is nonzero whenever no adjacent pair is inverse.

A reduced word of length $n\ge1$ is $5^{-n/2}Q$, where $Q$ is an integral
Hamilton quaternion of norm $5^n$. Its image modulo five is nonzero by
the preceding observation. Were $Q$ scalar, its integer scalar entry $t$
would satisfy $t^2=5^n$. For odd $n$ this is impossible; for even $n\ge2$
it forces $t=\pm5^{n/2}$, making the image zero modulo five. This
contradiction proves the assertion. Since the scalar unit quaternions
are $\pm1$, it proves both freeness statements.
\end{proof}

Use the faithful complex representation
\[
 \rho(t+x\mathbf i+y\mathbf j+z\mathbf k)=
 \begin{pmatrix}t+xi&y+zi\\-y+zi&t-xi\end{pmatrix}.
\]
Define the following matrices, all of whose real and imaginary parts are
rational:
\begin{equation}\label{eq:rationaldata58}
 U=\frac15\begin{pmatrix}-3+4i&0\\0&-3-4i\end{pmatrix},\quad
 V=\frac15\begin{pmatrix}-3&4\\-4&-3\end{pmatrix},\quad
 P_0=\frac12\begin{pmatrix}1&1\\1&1\end{pmatrix}.
\end{equation}
Here $U=\rho(a^2)$, $V=\rho(b^2)$, and $P_0$ is a rank-one density
matrix. The alphabet is exactly
\begin{equation}\label{eq:rationalalphabet58}
                  \mathcal A=\{I,U,U^*,V,V^*\}.
\end{equation}

\begin{theorem}[Explicit rational exact--noisy separation]
\label{thm:rationalseparation58}
For the fixed input \eqref{eq:rationaldata58}--\eqref{eq:rationalalphabet58},
let the terminal target be $F(g)=gP_0g^*$ and require each decoder value
to be a legal density matrix in $\mathcal D_2$. Error is Frobenius error.
At every horizon $N\ge0$ the exact minimum command-width profile is
\begin{equation}\label{eq:rationalprofile58}
             |S_t|_{\min}=2\cdot3^t-1\quad(0\le t\le N),
       \qquad W_{N,0}=2\cdot3^N-1.
\end{equation}
In fact the same profile is necessary, and is attained by the exact
machine, throughout the closed range
\begin{equation}\label{eq:robustendpoint58}
                      0\le\varepsilon\le625^{-N}/16.
\end{equation}
For each fixed $0<\varepsilon<1/\sqrt2$, constants $c_\varepsilon,
C_\varepsilon>0$ satisfy
\begin{equation}\label{eq:rationalnoisy58}
 c_\varepsilon N
       \le W_{N,\varepsilon}\le C_\varepsilon(N+1).
\end{equation}
At $\varepsilon\ge1/\sqrt2$ one command label suffices. The upper bound
in \eqref{eq:rationalnoisy58} has a uniform rational-table and fair-bit
construction, specified in Theorem~\ref{thm:rationalcompiler58}.
\end{theorem}
\begin{proof}
First, $\langle U,V\rangle$ is free of rank two: a nontrivial reduced
word in $a^{\pm2},b^{\pm2}$ expands to a nontrivial reduced word in
$a^{\pm1},b^{\pm1}$, and Lemma~\ref{lem:freequaternion58} applies.
If a member $g$ of this subgroup fixes $P_0$, then it commutes with $P_0$.
It consequently commutes with
\[
             \rho(c)=\bigl(I+2i(2P_0-I)\bigr)/\sqrt5.
\]
A nonempty reduced word $w$ in $a,b$ cannot commute with $c$:
$wcw^{-1}$ is reduced, has length $2|w|+1$, and cannot equal $c$ in the
free group. Thus the projective stabilizer of $P_0$ in $\langle U,V\rangle$
is trivial.

The subgroup is dense in $SU(2)$. Both $U$ and $V$ have trace $-6/5$.
If either had root-of-unity eigenvalues, its trace would be a rational
algebraic integer and hence an integer, a contradiction. Its powers are
therefore dense in its one-dimensional torus. The Lie algebra of the
closed generated subgroup contains the $\mathbf i$ and $\mathbf j$
directions and their bracket, the $\mathbf k$ direction. It is all of
$\mathfrak{su}(2)$; connectedness of $SU(2)$ implies density.

By freeness and the trivial stabilizer, the reachable pure states are in
bijection with the reduced word ball of radius $t$. Identity padding
makes this precisely the set reachable at length $t$. Its size is
$1+4\sum_{j=0}^{t-1}3^j=2\cdot3^t-1$. The extreme-orbit profile theorem
(Theorem~\ref{thm:extreme57}) gives \eqref{eq:rationalprofile58}. Its proof
uses one fixed suffix at each cut, not a query-dependent suffix.

We prove the stronger assertion \eqref{eq:robustendpoint58} directly.
Write $D=(I+x\cdot\sigma)/2$, with $x\in B_3$, for the Bloch coordinates
of a density matrix. Then
\begin{equation}\label{eq:blochnorm58}
                  \norm{D-D'}_{\mathrm F}=\norm{x-x'}/\sqrt2.
\end{equation}
The rotations induced by $U,V,U^*,V^*$ have entries in
$25^{-1}\Z$. At cut $t$, different reachable unit Bloch vectors have
coordinates in $25^{-t}\Z$, and hence mutual Euclidean distance at least
$25^{-t}$. Fix any suffix of length $N-t$. Its action is an isometry, so
the same separation holds for the corresponding terminal targets.
For each cut label let $z_i\in B_3$ be its conditional terminal decoder
mean after this fixed suffix. For a prefix with pure target $n$, error
at most $\varepsilon$ gives
\[
 \left\|\sum_i p_i z_i-n\right\|\le\sqrt2\varepsilon,
 \qquad
 \sum_i p_i\norm{z_i-n}^2
       \le2\left(1-\left\langle\sum_i p_i z_i,n\right\rangle\right)
       \le2\sqrt2\varepsilon.
\]
Some positive-mass label is therefore within
$\sqrt{2\sqrt2\varepsilon}$ of $n$. Under \eqref{eq:robustendpoint58},
twice this radius is at most
$\sqrt{\sqrt2/2}\,25^{-N}<25^{-t}$. These closed balls are disjoint
for different prefixes. They require distinct labels and prove the
pointwise lower bound at every cut. Deterministic exact storage gives
the simultaneous upper bounds.

For the positive-error converse, apply the spectral-gap theorem of
Bourgain--Gamburd \cite[Theorem~1]{BG2012} to the \emph{displayed} algebraic
generators $U,V$. The density just proved verifies its hypothesis.
The symmetric walk becomes an absolute $L^2$ contraction off the
constants after making it lazy, for example with mass $1/2$ at $I$ and
$1/8$ at each other letter. The legal density-matrix width theorem in
dimension two gives the lower bound in \eqref{eq:rationalnoisy58}.
The rational construction proved next gives its upper bound without
using a spectral gap. The center $I/2$ is at distance $1/\sqrt2$ from
every pure state, which gives the final threshold assertion; the lower
bound excludes bounded width at any smaller fixed error.
\end{proof}
This is a finite-input separation, not merely an existence statement in
an arbitrary-real advice model. Every length-$N$ target has real and
imaginary entries with denominator dividing $2\cdot25^N$, hence bit
length $O(N+1)$. Neither the exponential exact lower bound nor its robust
extension uses a spectral-gap theorem. The spectral-gap constant in the
noisy converse remains non-effective in the present proof; no numerical
value is inferred from finite word tests. The logarithmic interval in
\eqref{eq:rationalnoisy58} is stated explicitly rather than suppressed.

\section{A uniform rational qubit compiler}\label{sec:rationalcompiler58}
The preceding example also admits a constructive upper realization. The
compiler below constructs its rows rather than assuming that an optimal
realization has already been supplied. Its persistent labels and finite
table resources are accounted for separately.

\begin{lemma}[A rational spherical inner hull]\label{lem:rationalnet58}
For $m\ge2$ an integer, let $\mathcal V_m$ consist of the distinct points
\begin{equation}\label{eq:rationalnet58}
 v_{j,\ell,\eta}
 =\eta\frac{(2jm,\,2\ell m,\,m^2-j^2-\ell^2)}{m^2+j^2+\ell^2},
 \quad -m\le j,\ell\le m,\quad \eta\in\{-1,1\}.
\end{equation}
Put $r=1-m^{-2}$ and add the vertex $0$ as a distinguished label. Then
\[
 rB_3\subseteq\operatorname{conv}(\mathcal V_m)\subseteq B_3,
 \qquad K:=1+|\mathcal V_m|\le1+2(2m+1)^2\le13m^2.
\]
The points have $O(\log(m+1))$-bit rational coordinates. For any supplied
rational rotation $A\in SO(3)$ and any label $v$, the point $rAv$ has a
rational barycentric expression using $0$ and at most three members of
$\mathcal V_m$. Such an expression can be found by finitely enumerating
triples and solving rational $3\times3$ systems.
\end{lemma}
\begin{proof}
The stereographic parametrization
$\psi(s,t)=(2s,2t,1-s^2-t^2)/(1+s^2+t^2)$ has differential norm at
most two. Every point in the closed northern hemisphere has a preimage
in the unit disk, hence in $[-1,1]^2$. Rounding its coordinates to the
nearest multiples of $1/m$ gives a point of \eqref{eq:rationalnet58}
within distance $\sqrt2/m$. Negation treats the southern hemisphere.
For every unit $n$ there is thus $v\in\mathcal V_m$ with
$\langle n,v\rangle=1-\norm{n-v}^2/2\ge1-m^{-2}$.
The support-function criterion proves the inner inclusion. The stated
cardinality bound follows since $m\ge2$; duplicates are discarded.

The polytope contains a neighborhood of zero. The ray through a nonzero
$rAv$ reaches its boundary in a face. Triangulate a facet containing
that face using its vertices. A triangle containing the ray endpoint,
together with zero, gives the desired expression, allowing zero
coefficients on lower-dimensional faces. Its three nonzero vertices
are linearly independent: their affine plane is a supporting plane not
through zero. Enumerate all such triples, solve for their three
coefficients, and retain any solution with nonnegative coefficients and
sum at most one. The remaining mass is assigned to zero. Cramer's rule
gives rational coefficients; for $v=0$ use the point mass at zero.
All arithmetic is finite, and no selection theorem is needed.
\end{proof}

\begin{theorem}[Constructed finite-input upper realizations]
\label{thm:rationalcompiler58}
Let the alphabet consist of a fixed finite list of rational Bloch
rotations, and let the seed be the first coordinate unit vector. These
include the explicit experiment \eqref{eq:rationaldata58}. For a rational
$0<\beta<1$ and horizon $N\ge0$, choose
\begin{equation}\label{eq:compilerparameters58}
 m=\max\left(2,\left\lceil\sqrt{\frac{N}{1-\beta}}\right\rceil\right),
                 \qquad r=1-m^{-2}.
\end{equation}
There is an exactly correct rational stochastic realization of the
matrix target
$F_\beta(g)=I/2+\beta(gP_0g^*-I/2)$ on at most $13m^2$ command labels.
Each command row has at most four nonzero entries and is independent
of the command epoch within the chosen horizon.

For any rational $0<\varepsilon<1/\sqrt2$, the pure target $F_1$ admits
an error-$\varepsilon$ realization on $O(1+N/\varepsilon)$ labels with
rational legal matrix decoders and dyadic rows. It can be sampled with
at most $Nb$ fair bits, where it suffices to take
\begin{equation}\label{eq:compilerbits58}
 b=\left\lceil\log_2\frac{24(N+1)}{\varepsilon}\right\rceil.
\end{equation}
Both constructions are uniform algorithms polynomial in $N$,
$(1-\beta)^{-1}$ or $\varepsilon^{-1}$, and the input bit lengths. This
is not a polynomial bound in the binary length of $N$ or in the logarithm of the inverse accuracy.
\end{theorem}
\begin{proof}
The net includes the first coordinate unit vector, so initialization is
deterministic. For each letter $A$ and each net label $v_i$, select the
rational row in Lemma~\ref{lem:rationalnet58} whose barycentric mean is
$rAv_i$. The zero label stays at zero. Induction shows that the mean
Bloch vector of the label at time $t$ is $r^t A_{a_t}\cdots A_{a_1}e_1$.
At label $v_i$ use the matrix decoder
\begin{equation}\label{eq:rationaldecoder58}
                   D_i=\frac12\left(I+\frac{\beta}{r^N}
                                                   v_i\cdot\sigma\right).
\end{equation}
Bernoulli's inequality and \eqref{eq:compilerparameters58} give
$r^N\ge1-N/m^2\ge\beta$. Thus every decoder is positive semidefinite
of trace one. Its expectation is exactly $F_\beta$, for every length-$N$
word; the same rows work for all those words. Formula
\eqref{eq:rationaldecoder58} also gives $I/2$ at the zero label.

For the pure target choose $\eta=\min(\varepsilon,1/2)$ and
$\beta=1-\eta$. The amplitude loss contributes Frobenius error
$\eta/\sqrt2\le\varepsilon/\sqrt2$. Keep the rational decoders and
round each sparse probability row to denominator $2^b$ without enlarging
its support: round down the first at most three masses and give the last
mass the remainder. Its total-variation error is at most $3\cdot2^{-b}$.
Stochastic contraction bounds the terminal distribution discrepancy by
$3N2^{-b}$. The Frobenius diameter of $\mathcal D_2$ is $\sqrt2$;
therefore the additional error is at most
$3\sqrt2 N2^{-b}\le\varepsilon/4$ under
\eqref{eq:compilerbits58}. Since $1/\sqrt2+1/4<1$, the total is strictly
less than $\varepsilon$. One uniform $b$-bit integer samples each row.
There is no initialization draw and no sampled quantum output in this
count.

For completeness we give finite construction costs. Enumerating the net
uses $O(m^2)$ rational points. For each of $|\mathcal A|K$ rows,
enumerating at most $\binom{K-1}{3}$ triples and solving fixed-size
systems uses polynomially many rational operations. Cramer's rule shows
that the exact row coefficients have bit length
$O(\log(m+1)+B_{\mathcal A})$, where $B_{\mathcal A}$ bounds the supplied
rotation-entry bit lengths. The denominator $r^N$ has bit length
$O(N\log(m+1))$. Expanded decoders therefore have
$O(B_\beta+(N+1)\log(m+1))$ bits each. A naive full enumeration needs
$O(|\mathcal A|K^4)$ fixed-size rational solves, each of polynomial bit
cost. This proves the algorithmic assertion as well as the bit bounds.
\end{proof}
Sparse row tables require
$O(|\mathcal A|K(B_{\mathcal A}+\log(K+1)+\log(m+1)))$ bits before
rounding, and $O(|\mathcal A|K(b+\log(K+1)))$ bits after rounding.
Expanded matrix decoders require
$O(K(B_\beta+(N+1)\log(m+1)))$ bits. The persistent label uses
$\lceil\log_2 K\rceil$ bits; the horizon counter, temporary sampler, and
arithmetic working storage are separately charged. Unlike a statement
only about output bit length, the preceding proof gives a terminating
polynomial table-construction algorithm for this specific family. It
does not solve general minimum-width realization or certify a spectral
gap. The full legal density-matrix constraint is maintained at every
rounding step.

\begin{proposition}[Supplied rational unit seeds]\label{prop:seedextension59}
Theorem~\ref{thm:rationalcompiler58} remains valid for any supplied
rational unit Bloch seed, with the label bound replaced by $13m^2+1$.
\end{proposition}
\begin{proof}
Add the seed to the rational unit net if it is absent. The inner-ball
inclusion is unchanged, every contracted rotated vertex still has an
exact rational sparse row, and initialization is now deterministic at
the added seed. The same induction and terminal rescaling apply.
\end{proof}
The implementation accepts the supplied finite rotation list and unit
seed as exact rational data. Orthogonality, determinant one and unit
seed norm are checked by rational identities before table allocation.

\section{A rational experiment with explicit spectral constants}\label{sec:effective61}
The original rational pair in Theorem~\ref{thm:rationalseparation58}
now has a sharp linear fixed-error law, but its spectral constant is
still supplied existentially by Bourgain--Gamburd. We give a second
finite input for which every constant in a useful linear lower bound
is numerical. The deep input here is the classical
Lubotzky--Phillips--Sarnak theorem, not a new spectral-gap calculation.

Let $\sigma_1,\sigma_2,\sigma_3$ be the Pauli matrices in their usual
order, and put $D(x)=(I+\sum_jx_j\sigma_j)/2$ for $x\in\mathbb B_3$.
Thus $D(\mathbb B_3)=\mathcal D_2$ and
$\norm{D(x)-D(y)}_{\rm F}=\norm{x-y}/\sqrt2$.
Consider the following \emph{rational Bloch rotations}:
\begin{equation}\label{eq:lpsinput61}
 A_1=\frac15\begin{pmatrix}5&0&0\\0&-3&-4\\0&4&-3\end{pmatrix},\quad
 A_2=\frac15\begin{pmatrix}-3&0&4\\0&5&0\\-4&0&-3\end{pmatrix},\quad
 A_3=\frac15\begin{pmatrix}-3&-4&0\\4&-3&0\\0&0&5\end{pmatrix}.
\end{equation}
Take the alphabet
\[
 \mathcal A_{\rm LPS}=
       \{I,A_1,A_1^{\mathsf T},A_2,A_2^{\mathsf T},A_3,A_3^{\mathsf T}\}
\]
and the rational pure seed $e_1$, so $D(e_1)=P_0$ as in
\eqref{eq:rationaldata58}. The induced unitary lifts can be taken as
$(I-2i\sigma_j)/\sqrt5$. They are algebraic, not rational matrices;
the rational input in \eqref{eq:lpsinput61} describes the physical
Bloch action directly. This is a different alphabet from the rational
unitary pair of Theorem~\ref{thm:rationalseparation58}.

\begin{lemma}[The classical six-rotation spectral bound]\label{lem:lpsgap61}
Let $P_{\rm LPS}$ average uniformly over the six nonidentity letters
in $\mathcal A_{\rm LPS}$. On normalized $S^2$,
\begin{equation}\label{eq:lpsgap61}
 \norm{P_{\rm LPS}}_{L^2_0(S^2)\to L^2_0(S^2)}=\frac{\sqrt5}{3}.
\end{equation}
Their generated subgroup is dense in $SO(3)$.
\end{lemma}
\begin{proof}
The matrices in \eqref{eq:lpsinput61} are the adjoint actions of the
unit Hamilton quaternions
$(1+2\mathbf i)/\sqrt5,(1+2\mathbf j)/\sqrt5,(1+2\mathbf k)/\sqrt5$.
Together with their inverses they are exactly the norm-five set in the
Lubotzky--Phillips--Sarnak theorem \cite[Theorems~1.3 and~1.5]{LPS61};
see the precise restatement \cite[Theorem~1.1]{PLP61pre}, which gives
$2\sqrt p/(p+1)$ for this set at $p=5$. This proves
\eqref{eq:lpsgap61}. The cited theorem uses the theory of automorphic
forms and Deligne's theorem; no finite matrix enumeration replaces it.

Density can also be checked directly. Each rotation has cosine
$-3/5$. If its angle were a rational multiple of $2\pi$, twice this
cosine would be a rational algebraic integer, contrary to $-6/5\notin\Z$.
Its powers are dense in its circle subgroup. The three coordinate-axis
circle subgroups generate $SO(3)$, as follows either from their
infinitesimal generators and brackets or from Euler rotations. Hence
the closed generated subgroup is $SO(3)$.
\end{proof}

\begin{theorem}[An effective linear noisy law and an exponential endpoint]
\label{thm:effectivelps61}
For the fixed input \eqref{eq:lpsinput61}, the terminal target is
$D(A_w e_1)$, every decoder belongs to $\mathcal D_2$, and error is
Frobenius error. For $0<\varepsilon<1/\sqrt2$ put
$\zeta=1-\sqrt2\varepsilon$. Every admissible realization satisfies
\begin{equation}\label{eq:lpsoccupation61}
 \#\{1\le t\le N:|S_t|\le k\}
 \le 1+6\log(12k)+648\pi^2\frac{1-\zeta}{\zeta}k
 \le\frac{6500}{\zeta}k.
\end{equation}
In particular,
\begin{equation}\label{eq:lpslinear61}
 \frac{\zeta N}{6500}\le W_{N,\varepsilon}
                       \le52+\frac{104N}{\varepsilon}.
\end{equation}
For rational $\varepsilon$ in this interval the upper machine has
uniformly constructed rational legal decoders and dyadic stochastic
rows, each supported on at most four labels. It uses at most $Nb$
fair bits for hidden transitions, where
$b=\lceil\log_2(24(N+1)/\varepsilon)\rceil$.

At every horizon $N\ge0$, the simultaneous exact minimum profile is
\begin{equation}\label{eq:lpsexact61}
                  |S_t|_{\min}=5^t\quad(0\le t\le N).
\end{equation}
The same profile remains necessary and is attained throughout
\begin{equation}\label{eq:lpsrobust61}
                         0\le\varepsilon\le25^{-N}/16.
\end{equation}
For $\varepsilon\ge1/\sqrt2$, one label suffices.
\end{theorem}
\begin{proof}
For a spherical chordal cap of radius $h\le1$, normalized area is
exactly $h^2/4$. Consequently the support of a mixture of $k$ kernels
$K_h$ occupies at most $kh^2/4$ of each group orbit. Take $h=k^{-1/2}$.
In Lemma~\ref{lem:entropydefect61} we may take
\[
 \theta=1/4,\qquad \kappa^2=5/9,\qquad
 \gamma=(1-\kappa^2)(1-\theta)=1/3.
\]
Lemma~\ref{lem:smoothentropy61} gives $L_h=\log(12k)$ and $B_3=24$.
The proof of Theorem~\ref{thm:sharpoccupation61}, with these constants
and the normalization $\zeta=1-\sqrt2\varepsilon$, gives the first
inequality in \eqref{eq:lpsoccupation61} as a real bound before taking
an integer ceiling. For $k\ge1$, $\log(12k)\le3k$, and $\pi^2<10$.
Thus its right side is at most
$19k+6480(1-\zeta)k/\zeta\le6500k/\zeta$.
Counting all positive cuts proves the lower bound in
\eqref{eq:lpslinear61}.

For the upper bound apply the rational compiler of
Theorem~\ref{thm:rationalcompiler58} to the displayed seven-letter list
and the first-axis seed. At rational tolerance take
$\beta=1-\varepsilon/2$ and
$m=\max\{2,\lceil\sqrt{2N/\varepsilon}\rceil\}$.
It gives an exact rational realization of $D(\beta A_we_1)$ on at most
$13m^2\le52+52N/\varepsilon$ labels, with sparse common rows.
The target change has Frobenius norm $\varepsilon/(2\sqrt2)$.
The stated dyadic precision adds at most
$3\sqrt2N2^{-b}\le\sqrt2\varepsilon/8$, so the total is less than
$\varepsilon$. This proves the looser advertised numerical upper bound
and the fair-bit statement. For a real tolerance choose a rational
$\widetilde\varepsilon\in(\varepsilon/2,\varepsilon)$ and use the
same construction at $\widetilde\varepsilon$; the bound
$52+104N/\varepsilon$ follows. The algorithmic assertion uses rational
input, not an oracle for an unencoded real number.

We next prove the exact profile independently of the spectral bound.
Lemma~\ref{lem:freequaternion58} shows that the three norm-five
quaternion generators freely generate modulo scalars, hence their
adjoint rotations form a free group on $a,b,c$. The stabilizer of
$e_1$ in this free group is exactly $\langle a\rangle$. Indeed, a
rotation fixing $e_1$ commutes with the rotation $a$ about that axis.
In a free group the centralizer of this generator is its cyclic
subgroup: write a reduced word outside $\langle a\rangle$ as
$a^rva^s$, where $v$ is nonempty and begins and ends with neither
$a$ nor $a^{-1}$. Its conjugate of $a$ reduces to
$a^r v a v^{-1}a^{-r}$, which is not $a$. The converse, that every
power of $a$ fixes $e_1$, is immediate.

The orbit is therefore indexed by the right cosets of $\langle a\rangle$.
Each coset has a unique reduced representative with no terminal
$a^{\pm1}$, obtained by removing the maximal terminal power of $a$.
It has least word length in that coset. At length $\ell\ge1$ there
are $4\cdot5^{\ell-1}$ such representatives. Identity padding gives
exactly the ball of these cosets at command length $t$, whose cardinality
is $1+4\sum_{\ell=1}^t5^{\ell-1}=5^t$.
Theorem~\ref{thm:extreme57} applies because pure states are extreme in
$\mathcal D_2$, proving \eqref{eq:lpsexact61} at every cut; deterministic
orbit storage attains the entire profile at once.

For robustness, distinct cut-$t$ Bloch targets have coordinates in
$5^{-t}\Z^3$ and hence distance at least $5^{-t}$. One fixed suffix
of length $N-t$ preserves these distances. If $z_i\in\mathbb B_3$
is the conditional terminal decoder mean from cut label $i$,
wordwise error at a pure target $n$ implies
\[
 \sum_i p_i\norm{z_i-n}^2\le2\sqrt2\varepsilon.
\]
Some positive-mass label is within $\sqrt{2\sqrt2\varepsilon}$ of $n$.
Under \eqref{eq:lpsrobust61}, twice this radius is strictly less than
$5^{-N}\le5^{-t}$. The corresponding closed balls for distinct
prefix orbit points are disjoint, and so require distinct labels.
The exact construction attains this bound. No fixed positive tolerance
is asserted to lie in \eqref{eq:lpsrobust61} for all horizons.
Finally $D(0)=I/2$ is at Frobenius distance $1/\sqrt2$ from every pure
state, proving the one-label assertion.
\end{proof}

\begin{remark}[What is effective]\label{rem:effective61}
The input, orbit counts, robust tolerance, spectral norm, constants in
\eqref{eq:lpsoccupation61}--\eqref{eq:lpslinear61}, and rational compiler
are explicit. The constant $6500$ is conservative; it is not an optimal
leading constant. The spectral norm is imported from
\cite{LPS61,PLP61pre}, not inferred from finite harmonic tests.
For the earlier five-letter rational-unitary example,
Theorem~\ref{thm:sharpoccupation61} removes the logarithm but does not
supply a numerical Bourgain--Gamburd constant. Neither statement claims
a complete two-parameter phase diagram as the error tends to zero.
The finite-table and fair-bit resources refer to a classical hidden
label and a numerical matrix output, not quantum-state preparation.
\end{remark}

\section{The six-command sphere and variable accuracy}\label{sec:noidle62}
The full spectral measure in Lemma~\ref{lem:lpsgap61} already lives on
six nonidentity commands. The following result uses exactly those six
commands for the experiment itself, with no idle instruction.

\begin{theorem}[An explicit no-idle experiment]\label{thm:noidle62}
Let $A_1,A_2,A_3$ be the rational matrices in \eqref{eq:lpsinput61}.
Take exactly
\[
 \mathcal A_6=\{A_1,A_1^{\mathsf T},A_2,A_2^{\mathsf T},A_3,A_3^{\mathsf T}\},
 \qquad x_0=e_1,\qquad F(w)=D(A_we_1).
\]
Every decoder must lie in $\mathcal D_2$, and error is Frobenius error.
For $0<\varepsilon<1/\sqrt2$, put
$\zeta=1-\sqrt2\varepsilon$ and $\theta=(1-\zeta)/\zeta$.
Then every admissible width profile satisfies
\begin{equation}\label{eq:noidleoccupation62}
 \#\{1\le t\le N:|S_t|\le k\}
             \le1+6\log(12k)+648\pi^2\theta k.
\end{equation}
For $N\ge2$, the minimum peak width obeys the explicit joint bounds
\begin{equation}\label{eq:noidlejoint62}
 \min\left\{\frac{e^{(N-1)/12}}{12},
                   \frac{N-1}{1296\pi^2\theta}\right\}
 \le W_{N,\varepsilon}
 \le\min\{5^N,\ 52+104N/\varepsilon\}.
\end{equation}
At every horizon the exact simultaneous cut profile is $5^t$,
$0\le t\le N$, and the same profile is necessary and attained for
$0\le\varepsilon\le25^{-N}/16$. At
$\varepsilon\ge1/\sqrt2$, one label suffices.
For any fixed $\varepsilon_0<1/\sqrt2$ and sequence
$0<\varepsilon_N\le\varepsilon_0$ with
$\log(1/\varepsilon_N)=o(N)$,
\begin{equation}\label{eq:sixjointmatched62}
                  W_{N,\varepsilon_N}=\Theta(N/\varepsilon_N).
\end{equation}
The constants in this last comparison can be taken independent of the
sequence; its starting horizon can depend on the sequence.
\end{theorem}
\begin{proof}
The classical full spherical norm is $\sqrt5/3$ for the uniform law on
exactly $\mathcal A_6$ \cite[Theorem~1.1]{PLP61pre}. Use the drift proof
of Theorem~\ref{thm:driftoccupation62}, taking $h=k^{-1/2}$ and the
sharper sphere constants $\gamma=1/3$, $L_h=\log(12k)$, $B_3=24$
as in Theorem~\ref{thm:effectivelps61}. This gives
\eqref{eq:noidleoccupation62}, with arbitrary deterministic commands
in the gaps. It imports no uncomputed spectral constant.

For a full width bound $k$, at least one of
$6\log(12k)$ and $648\pi^2\theta k$ is at least $(N-1)/2$.
Inverting the corresponding inequality gives the lower bound in
\eqref{eq:noidlejoint62}. The sparse rational construction in
Theorem~\ref{thm:effectivelps61} applies unchanged to the six-letter
sublist; no row for the identity is needed. It gives the second upper
bound, with the same rational-decoder and dyadic-row guarantees for
rational tolerances.

For the exact profile, recall the independently proved free-group and
stabilizer calculation in Theorem~\ref{thm:effectivelps61}. The orbit is
the right-coset space of $\langle a\rangle$ in the free group on $a,b,c$,
where $a=A_1$. Every coset has a unique shortest representative $w$
with no terminal $a^{\pm1}$. If its length is $\ell\le t$, the group
word $w a^{t-\ell}$ has length $t$ and the same action on $e_1$.
In chronological execution this means applying the $a$'s first; they
fix the seed. There is no cancellation when $w$ is nonempty, and the
empty representative is padded by $a^t$. Thus even without an identity
letter, the length-$t$ reachable orbit is exactly the radius-$t$ coset
ball. Its size is $5^t$. The extreme-orbit theorem gives the exact
profile and deterministic orbit storage gives the upper bound $5^N$.

The denominator-five lattice and fixed-suffix isometry argument of
Theorem~\ref{thm:effectivelps61} does not require identity padding once
the orbit count is established. It gives the same robust interval.
The center $I/2$ gives the final constant-width assertion.

Finally $\theta\le C_{\varepsilon_0}\varepsilon_N$, while
$\log(N/\varepsilon_N)=o(N)$. The exponential alternative in the
minimum on the left of \eqref{eq:noidlejoint62} eventually exceeds a
constant times $N/\varepsilon_N$. The other alternative already has
that lower order. The upper bound proves \eqref{eq:sixjointmatched62}.
\end{proof}

\begin{remark}[Three separate accuracy regimes]\label{rem:threeaccuracy62}
The exact count and its robust interval, the explicit two-parameter
bounds \eqref{eq:noidlejoint62}, and the subexponential-accuracy
asymptotic \eqref{eq:sixjointmatched62} are different statements.
For example, $\varepsilon_N=N^{-a}$ gives width $\Theta(N^{a+1})$,
and $\varepsilon_N=e^{-N^b}$, $0<b<1$, gives
$\Theta(Ne^{N^b})$. Neither assertion extrapolates the exact plateau
to a fixed positive error or resolves its full crossover at
$\log(1/\varepsilon_N)\asymp N$. The constant $6500$ remains a safe
fixed-error simplification of \eqref{eq:noidleoccupation62}, not an
optimal leading constant. The seven-letter experiment and the older
rational-unitary pair remain separate retained examples.
\end{remark}

\section{A rational compiler for adaptive observation processes}
\label{sec:adaptivecompiler59}
The common-row construction also yields a causal approximation, but its
proof requires control of the full observation law. Correct terminal
means alone are not enough. We give that additional argument here.

The physical state is $x\in B_3$, initialized at a supplied rational
$x_0\in S^2$. A finite rational rotation list $\{A_a\}\subset SO(3)$
acts by $x\mapsto A_ax$ and emits a null symbol. A finite probe list has
rational affine probability laws
\begin{equation}\label{eq:affineprobes59}
 q_{j,y}(x)=b_{j,y}+\ip{\ell_{j,y}}x,\qquad
 \sum_y b_{j,y}=1,\quad \sum_y\ell_{j,y}=0,
 \quad b_{j,y}-\norm{\ell_{j,y}}\ge\eta>0.
\end{equation}
A probe emits a fresh draw and leaves $x$ unchanged. Controls may depend
on the entire observed transcript. A simulator has finite retained
labels and a nonnegative joint emission/transition matrix for each
control and output; these matrices may depend on horizon and epoch.
Let $\mathsf{PW}_{N,\varepsilon}$ denote its minimum peak label count
for one total-variation budget $\varepsilon$ on the entire transcript,
uniformly over all adaptive randomized policies. This paragraph defines
a causal model independently of the terminal matrix-output model.

Put $A=\max_j\sum_y\norm{\ell_{j,y}}^2$ and
$C_0=\max(3,\max_j(M_j-1))$. The lower probability bound in
\eqref{eq:affineprobes59} is used only for the relative-entropy upper
estimate, not for the finite-testing converse below.

\begin{theorem}[Adaptive rational and fair-bit realization]
\label{thm:adaptivecompiler59}
For rational $0<\varepsilon<1$ and integer $N\ge0$, choose
\begin{equation}\label{eq:adaptiveparameters59}
 m=\max\left(2,\left\lceil\sqrt{
                    \frac{4AN^2}{\eta\varepsilon^2}}\right\rceil\right),
 \qquad r=1-m^{-2},\qquad
 b=\max\left(0,\left\lceil\log_2
                 \frac{2C_0\max(1,N)}{\varepsilon}\right\rceil\right).
\end{equation}
There is an error-$\varepsilon$ causal realization on at most
$13m^2+1$ labels. Before probability rounding every rotation row is
rational, has at most four nonzero entries, and has barycentric mean
$rA_av$. Probe emissions leave the label unchanged. After rounding,
all emission and transition probabilities are dyadic, with at most $b$
fair bits per control and $Nb$ bits in total. The tables are stationary
within the chosen horizon.

An exact rational algorithm constructs these tables in time polynomial
in the numerical parameters $N,\varepsilon^{-1},A/\eta$, the explicit
input size, and the input bit lengths. In particular, for fixed probe
and rotation lists,
\begin{equation}\label{eq:adaptivewidth59}
                 \mathsf{PW}_{N,\varepsilon}
                    =O(1+N^2/\varepsilon^2).
\end{equation}
The bit count now includes the actual finite probe outputs as well as
hidden transitions; it does not describe quantum-state preparation.
\end{theorem}
\begin{proof}
Take the rational net $V_m$ from Lemma~\ref{lem:rationalnet58}, including
zero, and add $x_0$ if needed. This adds at most one label and cannot
shrink the convex hull. Initialize deterministically at $x_0$. For every
rotation and label $v$, the same exact triple search supplies a row with
mean $rA_av$. At probe $j$, emit $y$ with probability $q_{j,y}(v)$ and
retain $v$. All points lie in $B_3$.

Fix an arbitrary policy. On the same finite trajectory space of controls,
outputs and hidden labels, introduce an auxiliary law $P$: hidden labels
follow the just constructed rotation kernels, but every probe emits
$q_{j,y}(X)$ from the true physical state $X$, independently of the
hidden label conditional on the public history. Its public marginal is
exactly the target experiment. Let $Q$ be the constructed simulator law.
The policy kernels, initial label and hidden rotation kernels are the
same under $P,Q$; only probe emissions differ.

Under $P$, if the current physical and hidden points are $X,V$, both
$X$ and every nonzero hidden vertex are unit vectors. At a rotation
step, for $V\ne0$,
\begin{align*}
 \E[\norm{V'-A_aX}^2\mid X,V,a]
 &\le 2-2r\ip V X\\
 &=r\norm{V-X}^2+2(1-r).
\end{align*}
For $V=0$ choose its row to remain at zero; the same upper bound holds
because the left side is one and the right side is $2-r\ge1$.
The added seed is a unit vector, so these two cases exhaust the state
set. This is why a general interior seed has not been silently added
to the theorem. A probe changes neither point. Iteration, with the
control selected adaptively, gives
\begin{equation}\label{eq:processmoment59}
                  \E_P\norm{V_t-X_t}^2\le 2t/m^2.
\end{equation}
The inequality is unconditional; no posterior error bound after an
unlikely observation is assumed.

For probability vectors $p,q$ with $q_y\ge\eta$, the inequality
$\log z\le z-1$ gives
\[
 D(p\Vert q)\le\sum_y(p_y-q_y)^2/q_y.
\]
Consequently the conditional divergence at a probe is at most
$A\norm{X-V}^2/\eta$. The finite-trajectory relative-entropy chain rule
cancels the common policy and rotation kernels. Applying
\eqref{eq:processmoment59} at the probe times gives
\begin{equation}\label{eq:processKL59}
 D(P\Vert Q)\le \frac{2AN^2}{\eta m^2},\qquad
 \TV(P_{\rm pub},Q_{\rm pub})
       \le\sqrt{D(P\Vert Q)/2}
       \le \frac Nm\sqrt{A/\eta}\le\varepsilon/2.
\end{equation}
Here marginalization can only decrease divergence. These elementary
entropy facts apply to the joint classical trajectory, not to a quantum
state. One may verify the total-variation inequality by partitioning
into the set where $P\ge Q$ and its complement: the log-sum inequality
reduces divergence to a Bernoulli divergence, whose second derivative
in its first parameter is at least four. This yields
$D(P\Vert Q)\ge2\TV(P,Q)^2$. The chain rule follows by factoring each
trajectory probability into its conditional kernels. All probe
probabilities in $Q$ are positive, so absolute continuity causes no
exception. If $A=0$, every probe is constant and the divergence is zero.

Round each rotation row and each probe row by rounding all but its last
positive mass down to denominator $2^b$ and assigning the remainder to
the last. A row supported on $d$ points changes in total variation by
at most $(d-1)2^{-b}$. Thus every one-step joint kernel changes by at most
$C_0 2^{-b}$. Coupling the two label--transcript processes until their
first discrepancy proves, for every adaptive policy, an additional
path error at most $NC_0 2^{-b}\le\varepsilon/2$. Combining this with
\eqref{eq:processKL59} proves the asserted error. Each rounded row can
be sampled by one uniform $b$-bit integer and cumulative integer masses;
zero or deterministic rows may use fewer bits.

The net and an added rational unit seed are explicit. For every
command--vertex pair the exact enumeration in
Lemma~\ref{lem:rationalnet58} examines at most $\binom{k-1}{3}$ triples,
where $k\le13m^2+1$. The total is $O(|\mathcal A|k^4)$ constant-size
rational linear-algebra operations, in addition to the probe tables.
All determinants, comparisons and row identities involve rational
numbers with polynomially bounded bit lengths in the supplied data and
$\log m$; dyadic rounding adds $b$ bits. The computation of $m$ and $b$
can use integer comparisons, not floating-point decisions. This proves
the stated uniform, numerical-parameter complexity. It is a construction,
not a minimum-width algorithm, and its exhaustive fallback need not be
practical for large horizons.
\end{proof}

\subsection{An effective process converse for the rational free experiment}
Keep the rational rotations and seed of
\eqref{eq:rationaldata58}--\eqref{eq:rationalalphabet58}, now exposing the
three nondisturbing binary probes
\begin{equation}\label{eq:weakprobes59}
       q_{j,+}(x)=(1+\lambda x_j)/2,\qquad
       q_{j,-}(x)=(1-\lambda x_j)/2,
       \quad j=1,2,3,
\end{equation}
where $0<\lambda<1$ is rational. This defines a fixed rational classical
experiment. The probing law is not an implementation of nondisturbing
quantum measurement. Let $\mathsf{PW}$ refer to the causal model above,
not to terminal density-matrix width $W$.

\begin{theorem}[Effective repeated-observation lower bound]
\label{thm:processlower59}
Let $0\le\varepsilon<1$ and $0<\delta<1-\varepsilon$. For an integer
$t\ge0$ put
\begin{equation}\label{eq:testlength59}
 L_t=\max\left(1,\left\lceil18\lambda^{-2}625^t
                                      \log(6/\delta)\right\rceil\right).
\end{equation}
At every horizon $N\ge t+3L_t$, any error-$\varepsilon$ causal machine
has
\begin{equation}\label{eq:processcutlower59}
 |S_t|\ge(1-\varepsilon-\delta)(2\cdot3^t-1).
\end{equation}
If $\varepsilon+\delta<1/2$, it has the stronger lower bound
$|S_t|\ge2\cdot3^t-1$. More generally, the same bounds hold at every
cut $u$ with $t\le u\le N-3L_t$ after padding the prefixes by probes.
In particular, for each fixed $0<\varepsilon<1$, there are explicit
positive constants $c_{\varepsilon,\lambda},C_{\varepsilon,\lambda}$
such that, for all sufficiently large $N$,
\begin{equation}\label{eq:processgrowth59}
 c_{\varepsilon,\lambda}N^{\log3/\log625}
 \le\mathsf{PW}_{N,\varepsilon}
 \le C_{\varepsilon,\lambda}(N+1)^2.
\end{equation}
The converse uses no spectral-gap theorem. Its exponent is not asserted
to be sharp.
\end{theorem}
\begin{proof}
The rational freeness and stabilizer argument in
Theorem~\ref{thm:rationalseparation58} gives $2\cdot3^t-1$ different
reachable pure Bloch vectors at time $t$. Their coordinates lie in
$25^{-t}\Z^3$. Hence each distinct pair differs in some coordinate by
at least $25^{-t}$, and the associated plus-probabilities in
\eqref{eq:weakprobes59} differ by at least $\lambda25^{-t}/2$.
Use one common continuation: repeat each of the three probes $L_t$
times. Hoeffding's elementary Bernoulli bound and a union bound give
\[
 \mathbb P_x\left\{\max_j|\widehat q_j-q_{j,+}(x)|
                    \ge\lambda25^{-t}/6\right\}
 \le6\exp(-L_t\lambda^2 625^{-t}/18)\le\delta.
\]
The empirical neighborhoods for the different targets are disjoint.
At a cut label $z$, let $Q_z$ be the law of this same continuation.
Marginalizing prefix outputs and using total-variation correctness
shows that the continuation law from each prefix gives its own
neighborhood probability at least $1-\varepsilon-\delta$.
Writing this law as $\sum_z\alpha_x(z)Q_z$ and summing over the disjoint
neighborhoods yields \eqref{eq:processcutlower59}. If the probability
exceeds $1/2$, each neighborhood needs a different label assigning it
probability greater than $1/2$, giving the stronger conclusion.
The proof does not condition on the success of a previous observation.
Padding by arbitrary fixed probes leaves the target states unchanged
and proves the assertion at later cuts. Unused final controls may also
be probes. A terminal cut without enough remaining observations is not
claimed to satisfy this lower bound.

For the power lower bound take $\delta=(1-\varepsilon)/2$ and
$B=54\lambda^{-2}\log(6/\delta)+4$. Then
$t+3L_t\le B625^t$. For $N\ge B$, choose
$t=\lfloor\log_{625}(N/B)\rfloor$. Since
$3^t\ge(N/B)^{\log3/\log625}/3$ and $2\cdot3^t-1\ge3^t$,
\eqref{eq:processcutlower59} gives, for example,
$c_{\varepsilon,\lambda}=(1-\varepsilon)B^{-\log3/\log625}/6$.
The upper bound follows from Theorem~\ref{thm:adaptivecompiler59} with
$\eta=(1-\lambda)/2$ and $A=\lambda^2/2$. For irrational requested
$\varepsilon$, construct at any rational tolerance between
$\varepsilon/2$ and $\varepsilon$. All constants can thus be bounded
from the displayed data without computing a group spectral gap.
\end{proof}

The tests in this theorem consume horizon. The full-state terminal
exponential profile and the causal lower bound therefore express
different requirements. Conversely, the causal upper construction
controls one joint transcript law under adaptive control, not just its
one-time marginals. These process statements do not by themselves sharpen the terminal
spectral-gap converse. Its logarithmic loss is removed separately by
Theorem~\ref{thm:sharpoccupation61}.

\appendix
\section{An existential higher-dimensional endpoint}
\subsection{A fixed nonabelian family with an exponential exact endpoint}
The next example uses two established inputs, explicitly separated from
the realization argument. Breuillard and Gelander prove that a dense
subgroup of a connected semisimple real Lie group contains a dense free
subgroup on two generators \cite[Theorem~1.1]{BreuillardGelander57}; their
later correction supplies repairs and details for other parts of the
original proof \cite{BreuillardGelanderCorrection57}. The spectral-gap input
for dense algebraic unitary generators is \cite{BG2012}. No algorithm for
finding free generators, or for certifying a spectral-gap constant, is
asserted here.

\begin{lemma}[A line with trivial algebraic projective stabilizer]
\label{lem:transcendentalseed57}
Fix a real number $\tau$ transcendental over $\Q$ and put
\[
 v_\tau=\frac{(1,\tau,\ldots,\tau^{d-1})^{\mathsf T}}
                  {(\sum_{j=0}^{d-1}\tau^{2j})^{1/2}},\qquad
 P_\tau=v_\tau v_\tau^*.
\]
An invertible matrix with algebraic entries that preserves the line
$\C v_\tau$ is scalar.
\end{lemma}
\begin{proof}
If $Mv_\tau=\lambda v_\tau$, the eigenvalue $\lambda$ is algebraic because
it is a root of the characteristic polynomial of the algebraic matrix $M$.
Each row of $(M-\lambda I)(1,\tau,\ldots,\tau^{d-1})^{\mathsf T}=0$ is a
polynomial identity in $\tau$ with algebraic coefficients. Transcendence
over $\Q$ implies transcendence over the field of algebraic numbers: an
algebraic relation over a finite number field would make $\tau$ algebraic
over $\Q$ by transitivity of algebraic extensions. Every row polynomial
therefore has all coefficients zero. Thus $M=\lambda I$.
\end{proof}
One may take a fixed computable transcendental such as
$\tau=\sum_{n=1}^\infty10^{-n!}$. For completeness, its truncations have
rational denominator $10^{m!}$ and strictly positive tail at most
$2\,10^{-(m+1)!}$. For a nonzero integer polynomial $A$ of degree $D$, a
rational $p/q$ not among its finitely many roots satisfies
$|A(p/q)|\ge q^{-D}$. If $A(\tau)=0$, a bounded derivative near $\tau$
would instead bound $|A(p/q)|$ by a constant times $q^{-(m+1)}$ along these
truncations, a contradiction for large $m$. This also records precisely
why the chosen seed is not part of the rational-input specialization.

\begin{theorem}[Exponential exact width and polynomial noisy width]
\label{thm:freeendpoint57}
For every $d\ge2$, with a transcendental seed and existentially selected generators, there is a fixed five-letter algebraic unitary alphabet
\[
                       \{I,U,U^{-1},V,V^{-1}\}\subset SU(d)
\]
and a fixed rank-one seed $P_\tau$ as above such that the full matrix
interface $F(g)=gP_\tau g^*$, with legal decoder set $\mathcal D_d$, has
\begin{equation}\label{eq:freeexact57}
                      W_{N,0}=2\,3^N-1\qquad(N\ge0).
\end{equation}
For each fixed $0<\varepsilon<r_d=\sqrt{1-1/d}$, its Frobenius-error width
instead satisfies
\begin{equation}\label{eq:freenoise57}
 c_{d,\varepsilon}\left(\frac{N}{\log(N+2)}\right)^{d-1}
 \le W_{N,\varepsilon}\le C_{d,\varepsilon}(N+1)^{d-1}.
\end{equation}
At $\varepsilon\ge r_d$ one label suffices. For the same alphabet and seed,
every fixed amplitude $0<a<1$ admits the exact polynomial upper bound in
\eqref{eq:matrixupper57} and the corresponding lower bound in
\eqref{eq:matrixrate57}.
\end{theorem}
\begin{proof}
The subgroup of $SU(d)$ consisting of matrices with algebraic entries is
dense. One direct verification is to approximate a unitary matrix by a
nonsingular matrix with entries in $\Q+i\Q$, apply Gram--Schmidt, and
multiply its final column by the inverse determinant of the resulting
unitary matrix. All entries produced are algebraic; continuity near the
original unitary matrix gives density. Apply
\cite[Theorem~1.1]{BreuillardGelander57} to this subgroup. It contains
algebraic $U,V$ freely generating a dense free group $\Gamma$ of rank two.
The scalar matrices in $SU(d)$ form a finite group. Since a free group is
torsion-free, $\Gamma$ contains no nonidentity scalar matrix.
Lemma~\ref{lem:transcendentalseed57} now shows that the orbit map
$\Gamma\to\mathcal D_d$, $g\mapsto gP_\tau g^*$, is injective.

With the identity available for padding, the products at length $N$ are
exactly the elements of reduced free-word length at most $N$. There are
four choices for the first letter of a nonempty reduced word and three
for each subsequent letter. Thus the ball has cardinality
\[
               1+4\sum_{j=1}^N3^{j-1}=2\,3^N-1.
\]
Apply Theorem~\ref{thm:extreme57} to obtain \eqref{eq:freeexact57}. In fact the
entire pointwise minimum profile is $|S_t|=2\,3^t-1$.

Assign a symmetric positive law to the four nonidentity commands and
positive mass to the identity. The algebraic dense-generator theorem
\cite{BG2012}, with laziness, gives \eqref{eq:groupgap57}. Consequently
Theorem~\ref{thm:matrixwidth57} applies to this same fixed alphabet and
seed. It yields \eqref{eq:freenoise57}, the threshold $r_d$, and the asserted
strict-amplitude conclusion.
\end{proof}
This is a full legal-matrix-output statement. It does not resolve the
zero-error unit-amplitude problem for a single scalar planar readout:
that readout need not expose every orbit point as an extreme decoder
value. Nor does it identify a sampled measurement outcome with a density
matrix stored in a classical label. The distinction between a prescribed
convex output set and a larger simplex is essential here.

\begin{thebibliography}{99}
\item[]\textit{Stochastic semigroups and positive realization}
\bibitem{Flor} P. Flor, On groups of non-negative matrices, \emph{Compositio Mathematica} \textbf{21} (1969), no. 4, 376--382.
\bibitem{Schwarz58} \v S. Schwarz, On the structure of the semigroup of stochastic matrices, \emph{Publications of the Mathematical Institute of the Hungarian Academy of Sciences, Series A} \textbf{9} (1964), 297--311.
\bibitem{BF04} L. Benvenuti and L. Farina, A tutorial on the positive realization problem, \emph{IEEE Transactions on Automatic Control} \textbf{49} (2004), no. 5, 651--664.
\bibitem{Heller58} A. Heller, On stochastic processes derived from Markov chains, \emph{Annals of Mathematical Statistics} \textbf{36} (1965), no. 4, 1286--1291.
\item[]\textit{Automata, weighted series, and recognition}
\bibitem{Rabin} M. O. Rabin, Probabilistic automata, \emph{Information and Control} \textbf{6} (1963), no. 3, 230--245.
\bibitem{Paz} A. Paz, \emph{Introduction to Probabilistic Automata}, Academic Press, New York, 1971.
\bibitem{BR} J. Berstel and C. Reutenauer, \emph{Noncommutative Rational Series with Applications}, Encyclopedia of Mathematics and its Applications 137, Cambridge University Press, Cambridge, 2011.
\bibitem{FP} N. Fijalkow and C. Paperman, Monadic second-order logic with arbitrary monadic predicates, arXiv:1709.03117 (2017).
\bibitem{BP} A. Brodsky and N. Pippenger, Characterizations of 1-way quantum finite automata, \emph{SIAM Journal on Computing} \textbf{31} (2002), no. 5, 1456--1478.
\bibitem{AF57} A. Ambainis and R. Freivalds, 1-way quantum finite automata: strengths, weaknesses and generalizations, arXiv:quant-ph/9802062 (1998).
\bibitem{FSZ} Y. Feng, L. Song, and L. Zhang, Distribution-based bisimulation and bisimulation metric in probabilistic automata, arXiv:1512.05027 (2015).
\bibitem{Steinberg} B. Steinberg, Topological dynamics and recognition of languages, arXiv:1306.1468 (2013).
\item[]\textit{Compact groups and orbit geometry}
\bibitem{Folland} G. B. Folland, \emph{A Course in Abstract Harmonic Analysis}, second ed., Chapman \& Hall/CRC, Boca Raton, 2016.
\bibitem{Hall} B. C. Hall, \emph{Lie Groups, Lie Algebras, and Representations: An Elementary Introduction}, second ed., Graduate Texts in Mathematics 222, Springer, Cham, 2015.
\bibitem{Szarek57} S. J. Szarek, Metric entropy of homogeneous spaces, \emph{Banach Center Publications} \textbf{43} (1998), 395--410.
\bibitem{BG2012} J. Bourgain and A. Gamburd, A spectral gap theorem in $SU(d)$, \emph{Journal of the European Mathematical Society} \textbf{14} (2012), no. 5, 1455--1511.
\bibitem{BreuillardGelander57} E. Breuillard and T. Gelander, On dense free subgroups of Lie groups, \emph{Journal of Algebra} \textbf{261} (2003), no. 2, 448--467.
\bibitem{BreuillardGelanderCorrection57} E. Breuillard and T. Gelander, On dense free subgroups of Lie groups---revisited, arXiv:2605.20568 (2026).
\item[]\textit{Algebraic computation and complexity}
\bibitem{BPR} S. Basu, R. Pollack, and M.-F. Roy, \emph{Algorithms in Real Algebraic Geometry}, second ed., Algorithms and Computation in Mathematics 10, Springer, Berlin, 2006.
\bibitem{BJKP} V. D. Blondel, E. Jeandel, P. Koiran, and N. Portier, Decidable and undecidable problems about quantum automata, \emph{SIAM Journal on Computing} \textbf{34} (2005), no. 6, 1464--1473.
\bibitem{DJK} H. Derksen, E. Jeandel, and P. Koiran, Quantum automata and algebraic groups, \emph{Journal of Symbolic Computation} \textbf{39} (2005), nos. 3--4, 357--371.
\bibitem{deGraaf} W. A. de Graaf, Computing the Zariski closure of a finitely generated matrix group, arXiv:2505.02577v2 (2026).
\bibitem{NPSW} K. Nosan, A. Pouly, S. Schmitz, M. Shirmohammadi, and J. Worrell, On the computation of the Zariski closure of finitely generated groups of matrices, arXiv:2106.01853v6 (2025).
\bibitem{Shitov} Y. Shitov, A universality theorem for nonnegative matrix factorizations, arXiv:1606.09068v2 (2018).
\bibitem{Ramsey} F. P. Ramsey, On a problem of formal logic, \emph{Proceedings of the London Mathematical Society} (2) \textbf{30} (1930), no. 1, 264--286.
\bibitem{SJR2004}
S.~Singh, M.~R.~James, and M.~Rudary,
\emph{Predictive state representations: A new theory for modeling dynamical systems},
Proceedings of the Twentieth Conference on Uncertainty in Artificial Intelligence,
AUAI Press, 2004, 512--519. See Theorems~2 and~6; arXiv:1207.4167.
\item[]\textit{Entropy transport and explicit spherical spectral bounds}
\bibitem{PW61} Y. Polyanskiy and Y. Wu, Wasserstein continuity of entropy and outer bounds for interference channels, \emph{IEEE Transactions on Information Theory} \textbf{62} (2016), no. 7, 3992--4002. DOI: 10.1109/TIT.2016.2562630.
\bibitem{LPS61} A. Lubotzky, R. Phillips, and P. Sarnak, Hecke operators and distributing points on the sphere I, \emph{Communications on Pure and Applied Mathematics} \textbf{39} (1986), suppl., S149--S186. DOI: 10.1002/cpa.3160390710.
\bibitem{PLP61pre} A. Pinochet Lobos and C. Pittet, The exact convergence rate in the ergodic theorem of Lubotzky--Phillips--Sarnak, arXiv:1805.05261v2 (2019), Theorem 1.1. Published in expanded form as \emph{The exact convergence rate in the ergodic theorem of Lubotzky--Phillips--Sarnak and a universal lower bound on discrepancies}, \emph{L'Enseignement Math\'ematique} \textbf{67} (2021), nos. 1--2, 63--94. DOI: 10.4171/LEM/1003.
\item[]\textit{Nonhomogeneous Markov chains and observation tails}
\bibitem{Blackwell62} D. Blackwell, Finite non-homogeneous chains, \emph{Annals of Mathematics} (2) \textbf{46} (1945), no. 4, 594--599. DOI: 10.2307/1969199.
\bibitem{Cohn62} H. Cohn, On the tail $\sigma$-algebra of the finite inhomogeneous Markov chains, \emph{Annals of Mathematical Statistics} \textbf{41} (1970), no. 6, 2175--2176. DOI: 10.1214/aoms/1177696725.
\bibitem{GS62} M. Grabchak and I. Sonin, A zero-one law for Markov chains, arXiv:2011.04063 (2020), Theorem 1 and Section 5.1; published in \emph{Stochastics}, DOI: 10.1080/17442508.2021.1980569.
\end{thebibliography}

\end{document}
